% ============================================================
%  R. Nielsen, "Philosophisk Propædeutik i Grundtræk."
%  "Philosophical Propaedeutic in Outline."
%  Kjøbenhavn: Forlagt af den Gyldendalske Boghandling (F. Hegel);
%  Thieles Bogtrykkeri, 1857.
%  TRANSLATION (English) — from the verified Danish transcription.tex.
%
%  Method: prose translation of the whole book. Page-break comments
%  "% ---- printed p.N ----" are carried over from transcription.tex so the
%  two files stay aligned. Only natural-language prose and heading titles are
%  translated.
%
%  HOUSE STYLE / KEY TERMS (settled with the author, 2026-07):
%    Erkjendelse            -> knowledge / cognition
%    Erkjendelseslære       -> Theory of Knowledge   (Part I)
%    Videnskabslære         -> Theory of Science      (Part II)
%    Ideelære               -> Doctrine of Ideas      (Part III)
%    Sagerkjendelse         -> Factual Knowledge      (Part I.C)
%    Indledning / Slutning  -> Introduction / Conclusion
%    "i Grundtræk"          -> "in Outline"
%    Sensualisme/Idealisme  -> Sensualism / Idealism
%    Sperrsatz \emph{}      -> kept as \emph{} (italic emphasis)
%  Foreign-language quotations (German Hegel passages, Latin tags, Greek) are
%  kept VERBATIM and followed by a bracketed English rendering [ ... ].
%  Register: scholarly, moderately literal, preserving Nielsen's Hegelian flavour.
%
%  STATUS: IN PROGRESS. Done through printed p.95.
%        I. Theory of Knowledge (§§1--11, closed p.67). II. Theory of Science:
%        §12 intro; A. The Science of the Sciences (§§13--16: hindrances, science
%        vs. theology, subjectivity relative/infinite/absolute — Kierkegaard,
%        faith-vs-knowledge sophisms); B. The Particular Sciences (§17 learning
%        vs. science; §18 Fichte/Hegel/Apelt classifications; §19 history vs.
%        philosophy; §20 dead/living/ensouled forms; §21 auxiliary sciences,
%        faculties, theory/practice). C. Philosophy and the Particular Sciences:
%        §22 (Meno/anamnesis, a priori vs a posteriori); §23 (philosophy not a
%        subject-science — Eros/Poros/Penia; math & chemistry as keys; the
%        chemical trilogies); §24 (presentation & understanding); §25 (the
%        three-fold labour: assimilation / humanizing / totalizing — critical,
%        dialectical, critical-dialectical method) — II. CLOSED at p.117.
%        III. Doctrine of Ideas: §26 (What is an idea? — Kant on Plato's Ideas,
%        Notio, the teleology of the Kritik der Urtheilskraft, Schelling/Hegel,
%        the mechanism/purposiveness antinomy); A. The Idea in its Essentiality,
%        §27 (Plato's Ideas vs. concepts, the Theaetetus/Sophist critique — dense
%        polytonic Greek kept verbatim + bracketed English) through printed p.128.
%        §27 b) Ideas & sensible things (Parmenides — the One/Many α/β/γ, Philebus,
%        Timæus) + c) mutual relation of Ideas (Meno, Republic, Ideas-as-numbers);
%        §28 a) Aristotle's critique (third man, etc.) through printed p.139.
%        §28 b) Aristotelian principles (substance/matter/form) + c) the
%        Aristotelian entelechy (first mover / the Good / Reason — νόησις νοήσεως
%        νόησις); §29 a) the Dialectic of contradiction (Plato vs Hegel) through
%        printed p.150. NB: the p.129 orphan closing quote is reproduced as
%        printed → English quote balance carries a net −1 from p.129 (mirrors the
%        transcription). §29 a) (distinction vs identity dialectic; Hegel Antinomien/
%        Sprüchwörter), b) infinite negativity (Sextus/Skepticismus), c) speculative
%        form (Wissenschaft der Logik); closes A. B. Reality against the Idea: §30
%        (Apelt on Aristotle's εἶδος; Bacon's natural law), §31 (space/time/motion,
%        v=s/t) through printed p.161. §31 (rest) + §32 opening — Hegel
%        Naturphilosophie critique of motion (fall-law, Taylor series, planetary
%        orbits — Ptolemy/Copernicus/Kepler; teleology — Bacon/Leibniz/Laplace)
%        through printed p.172. Heavy German quotes verbatim + [English]; Latin/
%        French/Greek glossed. Quote balance now +1 (open−close): p.129 −1 offset
%        by the p.162–163 repeated-open quote +1 and the p.172 trailing-open +1
%        (closes on p.173). §32 (rest: Erde/Magnet quote) + §33 (Timæus solids;
%        crystal teleology; Eschricht; Feuerbach/Moleschott/Büchner/Vogt) — closes
%        B. C. The Validity of the Idea: §34 (letterspaced opening; Büchner/Kant on
%        materialism; analytic/synthetic — „5+7=12"; the Idea infinitely near)
%        through printed p.183. Quote balance back to 0 (trailing p.172 quote closed;
%        p.129 −1 offsets p.162–163 +1) — mirrors transcription 374/374. Resume at
%        §34 (printed p.184). See RESUME-NOTES.md.
% ============================================================
\documentclass[12pt]{book}
\usepackage[utf8]{inputenc}
\usepackage[T1]{fontenc}
\usepackage[english]{babel}
\usepackage[protrusion=true,expansion=true]{microtype}
\usepackage[margin=1.4in]{geometry}
\usepackage{setspace}
\usepackage{amsmath}
\usepackage{libertinus}
\usepackage{libertinust1math}
\usepackage{textalpha}   % polytonic Greek typed directly (occurs in Part III)
\usepackage{fancyhdr}
\usepackage[colorlinks=true,linkcolor=black,urlcolor=black,
  pdftitle={Philosophical Propaedeutic in Outline},
  pdfauthor={Rasmus Nielsen}]{hyperref}
\setlength{\headheight}{15pt}
\pagestyle{fancy}
\fancyhf{}
\fancyhead[LE,RO]{\thepage}
\fancyhead[RE]{\textit{Philosophical Propaedeutic}}
\fancyhead[LO]{\textit{\leftmark}}
\setlength{\parindent}{1.5em}
\setlength{\parskip}{0pt}
\onehalfspacing

\begin{document}

% ============================================================
% TITLE PAGE — PDF 9 (printed p.1)
% ============================================================
\begin{titlepage}
\centering
\vspace*{2.2cm}
{\Huge\bfseries Philosophical Propaedeutic\par}
\vspace{1.1cm}
{\LARGE in\par}
\vspace{0.9cm}
{\LARGE\bfseries Outline\par}
\vspace{1.1cm}
{\large by\par}
\vspace{0.9cm}
{\LARGE\bfseries R.\ Nielsen,\par}
\vspace{0.2cm}
{\normalsize Professor of Philosophy.\par}
\vspace{2.2cm}
\rule{4cm}{0.4pt}\par
\vspace{1.6cm}
{\large\bfseries Copenhagen.\par}
\vspace{0.4cm}
{\normalsize Published by the Gyldendal Bookshop (F.\ Hegel).\par}
\vspace{0.3cm}
{\small\emph{Thiele's Printing House.}\par}
\vspace{0.6cm}
{\large 1857.\par}
\end{titlepage}

\frontmatter

% MOTTO — Bacon, Novum Organum I, aph. LIV. Kept verbatim (epigraph), with an
% English rendering beneath, per the bracketed-translation house style.
\thispagestyle{empty}
\vspace*{0.34\textheight}
\begin{quote}
\itshape\noindent Adamant homines scientias et contemplationes particulares, aut
quia autores et inventores se earum credunt, aut quia plurimum in illis operæ
posuerunt, iisque maxime assueverunt. Hujusmodi vero homines, si ad philosophiam
et contemplationes universales se contulerint, illas ex prioribus phantasiis
detorquent et corrumpunt.
\par\vspace{6pt}\normalfont\hfill\textsc{Franc.\ Baco.}
\par\vspace{14pt}
\small\noindent [Men grow attached to particular sciences and speculations,
either because they fancy themselves their authors and inventors, or because
they have bestowed the greatest labour upon them and become most habituated to
them. But if men of this sort betake themselves to philosophy and universal
speculations, they distort and corrupt these out of their former fancies.
--- Francis Bacon.]
\end{quote}

\tableofcontents
\mainmatter

% ============================================================
% BODY BEGINS — printed p.3 (PDF 11)
% ============================================================

% ---- printed p.3 (PDF 11): Indledning / Introduction ----
\begin{center}
  {\Large\bfseries Introduction.}\par\vspace{4pt}
  \rule{2.5cm}{0.4pt}
\end{center}
\phantomsection
\addcontentsline{toc}{section}{Introduction}
\markboth{Introduction}{Introduction}

\noindent The philosophical propaedeutic sets itself to answer two fundamental
questions: \emph{what is philosophy?} \emph{what validity can be ascribed to
philosophy as a science?}

To answer these questions the propaedeutic must characterize philosophy by
pointing out the shapes which it has assumed at its various stages of
development; the general fundamental concept will then be determined more
closely as it runs through significant stages. The propaedeutic, in letting
philosophy come to be for contemplation in its own way, must everywhere take up
something historically given. But, although the historical is made use of, we
can here give neither an epitome of the history of philosophy in its entirety,
nor a separate detailed treatment of individual, arbitrarily chosen principal
parts on their own account; the historical must be grasped and used in such a
way that, within the connection determined by the nature of the matter, it
illuminates the tasks that arise, and serves the totality.

To characterize philosophy as a science in its relation to the other sciences,
it is doubtless necessary to enter more closely into an investigation of that
which brings about the distinctive character of the various sciences, and thus
to point out the principle of division on which their separation essentially
% ---- printed p.4 (PDF 12) ----
rests; but since the determination of what is universally scientific is and
remains the main thing, the treatment of the mutual relation of the diverse
sciences cannot take up the particulars of that ramification to such a degree
that there should here be any question of an encyclopædia of the sciences.

Insofar as philosophy itself is a science articulated into definite members, it
will likewise, for the sake of the fundamental concept, be necessary to point
out, within this science's own domain, the inner connection of the larger
divisions with the articulating principle; but since every survey in which the
individual parts are not at the same time treated with some detail necessarily
becomes superficial, the propaedeutic must, instead of setting itself the task
of an encyclopædic presentation of the disciplines of philosophy, rather lay out
its principal parts in such a way that, by a more detailed treatment of each
division on its own, it may find occasion to illuminate the systematic structure
from one or another definite side, and thereby gradually open an insight into
the architectonic relations of the whole.

In order to comprise all the separate considerations under one comprehensive
point of view, the propaedeutic begins with the conceptual determination that
philosophy is a scientific theory of consciousness.

The theory of consciousness shows what it means to become conscious of
something, and must to that extent distinguish between consciousness itself and
that which comes to consciousness. This distinction in turn demands the union of
the distinguished sides. An empty consciousness, a consciousness that is
conscious of nothing, is a contradiction: consciousness must have a content.

Although consciousness and its content are to that extent essentially the same
--- since the passivity and activity of consciousness are always determined in
likeness with that which, under this passivity and activity, is taken up into
consciousness --- they are nevertheless in reality so different that the
scientific orientation depends on whether attention is directed especially
either to
% ---- printed p.5 (PDF 13) ----
the consciousness determined by the content, or to the objectivity of the
content set over against consciousness.

Since the side of consciousness is everywhere, and must be, the predominant
thing in philosophy, it is necessary to investigate the fundamental relation
a) between consciousness and its content in general, b) between consciousness
and the particular content, c) between consciousness and the absolute content.
The main parts of the propaedeutic corresponding to this are then I) Theory of
Knowledge, II) Theory of Science, and III) Doctrine of Ideas.

% ---- printed p.5 (PDF 13): I. Erkjendelseslære / Theory of Knowledge ----
\begin{center}
  {\large\bfseries I.}\par\vspace{4pt}
  {\Large\bfseries Theory of Knowledge.}\par\vspace{4pt}
  \rule{2.5cm}{0.4pt}
\end{center}
\phantomsection
\addcontentsline{toc}{section}{\texorpdfstring{I.\quad Theory of Knowledge}{I. Theory of Knowledge}}
\markboth{Theory of Knowledge}{Theory of Knowledge}

\begin{center}{\bfseries \S\,1.}\end{center}

\noindent Knowledge presupposes identity of consciousness and content; as the
object is, so too must consciousness be: ``the like is understood by the like.''
Insofar as consciousness is sensing, the object is sensible; insofar as it is
thinking, thinkable.

The essential difference between sensation and thought can be recognized by the
difference between the Positive and the Negative. The object of sensation is,
without regard to whether the sense is outer or inner, always something
immediate, existing in the moment, and to that extent a Positive: the Negative
cannot be sensed. By suspending immediacy and positing negation, thought arises;
one can think \emph{about} the immediate, but its immediacy cannot be immediately
thought. Insofar as thought begins with the immediate, the understanding shows
itself as the intellectual sense. The sensing understanding is not a thinking,
but the immediate presupposition of thinking, a thinking in possibility: actual
thinking begins with the Negative.
% ---- printed p.6 (PDF 14) ----
In its pure, qualitative opposition to the immediacy of sensation, the negation
of thought is itself immediate. The immediately Positive and the immediately
Negative bound one another; by this bounding they mutually determine each other.
When it is asked what the Positive is, the answer must be given by means of a
Negative, \textit{omnis determinatio est negatio} [all determination is
negation]; likewise, on the other hand, the Positive must be presupposed when
negation and the Negative are to be determined.

As sensing, consciousness is essentially passive; here it holds that ``all our
knowing is a finding, receiving, and beholding (\textit{scimus, quia accepimus}
[we know because we have received])''; as thinking, consciousness is essentially
self-active. As regards self-activity, all our knowing (Cf. Sibbern:
\emph{Om Erkjendelse og Granskning}) is ``a doing and producing (\textit{scimus,
quia facimus} [we know because we make])''. The sensible is objectified in
representation and image, the thinkable in insight and conception.

Sensation is not merely heterogeneous with thought: what can be sensed cannot be
thought, and conversely, what can be thought cannot be sensed; but the
individual senses are also mutually heterogeneous: what can be seen cannot be
heard, what can be heard cannot be felt. Nevertheless these opposites too
determine one another mutually. The audible can be determined by the visible
(sound-figures), the sensible by the thinkable. In the passivity of the sensing
consciousness there is an activity, in the thinking activity there is a
passivity: the representation is obtained only insofar as consciousness itself
forms it; thought determines only insofar as it is led, compelled, and
determined by that which is thought; the sensing consciousness is thinking, the
thinking sensing; the acting I is passive, the passive acting (\textit{scientia
et potentia in idem concidunt} [knowledge and power coincide in one and the
same]).

Since it is now one and the same consciousness that is both passive and active,
both sensing and thinking, and since one and the same object can have both
sensible and
% ---- printed p.7 (PDF 15) ----
thinkable properties, it follows from this that there must not only be an
infinite reciprocal relation between sensation and thought, but that this
reciprocal relation must even be able to set itself up so differently that now
it is thought that is absorbed into sensation, now conversely sensation that
vanishes into thought.

Insofar as thought is absorbed into sensation in such a way that only the
sensible in the object is granted validity, we get a one-sided sensuous
tendency: \emph{Sensualism}; the corresponding one-sidedness then consists in
this, that sensation is grasped as the fleeting presupposition of thought, so
that only the thinkable in the object, the negative essence, the pure concept,
is acknowledged as the True: \emph{Idealism}. The infinite reciprocal relation
of sensation and thought will therefore first become evident when Sensualism and
Idealism have each contributed their share to a closer determination of the
truth in a double-sided \emph{Factual Knowledge}.

\begin{center}
  {\large\bfseries A.}\par\vspace{2pt}
  {\bfseries Sensualism.}\par\vspace{4pt}
  \rule{2cm}{0.4pt}
\end{center}
\phantomsection
\addcontentsline{toc}{subsection}{\texorpdfstring{A.\quad Sensualism}{A. Sensualism}}

\begin{center}{\bfseries \S\,2.}\end{center}

\noindent The consciousness of the individual begins with an immediate unity of
sensation and thought; but the advance of knowledge is conditioned by thought's
being brought out on its own in its opposition to sensation. Now so long as the
equilibrium between the two expressions of consciousness is preserved, so long as
the agreement between the sensible and the intellectual is still maintained,
individuality is happy, and its cultivation beautiful. But with growing
abstraction, as the imagination really begins to work over and transform
% ---- printed p.8 (PDF 16) ----
the elements of sensation, thought becomes restless, a tension sets in, and the
equilibrium is lost. The lost equilibrium --- whose mark is that now the
intellectual, now the sensible, now an overwrought spirituality, now a bold
naturalness, has gained the upper hand --- will then be re-established, and the
equilibrium thus re-established be lost again.

Corresponding transitions are seen in the history of philosophy.

Greek philosophy begins among the Ionians (Thales, Anaximander, Anaximenes) with
a naïve cognition of the non-sensible in the sensible; thought seeks the
principle of existence, and also divines that the principle must be ``something
unbounded,'' yet finds it most nearly represented by a given sensible material:
``Water,'' ``Air.'' This going-together of thought with sensuousness is not yet
Sensualism. Thought is here so far from feeling the pressure of the senses'
preponderance that it is precisely on the point of raising itself above the
parcelling-out of the sensible; it is Oneness that thought has here grasped with
youthful courage, as it first attempts a soaring flight toward the universal.

Only when thought has left immediacy behind, has consumed its content, and is
weary of mirroring itself in itself, does Sensualism emerge. Sensualism is a
conscious opposition, an express counterweight to the empty conceptual essence:
only after Greek thought (Pythagoreans, Eleatics, Plato) had volatilized the
material of sensation, held fast to the pure essence, and tried its hand at all
manner of abstractions, does Aristotle assert experience, and the repressed
Sensualism then properly comes forward in the theory of knowledge of the Stoics
and Epicureans.

The stages which in this respect characterize the directions of the movement in
Greek philosophy repeat themselves on an even larger scale in the universal
history of philosophy.

In contrast to the Middle Ages, antiquity on the whole exhibits the beautiful
equilibrium and natural unity of sensation and thought; but in the Middle Ages
abstract
% ---- printed p.9 (PDF 17) ----
spirituality commits such an overreach that a reaction thereby called forth has
been able to lead the modern age to a denial of spirit whose extent, sharpness,
and inner transparency is unknown among the ancients.

In the Middle Ages, when the natural is only something corruptible, the sensible
only something sinful, the present only a deception, it is naturally impossible
to know things as they are. Frightened by fancy and feeling, the thought of
science enters the service of faith, and is cowed by the supernatural; the cowed
thought finds only a formal freedom in the logical. Hence the emptiness, the
subtlety, the word-splitting, and the scientific barrenness in which Scholasticism
ends. Instead of grasping a reality, and nourishing itself on the working-up of a
digestible material, reflection fed upon itself, and was tormented by hunger.

The efforts that scientific hunger called forth drew minds to nature in two
mutually opposed directions, a speculative and an empirical; yet in such a way
that something of the one can be found again in the other.

\emph{The speculative direction.} Insofar as there can be any talk of a
speculative Sensualism, we find such among the Italian philosophers: Telesio,
Cardano, Bruno, and others. With a decided predilection for the forms of ancient
philosophy, these thinkers were inspired by the assurance of nature's identity
with spirit, by the contemplation of that divine life which animates the
world-organism, the living unity that works through all things.

That Aristotelian categories are applied to the interpretation of natural
phenomena is not in and of itself Sensualism; the Scholastics did the same, in
good understanding with the Hierarchy. The Sensualistic, on the contrary, comes
to light in the boldness, frankness, and unconditionedness with which an
omnipresent natural life displaces the Beyond, and sets itself in the place of
the supernatural. The realm of reason with which Giordano Bruno deals
(\textit{De la causa, principio ed uno})
% ---- printed p.10 (PDF 18) ----
is at the same time a realm of the senses, and the wonders to which Vanini
(\textit{De admirandis Naturæ, reginæ Deæque mortalium, arcanis}) points are owed
not to ``the eternal King,'' but to ``the Queen of mortals.''

\emph{The empirical path.} By sober research, by careful observation of the
phenomena (Tycho Brahe) and corresponding calculations, that path was at last
marked out which alone could lead to a true knowledge of nature. The law of
falling bodies on earth (Galileo) and the planetary orbits in the heavens
(Kepler) showed clearly enough that the path had been found. The investigators
who thus understood not merely how to open the eye, but also how to arm the eye
(the telescope), then, in the name of science, asserted the right of the senses,
and set in force facts whose incontrovertibility must at length subdue the
imaginings of the Middle Ages, and shake Scholastic metaphysics in its
foundations.

Were such efforts to be grasped, judged, and presented in a theory of knowledge
framed to the needs of natural science, then it would doubtless above all be a
matter of impressing two chief demands: thought's liberation from abstract
theories (from \textit{anticipationes naturæ} [anticipations of nature]), and the
leading-back of knowledge to true experience (\textit{interpretatio naturæ}
[interpretation of nature]). By working against formalism, emphasizing
``Induction,'' and bringing the concept of natural law into harmony with the
demand of physics, Francis Bacon of Verulam (†1626; \textit{Novum organum
scientiarum}), though without a natural-scientific initiation, has nonetheless,
in inspired traits, pointed out a new foundation, and marked the transition to a
new age.

\begin{center}{\bfseries \S\,3.}\end{center}

Advancing Empiricism soon discovered the misinterpretation of the a priori that
lay hidden in the doctrine of ``innate ideas.'' Against this misinterpretation
J. Locke (†1704) directed his ``philosophical essay on the human understanding''
(\textit{An essay concerning human understanding}
% ---- printed p.11 (PDF 19) ----
\textit{in four books}), whose content we may here consider under the following
points of view: the origin of ``ideas,'' the essence of ``ideas,'' and the
knowledge of things.

\emph{a) The origin of ``ideas.''} There are no innate ideas. Even if it could be
shown that there were propositions about whose universality all must agree, this
would still prove nothing, so long as the universal acceptance of these
propositions could be explained in another way. But so far are such propositions
from being found, that nations at different levels of culture are even at odds
about the very most important practical principles. If one sets forth the
commandment ``Whatever you would have others do to you, do also to them,'' then
everyone who has not heard the commandment before, but thinks it over, will of
course ask: why? But to ask for the grounding of a commandment that is innate is
a contradiction.

As regards the theoretical propositions, such as ``Whatever is, is; it is
impossible that one and the same thing can at the same time both be and not be,''
children, idiots, and the uneducated are wholly unacquainted with their meaning.
But if such principles really were innate to the understanding, then all would
have to be aware of them from earliest childhood, for to be in the understanding
is the same as to be in knowledge.

Likewise it is an evasion when, in order to defend innate ideas, one appeals to
the fact that these ideas immediately become evident to consciousness as soon as
a person begins to use his reason. The relation is precisely the reverse: it is
not universal but particular propositions, such as ``Sweet is not bitter,'' of
which one first of all becomes conscious when one begins to think. That
theoretical and practical principles should be innate is on the whole just as
absurd as that arts and sciences should be innate. It is therefore impossible to
be born with ideas;
% ---- printed p.12 (PDF 20) ----
the child's soul is like a clean sheet of paper, as empty as a \textit{tabula
rasa} [a blank tablet].

Whence, then, does the understanding get its ideas? From experience. But
experience is of a twofold kind: insofar as it is made by an observation of
external objects conducted through the senses, we call it sensation,
\emph{Sensation}; insofar as it rests on observation of our own understanding's
activity, we call it the inner sense, \emph{Reflection}. Sensation and Reflection
are to be regarded as ``the only channels that convey knowledge to the
understanding, the only windows through which the light of ideas falls into the
dark chamber of the soul.''

\emph{b) The essence of ``ideas.''} The ideas coming from experience are partly
simple, partly compound.

The simple ideas furnish the material of knowledge, are, so to speak, the letters
in the alphabet of existence. In that these simple ideas press themselves upon
consciousness from without, the soul is passive, and takes up the impressions of
the surrounding world as the mirror takes up
% [errata applied (Rettelser p.12): print ``the objects' image''; read ``the object's image'']
the object's image; if, on the other hand, they are obtained through Reflection
--- as the idea of thought, the idea of will --- the soul is active. The simple
ideas are partly such as come to the understanding through a single sense: colors
through the eye, tones through the ear, impenetrability through the sense of
touch; partly such as are conveyed to it through several senses, as when space
and motion are apprehended both by touch and by sight.

By the connections and compositions of the simple ideas the soul can now know the
objects of the surrounding world, ``just as in language, by manifold compositions
of letters, one can form syllables and words''; but, in that the compound,
complex ideas are referred to three chief classes --- modes, substances, and
relations --- it becomes above all the task to give an account of the concept of
substance.

The nominal essence indicates the relation between genus and species; the real
essence, on the other hand, denotes the inner
% ---- printed p.13 (PDF 21) ----
peculiarity of the existing thing. In the simple ideas and modes the real and the
nominal essence are one and the same; with respect to ``substances'' they are
different. That a surface bounded by three straight lines is a triangle expresses
the unity of the real and the nominal essence. It is otherwise with the part of
matter that forms the ring on my finger; on the real peculiarity of the
non-sensible parts depend the properties of its color, weight, softness,
hardness, and so on. After the observed properties the name ``Gold'' is formed,
an expression that denotes the nominal but not the real essence.

By sensation and reflection we experience how a number of simple ideas come
together in a peculiar connection, presuppose on that account a necessary ground
of connection, and refer this ground to a substratum subsisting on its own,
namely \emph{substance}; but is it now possible to know substance?

\emph{c) The knowledge of things.} As the bearer of the qualities by which the
simple ideas are produced in us, substance can be known only in this, that a
distinction is made between the properties that necessarily belong to the thing in
itself, and prove the agreement of the simple ideas with the objects
(\textit{qualitates primariæ} [primary qualities]), and such properties --- color,
tone, taste, smell, warmth, cold --- as do indeed proceed from the things' action
upon the soul, but are nevertheless wholly different from the objective essence of
the things (\textit{qualitates secundariæ} [secondary qualities]).

In this way Locke would determine the opposition between real truth
(\textit{real truth}) and mental truth (\textit{mental truth}); but although
substance is distinguished from all other complex ideas precisely in this, that it
has its archetype outside us --- one does not connect a sheep's bleating with the
representation of a horse ---, and precisely therein makes known its objective
reality, it is nevertheless impossible for us to penetrate into substantiality
itself.
% ---- printed p.14 (PDF 22) ----
What we call the essential does not bear upon the individual. ``It is doubtless
necessary that I be what I am: God and nature have formed me thus; but I know
nothing that is essential to me, the individual, as individual: an accident, an
illness, can bring about incredible changes.'' So the essential bears, then, upon
the species; but it can by no means be the real essence that determines the
species, since the real essence is unknown to us. We are referred to a collection
of sensible qualities that we can observe; but even when our observations are
conducted with the greatest possible accuracy, our knowledge is just as far from
corresponding to that inner constitution of the thing from which these qualities
follow ``as a peasant's representation of the clock in Strasbourg is far from
corresponding to the inner, artful structure of that remarkable machine.''

What Locke had begun, Condillac (†1780; \textit{Essai sur l'origine des
connaissances humaines}) continued. By letting the two sources of knowledge ---
sensation and reflection --- flow together in such a way that the latter vanishes
into the former, and by laying a reinforced emphasis on sensation, Condillac
arrived at the result that the so-called inner sense is only a reflex of the
outer, and that consequently all expressions of spirit represent only modified
sense-impressions. Once Sensualism, thus carried through, has first shown that the
True is only something sensible, it lies near to reverse the proposition and say:
only the sensible is the True. Here is a materialism for which everything
spiritual is a figment (La Mettrie).

\begin{center}{\bfseries \S\,4.}\end{center}

From presuppositions by which a one-sided standpoint is expressly designated,
consequences can be developed partly in a coordinate, partly in a superordinate
direction. To consequences in a direction coordinate with the standpoint belongs
everything that unfolds its character, and raises its significance
% ---- printed p.15 (PDF 23) ----
to an ever wider compass. In this sense Condillac and La Mettrie have driven the
consequences of Sensualism to the extreme.

To consequences in a direction superordinate to the standpoint belong those that
lead out beyond the standpoint itself, and thus either establish or at least
introduce a new standpoint. In this sense Hume's (†1776) Scepticism (\textit{A
treatise of human nature}) is to be grasped as a higher consequence of Lockean
Sensualism.

To test the validity of concepts, one must --- herein Hume agrees with Locke ---
hold to experience. Those concepts that can prove their origin from
sense-impressions (\textit{impressions}), or at least from the reflections
corresponding to these, ought to be acknowledged; those that cannot, are of
spurious descent, and ought to be rejected.

According to Locke, ``the idea of cause and effect'' arises when, in considering
the constant change of things, we observe how certain things --- be they
substances or qualities --- begin to exist through the law-governed activity of
another. But what validity now has this observation? whence do we know that two
things stand to one another as cause to effect? \textit{A priori} we cannot know
it; for a priori knowledge is analytic, but the knowledge of the causal concept
must be synthetic, since the effect is really something other than the cause.
Neither do we know it from experience; for experience shows us only that the one
fact partly follows the other in time (succession),
% [errata applied (Rettelser p.15): print ``Continuity''; read ``Contiguity'']
partly touches the other in space (contiguity); but no experience can teach us
that the one essentially follows from the other, as the effect from its cause.

To assert the validity of substance Locke had already referred to habit; this
indication Hume follows: ``Our inferences from experience are grounded solely on
habit; we are accustomed (\textit{the mind is carried by habit}) to see the one
thing follow partly along with, partly after the other, and then habituate
ourselves thereby to the representation that the one
% ---- printed p.16 (PDF 24) ----
% [errata applied (Rettelser p.16): print ``Continuity-''; read ``Contiguity-'']
must be an effect of the other. Out of the relation of contiguity and succession
our imagination forms the causal relation: \textit{cum hoc aut post hoc, ergo
propter hoc} [with this or after this, therefore because of this].''

That the threads of the investigation here run together in the question of the
causal concept as in their nodal point, shows the sureness with which the task has
been posed.

Once the causal concept has fallen, then the categories --- force, substance,
necessity, truth --- fall as of themselves; the bonds of essence that were to hold
things together must now all burst.

As proceeding from a cooperation of sensation and reflection, the concept
``force'' might, in Locke's opinion, well have its validity; but how is it now
possible to know force? It might at first glance seem as though the relation
between soul and body furnished, in this respect, a reliable presupposition. In
that certain bodily organs involuntarily obey our will, spirit does indeed act as
force, and we need therefore only to become conscious of our voluntariness in
order to become conscious of an expression of force. This is, however, a
deceptive reflection; partly we do not know the means by which spirit acts upon
the body; partly it is not, after all, all the body's organs whose activity is
subject to the will. Were force to be an object of experience, then --- what is
impossible --- cause and effect would also have to be an object of experience.

By presenting substance as an Unknown, Locke had already introduced the doubt
about its existence. How are we now to know substance? Were the validity of
substantiality to admit of proof, then the inner necessary connection of the
properties that in substance are united into unity would have to admit of proof;
but with causality the necessary connection (\textit{this idea of a necessary
connexion}) has fallen away. To refer to the difference between original and
derived properties does not help. The derived properties are all subjective, and
thus have only
% ---- printed p.17 (PDF 25) ----
existence in our apprehension; as for the original ones, they let themselves at
last be resolved into one single one, namely ``Solidity,'' but Solidity is an
empty word, since we, when all the other properties are suspended, cannot possibly
make the Solid an object of any experience.

In reliance on the ``inner sense,'' Locke had still held fast to ``the idea of
will,'' and acknowledged the self. But how should it now be possible to assure
oneself of the reality of the I? Our feelings, thoughts, and representations
cannot possibly be regarded as expressions of force of an I, when the category
expression-of-force has no validity; they cannot possibly be referred to an acting
self, as to their true cause, when the causal concept has no truth.

The remnant of reality that still attached, in Locke, to the concept of truth
(\textit{real truth}) has thus vanished: truth is gone, probability has stepped
into its place. The I, the self, is then in reality nothing other than a figment,
which arises in this, that an aggregate of many rapidly successive representations
is laid under a substratum.

That a philosophy which subordinates thought to sensation, even to such a degree
that concepts must lose their universal validity, that the laws of things must
vanish before the rules of the association of ideas, the ground of reason before
habit, truth before probability, the self-knowledge of the I before a deluding
figment, is in itself neither tenable nor satisfying, goes without saying;
nevertheless Hume's Scepticism, in contrast to dogmatic Materialism, is a higher
power, a spiritual force whose attack upon thought can transform itself into a
defense of thought.

\begin{center}\rule{2.5cm}{0.4pt}\end{center}
% ---- printed p.18 (PDF 26) ----
\begin{center}
  {\large\bfseries B.}\par\vspace{2pt}
  {\bfseries Idealism.}\par\vspace{4pt}
  \rule{2cm}{0.4pt}
\end{center}
\phantomsection
\addcontentsline{toc}{subsection}{\texorpdfstring{B.\quad Idealism}{B. Idealism}}

\begin{center}{\bfseries \S\,5.}\end{center}

\noindent If Sensualism can be recognized by this, that it examines the senses in
order to judge thought, then Idealism can in turn be recognized by this, that it
makes thought the judge of the senses; if Sensualism sets things above
consciousness, then Idealism reverses the relation, and sets consciousness above
things.

In that Descartes (†1650; \textit{Principia philosophiæ}) begins with doubt, he
wishes, like Bacon, to cast off prejudices, and free the mind from blind trust in
authority; but, in that by doubting everything (\textit{de omnibus dubitandum est}
[everything is to be doubted]) he can also
draw into doubt the validity of natural things and of experience as well, he goes
further than Bacon.

Since doubt is a work of thought, thought is, in the midst of doubt, raised above
doubt; everything can be doubted, but the validity of the doubting thought cannot
be drawn into doubt; to doubt doubt is to confirm doubt.

Empiricism, which begins with an unconditional reliance on the immediate
sense-impression, ends consistently in Scepticism; for Scepticism is an activity
of spirit in which thought, after having stared itself blind on the sensible,
incessantly comes to itself as a stranger, tears itself down, and cannot know its
own essence.

In the idealistic doubting of everything, thought begins precisely by finding
itself, in that it begins by removing the foreign, and shaking off the appendage
of tradition. This doubting is then so far from losing itself in Scepticism that
the thinker rather wishes to assure himself of the validity of everything that
withstands the testing of thought.
% ---- printed p.19 (PDF 27) ----
The thinker then not only sees the unity of his thought with being
(\textit{cogito ergo sum} [I think, therefore I am]); but also gains certainty
that consciousness is borne, penetrated, and filled by an absolute being, by the
Unconditioned, whose image is in the soul, whose forms (\textit{ideæ innatæ}
[innate ideas]) shine with a clarity that can proceed only from the Perfect.

In Descartes, however, Idealism is still restricted in this, that over against the
being that is in thought there stands a material being opposite thought; to
suspend this substantial separation, and transform thought and being into
attributes in one and the same substance, is an advance owed to Spinoza (†1677).
Insofar as the question here is of the validity of thought and its identity with
being, Spinozism is decidedly Idealism.

As regards the soul's a priori self-activity, Leibniz (†1716) went a step
further. In his ``Nouveaux essais'' he works against Lockean Sensualism, and frees
the assumption of ``innate ideas'' from the misunderstanding that could justify
Locke's objections. To the Sensualists' \textit{nihil est in intellectu, quod non
antea fuerit in sensu} [nothing is in the intellect that was not previously in
sense], there is then added a \textit{nisi intellectus ipse} [except the intellect
itself]. That ``ideas'' are innate must not be misinterpreted as though at birth
they lay ready-made in the soul, since on the contrary they are in it only
according to disposition. This view is in essential connection with the doctrine
of monads; as monad the soul is originally self-active, can receive no impression
from without, but on the contrary develops a world in its interior.

Whereas Locke compared the soul to an unwritten tablet, Leibniz compares it to
``a block of marble in which the veins preform the shape of the statue.'' The
principle of contradiction (\textit{principium contradictionis}), the proposition
of sufficient reason (\textit{principium rationis sufficientis}), and the
so-called \textit{principium indiscernibilium} belong among ``the innate ideas.''
% ---- printed p.20 (PDF 28) ----
How Idealism, by thinking over the relation between the senses and the sensed, and
pointing out the contradictions that everywhere accompany the cognition of an
external, can transform the whole outer world into a world of spirit, has been
shown in particular by G. Berkeley (†1754; \textit{A treatise concerning the
principles of human knowledge}; and \textit{Three dialogues between Hylas and
Philonous}). --- Convinced that it is the assumption of matter's unconditional
existence that especially favors Scepticism and Atheism, Berkeley reckons this
assumption among the most meaningless and confused that can ever arise in a human
head.

Sensible things are such as are immediately (\textit{immediately}) apprehended by
the senses, whereas the cause of things, which cannot be sensed, must expressly be
thought, and, in being thought, be in consciousness. If I see that the one part of
the sky is red, the other blue, then I infer that there must be a cause of this
color-difference; but the cause cannot be said to be a sensible thing (\textit{or
perceived by the Sense of Seeing}). Nor can I, when I hear various tones, be said
to hear the cause of the tones; just as little as I, when by feeling I experience
that a thing is hot or heavy, can be said to feel the cause of the heat or the
weight.

Sensible things are aggregates of sensible qualities (\textit{Combination of
sensible Qualities}); these sensible qualities are not outside the soul, but in
the soul. As little as one can ascribe to the needle a sensation of the pain that
its prick causes us, so little can one ascribe warmth to the fire. To have a high
degree of warmth, to be hot, is to be in pain; were the fire really hot, then it
would have to burn itself as it burns us, it would then have to be in pain; but
that is absurd.

Not only of the derived properties, but of the original ones too, it holds that
they exist only in our consciousness. Were properties such as extension, figure,
weight, density,
% ---- printed p.21 (PDF 29) ----
motion something outside consciousness, then one and the same thing would have to
be both large (namely when seen near) and small (insofar as seen at a great
distance), both smooth and rough, both straight and crooked, and the same motion
both slow and fast; but that is absurd.

Just as absurd is the opinion that the named properties should have a substratum;
the substratum on which extension and the other properties should depend would,
in order to bring about extension, have to have an extension already beforehand,
and so on to infinity.

The infinite divisibility of matter is generally assumed. There is then in every
finite matter an infinite number of parts that the senses do not yet apprehend;
but the more the eye is sharpened, the more parts it sees, and the larger each
part becomes: could the eye be sharpened to infinity, each smallest part would
become infinitely large. This materialism contradicts itself; for it follows from
this that every body is in itself infinitely extended, that is to say: without any
shape or figure. It is therefore not an essence of matter, but consciousness
itself, that forms all changes in the bodies of which the visible world
(\textit{the visible World}) consists: matter is an idea.

As for mathematical natural philosophy, mathematics forgets to give a sufficient
account of its principles. Arithmetic reckons with units (\textit{Unities}), but
is not able to show what ``the unit'' really is; in geometry the infinite
divisibility of finite extension plays a chief role; but finite extension is, as
already shown, only an idea in the soul. Since it is now evident that no one is
able to separate infinitely many parts either in a finite line or in a finite
surface or in a finite body, everyone can, by making clear to himself the idea of
extension, convince himself that the alleged infinity is not present.
% ---- printed p.22 (PDF 30) ----
What in everyday speech we call things, we must then, according to the given
explanations, call ideas. It is doubtless a hard saying, if for daily use we
should eat ideas, drink ideas, and clothe ourselves in ideas; but usage is one
thing, truth another. That things outside us are transformed into ideas in our own
soul is not to be understood as though the sensible representations rested on our
own arbitrariness. We can arbitrarily open and close our eyes, but the shapes and
colors that show themselves to our opened eyes affect us with necessity.

Our sensible sensations and representations must doubtless be referred to a cause
independent of us. But since cause and effect must be of like kind, and since a
thing cannot possibly impart to us what it does not itself possess, it is just as
absurd to assume that a thing that does not itself have sensation should be able
to arouse sensation, as that a brain-impression should be able to become a
representation. There exist therefore no bodies, but only spirits; it is not a
material cause, but an almighty spirit, that through all sense-images,
representations, sensations, and impressions lets the finite spirits impart
themselves to one another.

By thus explaining the sensible and abolishing the material, this Idealism has
transformed itself into a Spiritualism that annihilates nature, appeals to the
supernatural, and culminates in Theology. \textit{It is therefore plain, that
nothing can be more evident to any one, that is capable of the least Reflexion,
than the Existence of GOD, or a Spirit who is intimately present to our Minds,
producing in them all that variety of Ideas or Sensations, which continually
affect us, on whom we have an absolute and intire Dependence, in short, in whom we
live and move, and have our Being.} (\textit{Principles of hum. knowl. London
1734. S. 166}).
% ---- printed p.23 (PDF 31) ----
\begin{center}{\bfseries \S\,6.}\end{center}

Among the many mutually crossing approaches to a theory of knowledge, Scepticism
and Dogmatism meet as conflicting results. In Scepticism (Hume) thought is at
variance with itself, and experience stands as an unexplained fact; in Dogmatism
(Chr. Wolff) the unity of being with thought is grasped so immediately that the
rigid schema of the understanding does indeed dazzle by its proofs, but lacks the
grounding. Against these mighty one-sidednesses I. Kant (†1804; \textit{Kritik der
reinen Vernunft}) comes forward.

Descartes's Idealism, which makes the existence of objects in space unprovable and
doubtful, Kant calls the problematic; Berkeley's, which presents the external
existence of things as false and impossible, he calls the dogmatic. Dogmatic
Idealism is unavoidable when space is to be regarded as a property belonging to
things in and for themselves; for then every thing for which such a space should
furnish the condition is an absurdity. The ground of this Idealism is suspended in
``the transcendental Aesthetic.'' Problematic Idealism, on the other hand, is
first refuted when proof is furnished that even our inner experience, not doubted
by Descartes, is possible only under the presupposition of outer experience.

Although the Kantian critique of reason, as a systematic whole in itself, has by
no means come into being out of finite regard for this or that view, it can
nonetheless, if one will, be grasped in such a way that the two parts of ``the
transcendental Doctrine of Elements,'' namely the Aesthetic and the Logic, may
chiefly be said to be directed against Scepticism, its third part, namely the
Dialectic, against Dogmatism, and ``the transcendental Doctrine of Method''
against both. We therefore treat, in this connection, the Kantian theory of
knowledge insofar as it characterizes itself, by its opposition to the named
extremes, as a ``\emph{transcendental}'' Idealism. Thus:
% ---- printed p.24 (PDF 32) ----
\begin{center}{\bfseries 1)\quad Against Scepticism.}\end{center}

That all our knowledge begins with experience is beyond all doubt; but from this
it does not yet follow that knowledge in each and all has its source from
experience. There are a priori sensible intuitions, a priori categories of the
understanding, and a priori ideas of reason.

a) ``The transcendental Aesthetic'' deals with the pure sensible intuitions:
space and time. Space is the form of the outer sense, time of the inner. If we
abstract from the objects outside us, we get space back as the pure, universal
form; if we abstract from the changes observed by the inner sense, empty time
remains.

That intuitions of space and time are a priori can be shown partly from a
metaphysical, partly from a transcendental point of view. The metaphysical
consideration proves that the intuitions space and time cannot possibly be
borrowed from experience, since one must precisely presuppose them in order to be
able to make an experience; it proves further that space and time are not
concepts but intuitions: the individual spaces are contained as parts in universal
space, likewise the individual times in universal time, whereas the universal
concepts, on the contrary, always comprise the particular ones as members under
themselves, never as parts within themselves. The transcendental consideration
proves the a priori origin of space and time in an indirect way, in that it shows
that certain sciences, such as, for example, pure mathematics, would be impossible
if there were no universal and necessary propositions about space and time; but
universal and necessary propositions are never had from experience, they are had
only \textit{a priori}.

b) ``The transcendental Analytic'' treats the pure concepts of the understanding.
The understanding is, negatively determined, a non-sensible faculty of cognition;
outside intuition there is only a cognition by concepts. To intuition as
receptivity
% ---- printed p.25 (PDF 33) ----
corresponds the spontaneity of thought: intuitions rest on ``affections,''
concepts on ``functions.'' Intuition bears immediately on the object, the concept,
on the contrary, mediately by means of another representation. The understanding
judges; the judgment is a mediate cognition, a representation of a representation.

When we then abstract from the content of a judgment, and look only at the bare
form of the understanding, we find that the function of thought lets itself be
brought under four titles, each of which again contains three moments. In
judgments there lie, namely --- insofar as they are divided, according to
Quantity, into universal, particular, and singular --- the categories: totality,
plurality, unity; insofar as they are divided, according to Quality, into
affirmative, negative, (limiting), the categories: reality, negation, limitation;
insofar as they are divided, according to Relation, into categorical, hypothetical,
and disjunctive, the categories: substance and inherence, cause and effect,
reciprocity; and finally, when they are distinguished, according to Modality, into
problematic, assertoric, and apodictic, the categories: possibility, actuality,
and necessity.

That these twelve categories, from which the rest let themselves be derived, are
not owed to experience, but are a priori concepts, is seen from this, that these
too are just as necessary conditions for the possibility of an experience as the
pure intuitions. In and for themselves the categories are empty forms that get
their material conveyed only through intuition; only when the sense-content of
intuition is absorbed into the form of the understanding does experience come
about. Yet a difficulty meets us here. The objects, which are of sensible nature,
are heterogeneous with the concepts: how then does it come about that a posteriori
sense-objects can be subsumed under pure concepts of the understanding, and
judgments \textit{a priori} be formed therefrom? This happens in that a third,
which on the one side is a priori, on the other side sensible, steps between,
namely, the pure intuitions: space and time. In order in this respect to bring the
concept into intimate union with intuition, the imagination forms ``the tran-
% ---- printed p.26 (PDF 34) ----
scendental schema.'' Thus Quantity has for its universal schema the time-series or
number, so that the pure concept of magnitude meets intuition, in that we in
imagination produce units, and thereby go from unity through plurality to
totality. Quality's schema is the time-content: the Real fills a given time, and,
when the content of reality is gone, negation is intuited in the emptiness of
time. The schema of the categories of relation is the time-order, for
substantiality is the persistence of the Real in time, the members of causality
follow one another with necessity in time, and reciprocity is the regular
coexistence of two or more expressions of substance in time. The schema of
Modality is the time-comprehension: the schema of possibility signifies a
representation's agreement with the conditions of time in general, the schema of
actuality an object's existence in a given time, and that of necessity an object's
being at all times.

Since the objects we can observe all together fall in time, they must necessarily
also fall under those schemata of time. This is proved in ``the Analytic of
Principles.'' At the head of the analytic judgments stands the principle of
contradiction; but the highest principle of the synthetic judgments lays down that
every possible experience ``must stand under the conditions that are necessary for
bringing the manifold impressions of intuition to synthetic unity.'' α) All
phenomena are then, since they can be apprehended only under the forms of space
and time, extensive magnitudes; ``the Axioms of Intuition.'' β) All phenomena are,
since every sensation can be strengthened and weakened in a definite degree,
intensive magnitudes; ``the Anticipations of Perception.'' γ) All phenomena are,
according to the schema of the categories of relation, subject to necessary
determinations of time; the substantial remains in time; the accidental changes in
time; the temporal change of accidents expresses a relation between cause and
effect; ``the Analogies of Experience.'' δ) Finally all phenomena are, in
agreement with the schema of the categories of modality,
% ---- printed p.27 (PDF 35) ----
determined in such a way that what agrees with the formal conditions of experience
is possible; what agrees with its material conditions, actual; what exhibits
universally valid conditions of experience, necessary; --- ``the Postulates of
Empirical Thought.''

c) ``The transcendental Dialectic'' investigates the ideas of reason.

Our cognition begins with the senses, passes thence over to the understanding, and
ends in reason as in its highest. If the understanding is the faculty of bringing
phenomena under a unity by means of rules, then reason is in turn the faculty of
bringing the rules of the understanding to unity under principles. The
understanding develops the conditioned, reason points out the unconditioned; the
understanding furnishes the premises, reason draws out the conclusion. From the
three logical inferences of reason, namely: the categorical, the hypothetical, and
the disjunctive, the three speculative ideas of reason are derived: the idea of
the soul as a thinking substance (the psychological idea), the idea of the world
as the aggregate of all phenomena (the cosmological idea), the idea of God as the
highest condition for the possibility of all (the theological idea).

That the ideas of reason are not abstracted from experiences is seen immediately
from this, that in the whole world of experience no object can be pointed out that
corresponds to any of those ideas. That these ideas nonetheless have validity,
namely a priori validity, is understood from this, that they are the indispensable
lodestars of the understanding, the regulative principles of cognition. We order
our subjective faculties by means of the idea of the soul, the infinite series of
conditions by means of the idea of the world, the ascending steps of causes by
means of the idea of God.

\begin{center}{\bfseries 2)\quad Against Dogmatism.}\end{center}

The a priori is, according to Kant, solely in our consciousness, and has nothing
to do with the essence of things. Already ``the
% ---- printed p.28 (PDF 36) ----
transcendental Aesthetic'' warns us against regarding space and time as being
outside us. Were one to set up the proposition that things are in space and time,
then it would be a delusion; not in themselves, but only for us are things in
space and time, and that in such a way that the phenomena for the outer sense are
in space as well as in time, the phenomena for the inner sense only in time. Just
as emphatically ``the transcendental Analytic'' impresses that the categories of
the understanding have universal validity only insofar as they hold for us, in
that we apply them to things as objects of experience, but not to things in and
for themselves. In that we see objects through subjective forms, and apprehend
them in subjective concepts, we are like those who see the colors of things
through painted spectacles.

The discord between thinking and being that Kant thus brought about by impressing
from first to last the separation of the thing as object of experience and the
thing in itself (\textit{das Ding an sich} [the thing in itself]) --- a separation
that is further illuminated by the way in which, with respect to \textit{Phænomena}
and \textit{Noumena}, he demonstrates ``the Amphiboly of the Concepts of
Reflection'' --- is well suited to give uncritical Dogmatism a fundamental blow;
this blow becomes especially palpable in the critique that the ideas of reason
must undergo in ``the transcendental Dialectic.''

The Dialectic, ``\textit{eine Logik des Scheins} [a logic of illusion],'' is to
clear up the illusions in which reason, by applying the categories of the
understanding to the unconditioned, must necessarily entangle itself. This
``transcendental illusion,'' this optical self-deception of reason, shows itself
in the different ideas of reason from a different side.

a) The Paralogisms of pure reason. Rational psychology makes the soul into ``a
soul-thing with the attribute of incorporeality (immateriality); into a single
substance with the attribute of incorruptibility (incorruptibility); into a
numerically identical intellectual substance with the predicate of personality;
% ---- printed p.29 (PDF 37) ----
into a spaceless, thinking being with the predicate of immortality.'' These
propositions of rational psychology are sheer fabrications. At the ground of the
four paralogisms belonging here can lie nothing other than the simple, in and for
itself empty representation: ``I think''; but this ``I think'' is neither
intuition nor concept, it is only an act of consciousness that accompanies, bears,
and connects all my representations and concepts. That I who think can always be
regarded as subject, a something that does not merely, as a predicate, hang upon
thought, is an apodictic, indeed even identical, proposition; but it does not
signify that I as object am a being or substance subsisting for myself. The proofs
of immortality rest on deceptive inferences. ``I can indeed separate my pure
thinking ideally from my body; but from this it by no means follows that my
thinking can also persist really, after being separated from the body.''

b) The Antinomies of pure reason. To determine the chief sides from which the idea
of the world shows itself, one is guided by the categories given in the table. In
respect of the Quantitative, there is then here something to lay down about the
totality of the world's times and spaces, and as regards Quality, something about
the divisibility of matter; in respect of Relation the causality of the effects
given in the world must be pointed out, and, as regards Modality, the conditions
of the contingent are pursued until at last they are absorbed into the necessity
of the whole.

At each of these points, assertion stands against assertion (antithesis against
thesis) with proof of equal force. Against the thesis: ``The world has a beginning
in time and is enclosed within the limits of space,'' holds the antithesis: ``The
world is infinite in respect both of time and of space''; against the thesis:
``Every composite substance in the world consists of simple parts,'' the
antithesis: ``No composite thing in the world consists of simple parts''; against
the thesis: ``Besides causality according to natural laws there holds in the world
a
% ---- printed p.30 (PDF 38) ----
causality of freedom,'' the antithesis: ``there is no freedom, but everything in
the world happens solely according to natural laws''; against the thesis: ``To the
world belongs something that, either as a part of it or as its cause, is a simply
necessary being,'' the antithesis: ``There exists no simply necessary being as
world-cause, neither in the world nor outside it.''

These cosmological antinomies then show us the dialectical struggle, the
self-contradiction, in which reason entangles itself when by the idea of the world
it wishes to think the objective essence of the world.

c) The Ideal of pure reason. ``The category'' is empty, but can nonetheless find
its object in experience; ``the idea,'' which has no object of experience at all,
is thus more remote from objective reality than the category. ``But still more
than the idea, that seems to be removed from objective reality which I call the
Ideal, and by which I understand the idea, not merely \textit{in concreto} [in the
concrete], but \textit{in individuo} [in the individual], i.e. as a single thing
determinable and determined by the idea alone.'' After having shown how reason
determines its Ideal in different ways, namely as primal being (\textit{ens
originarium} [original being]), insofar as it has it in itself, as the highest
being (\textit{ens summum} [highest being]), insofar as it has it above itself,
and as the being of beings (\textit{ens entium} [being of beings]), insofar as
everything is conditioned by it, Kant passes over to an attack on the
``speculative reason'' of Theology, and in his critique of the dogmatic proofs of
the existence of God drives the separation of thinking and being to the extreme.

Among the three proofs, the ontological, the cosmological, and the
% ---- printed p.31 (PDF 39) ----
physico-theological, the emphasis must doubtless fall especially on the first.
``On the impossibility of an ontological proof of the existence of God'' the
following should here be brought out. The proof is in brief this: ``The most real
being is possible; among the many realities belongs also existence; it therefore
lies in the concept of the most real being that this being itself must necessarily
exist.'' One clarifies this by examples and dazzles by examples. When I suspend the
predicate in an identical judgment, and keep the subject, then quite rightly a
contradiction arises, and therefore one says that the predicate belongs to the
subject with necessity. But if I suspend the subject together with the predicate,
then no contradiction arises; ``\textit{denn es ist nichts mehr, welchem
widersprochen werden könnte} [for there is then nothing more that could be
contradicted].'' To posit a triangle and yet suspend its three angles is a
contradiction; but to suspend a triangle together with its three angles is no
contradiction. So it is with the absolutely necessary being. God is almighty;
almightiness cannot be suspended when you posit God; but when you say: God is not,
there is neither almightiness nor any other predicate. Being is, as copula in the
judgment, no real predicate. The actual contains in the concept no more than the
merely possible. A hundred actual dollars contain not the least more than a
hundred possible ones; for in the opposite case the concept would not be adequate.
A concept loses not a single one of its properties because its object does not
exist, and gains not a single property from its object's being found existing.

``On the so famous ontological (Cartesian) proof, then, all trouble and labor is
wasted; a person would, by mere ideas, become no richer in insight than a merchant
in fortune by adding a few zeros to his cash balance.'' From thinking to existence
there is no speculative transition.

\begin{center}{\bfseries 3)\quad Transcendental Idealism.}\end{center}

Since Kant expressly understands by ``a transcendental cognition'' one that ``does
not properly have to do with objects, but only with our cognition of objects,
insofar as this is to be possible \textit{a priori},'' it is already seen from
this why he himself calls his Criticism a transcendental philosophy. This
philosophy will now doubtless in its way assert the validity of the world of
experience; but so would
% ---- printed p.32 (PDF 40) ----
Berkeley's Spiritualism too, in its way. The manner in which the moment of
subjectivity becomes the overreaching thing in all our cognition characterizes the
Kantian transcendental philosophy as Idealism.

In what relation this transcendental Idealism has expressly set itself both to
Scepticism and to Dogmatism is expressed especially in ``the Doctrine of Method,''
which treats the pure reason's \emph{Discipline}, \emph{Canon},
\emph{Architectonic}, and \emph{History}.

``Were one,'' it says, ``to ask the cold-blooded David Hume, properly created for
the equilibrium of judgment, what then moved him, by scruples that only a
laborious brooding could devise, to undermine the imagination so comforting and
useful to men, that their rational insight suffices when the concept of the
highest being is to be determined, then he would probably answer: Nothing, but the
intention of bringing reason further in self-knowledge, and therewith a certain
displeasure at the compulsion one would impose on reason, in that one prides
oneself on it, and yet hinders it from making a candid confession about the
weaknesses that it can discover in itself only by serious examination.''

Unsatisfying as Scepticism is in and for itself, so justified is it over against
an uncritical Dogmatism that does not know what the understanding is capable of,
and has not tested the limits of the possibility of cognition according to
principles. ``\textit{Und so führt der Skeptiker, der Zuchtmeister des
dogmatischen Vernünftlers, auf eine gesunde Kritik des Verstandes, ja der Vernunft
selbst} [And thus the sceptic, the disciplinarian of the dogmatic rationalizer,
leads to a sound critique of the understanding, indeed of reason itself].''

In pure cognition as science, reason must on the contrary always proceed
dogmatically, i.e. conduct strict proofs from secure principles \textit{a priori};
but to conduct strict proofs from secure principles \textit{a priori} is not the
same as (\textit{Wolffs methodus mathematica} [Wolff's mathematical method]) to
imitate mathematics. Mathematics's thoroughness rests on definitions, axioms, and
demonstrations; philosophy's on explica-
% ---- printed p.33 (PDF 41) ----
tions, deductions, and acroamatic proofs. It does not befit philosophy
``\textit{mit einem dogmatischen Gange zu strotzen und sich mit den Titeln und
Bändern der Mathematik auszuschmücken, in deren Orden sie doch nicht gehört} [to
strut with a dogmatic gait and adorn itself with the titles and ribbons of
mathematics, in whose order it does not after all belong].'' The apodictic
propositions (be they now provable, or else immediately certain) divide into
Dogmata and Mathemata. A direct synthetic proposition from concepts is a
\emph{Dogma}, whereas a similar proposition by construction of concepts is a
\emph{Mathema}. Since there is now in the speculative use of pure reason not a
single Dogma, every dogmatic method too --- be it now borrowed from the
mathematician, or in another way become a peculiar mannerism --- is in itself
unfitting. ``\textit{Denn sie verbirgt nur die Fehler und Irrthümer, und täuscht
die Philosophie, deren eigentliche Absicht ist, alle Schritte der Vernunft in
ihrem klärsten Lichte sehen zu lassen} [For it only conceals the faults and
errors, and deceives philosophy, whose proper aim is to let all the steps of
reason be seen in their clearest light].''

\begin{center}{\bfseries \S\,7.}\end{center}

``In every experience two factors are set against each other: the I and the thing,
the intelligence and its object. If one abstracts from the I, then one gets a thing
\emph{in and for itself}, and must then grasp the representation as a product of
the object; if one abstracts from the object, one gets an I \emph{in and for
itself}, and must then hold to the self-active consciousness that produces its own
content. The former is \emph{Dogmatism}, the latter \emph{Idealism}. These two are
mutually incompatible; a third is not given: here one must choose. To judge between
Dogmatism and Idealism, let one note the following: the I occurs in consciousness,
the thing in itself is a fiction of consciousness. In that Dogmatism is to explain
the origin of the representation from the thing in itself, it performs the dubious
feat of going behind consciousness, and invoking being; but being effects only
being, not consciousness; therefore only that standpoint can be true which is
chosen, not in being, but in consciousness,
% ---- printed p.34 (PDF 42) ----
not in the object, but in the I itself.'' In this spirit J. G. Fichte (†1814) has
drawn a consequence of Kantianism and (\textit{Erste Einleitung 1797}) introduced
his ``\textit{Wissenschaftslehre} [Doctrine of Science].''

While Fichte dissolves ``the thing in itself,'' holds fast to the I, and carries
through \emph{subjective Idealism}, the Kantian doctrine is illuminated from an
opposite side, in that J. F. Herbart (Metaphysik) fixes attention on ``the
thing,'' grasps even the I as a thing for itself, a ``real with many properties''
(Psychologie), and in his way develops a \emph{critical Realism}. ``The material
of experience is the foundation of philosophy; but as a mere given it still stands
outside philosophy. What then is philosophy's first step, the act with which it
begins? Experience has forced upon us concepts that do not let themselves be
thought. They are doubtless thought in a way by everyday understanding; but it is
an obscure and confused thinking that does not see through the conflicting
distinguishing marks. The sharpened thinking, metaphysics's logical analysis,
discovers in the concepts of experience (space, time, motion, and so on)
contradictions, conflicting, mutually excluding marks. What is to be done here?
These concepts cannot be thrown away; they are given to us, and we have only the
given to hold to. Neither can they be kept, since they are unthinkable, in
themselves impossible. The only way out is: we must rework them. \emph{The
reworking of the concepts of experience is the task of Speculation.}''

Returning, after this indication of the Herbartian direction, to Fichte, we must,
in order to understand subjective Idealism, well imprint on ourselves the highest
principles of ``the Doctrine of Science''; for in these principles ``the Doctrine
of Science'' is theory of knowledge.

α) \emph{Thesis.} If one grants that the proposition $A = A$ is self-evident, then
one grants at the same time consciousness's power to posit something. One posits
not that $A$ is, but only that, if $A$ is, then it is $A$.
% ---- printed p.35 (PDF 43) ----
If we now, instead of looking at what is posited, look at the positing one, or
rather: at the identity of the positing one and what is posited, then the
proposition $A = A$ transforms itself into I $=$ I. Whereas instead of $A = A$ we
could not say that $A$ is, we can now, instead of I $=$ I, say: I am. This
self-grounded is the ground of all action in the human spirit, and thus the
original distinguishing mark of the activity of consciousness in and for itself.

β) \emph{Antithesis.} The proposition Non-$A$ not $=A$ is, as a free primal act,
according to its form likewise unconditioned, but on the contrary, according to
its content, its matter, conditioned, because the antithesis presupposes the
thesis, because when a Non-$A$ is to be posited, an $A$ must already beforehand
have been posited. By the first principle, $A = A$, the form of the primal act was
a positing; by the second it is a counter-positing. What Non-$A$ is, is not yet
known; only this is known, that Non-$A$ is the opposite of $A$. But now $A$ is
posited by the I; here only an I is presupposed: what is opposed to the I is thus
Not-I.

γ) \emph{Synthesis.} The third principle is, since it is determined by the two
foregoing, almost entirely capable of receiving a proof. The task for the act
expressed by the third principle is given with the two foregoing, but not the
task's solution. On the one side the I is completely suspended by the Not-I: the I
cannot be posited insofar as the Not-I is posited. On the other side the Not-I,
posited by the I, is itself an I, posited in consciousness. The I is thus not
suspended by the Not-I; it is suspended, and yet nevertheless not suspended. This
is the contradiction. For the solution of this contradiction we must seek to find
an $X$ in which the mutually excluding principles can be united. But how can being
and not-being, reality and negation, be united with one another without
annihilating one another? Answer: by mutual limitation! The $X$ asked for is thus a
barrier, and self-limitation the primal act of the I.
% ---- printed p.36 (PDF 44) ----
Insofar as the essence of I-hood, the pure I, is grasped as a principle that
develops itself in the empirical I's reciprocal action with a world of barriers,
subjectivity has to such a degree gained the preponderance that herein is
contained a natural summons to lay an emphasis of ideality on the objective too.

In his work ``\textit{Von der Weltseele}'' F. W. J. Schelling (†1854) has, still
with Fichte's presuppositions, sought to show how the concept of matter
necessarily proceeds from the nature and intuition of the human spirit. ``The mind
(\textit{das Gemüth} [the mind]) is the unity of an unlimited and a limiting force.
Absence of barriers would, as well as absolute limitedness, make consciousness
impossible. Only in that the force that strives out into the unlimited is limited
by the opposite, while the limited one is conversely freed from its barriers, is
feeling, observation, cognition thinkable. A corresponding thing is seen in nature.
Not matter as matter, but the forces whose unity matter constitutes, is here the
first. Matter is the ever-becoming product of attraction and repulsion. Matter is
what it is by force; but force is the immaterial in matter, force in nature is that
which can be compared with spirit. Nature and spirit must originally be united in a
higher identity. The spirit's organ for the apprehension of nature is intuition; it
is intuition that takes into possession the space bounded and filled by attractive
and repulsive forces.''

With these thoughts Schelling has already trodden the objective ground of natural
philosophy. ``All knowledge,'' it says (\textit{System des transcendentalen
Idealismus}), ``rests on the agreement of a subject with an object. The aggregate
of the merely objective is nature; the aggregate of the merely subjective is the I
or the intelligence. For the union of both sides two ways are possible; either one
makes nature the first, and asks how the intelligent comes to it,
% ---- printed p.37 (PDF 45) ----
i.e. one strives to interpret it in pure forms of thought --- natural philosophy;
or one makes the subject the first, and asks how the objects proceed from the
subject --- transcendental philosophy. All philosophy must aim either at making an
intelligence out of nature, or a nature out of the intelligence. Just as
transcendental philosophy has to subordinate the real to the ideal, so natural
philosophy must seek to explain the ideal from the real. But both are only the two
poles of one and the same knowledge, which mutually seek each other; therefore one
must also, when one starts from the one pole, necessarily come to the other.''

``All the parts of philosophy must be presented in one continuity, and the whole
of philosophy must be presented as what it is, namely as the continuous history of
consciousness, for which what is deposited in experience serves only, so to speak,
as monument and as document.''

Schelling's transcendental Idealism can now be regarded as forerunner of Hegel's
``\textit{Phänomenologie des Geistes} [Phenomenology of Spirit]''; what Schelling
promised, Hegel (†1831) fulfilled.

``Consciousness has in general, according as the object shows itself, three stages.
The object is namely either an object standing over against the I, or it is the I
itself, or something objective that nevertheless belongs just as much to the I, to
thought.'' The three stages are thus: Consciousness, Self-consciousness, Reason.

a) \emph{Consciousness.} At the first stage consciousness is sensing, observing,
understanding.

α) \emph{The sensing consciousness is immediate certainty of an external object.}
The expression for such an object's universality is, \emph{that it is, and is this,
now according to time, here according to space}. The object is thereby wholly
different from all other objects, and for itself sufficiently determined. But this
Now is no more, another has stepped into its place, which however is also a Now:
the Now
% ---- printed p.38 (PDF 46) ----
is, in its vanishing, something abiding, namely the universal. This Here, which I
mean and point out, has a right, a left, an above, a below, a behind, a before, to
infinity. This Here is thus not the single, immediately determined, but: now here,
now here, now here! an aggregate of Many, and in these Many nevertheless the One,
thus the universal.

β) The universal, insofar as it still hovers between sensuousness and reflection,
is object for the \emph{observing} consciousness; \emph{here is the thing with its
properties}. The sensible properties are at once immediately each for itself, and
yet in relation to one another, belong to the thing as this single thing, and are
yet so universal that they, separated from one another, can be found in manifold
other things. Thing and properties change, and yet have in this change their
subsistence: the changeable is an outer, the subsisting an inner.

γ) Insofar as the object is an inner with its outer, a force in its expression, it
is there for the \emph{understanding} consciousness. The understanding sees in the
changing outer of things a series of phenomena whose inner consists in a system of
abiding relations: the laws. Things are under the laws, and the laws are the laws
of things; things and laws are two sides of the same: the True, the essence of
things, is thought, the concept.

b) \emph{Self-consciousness.} In that consciousness turns toward thought as
thought, toward essence as essence, toward the inner as inner, it gets its own
reflection and thereby itself as object. Self-consciousness has in its formation or
movement three stages: desire, the relation of ruling and serving, universal
self-consciousness.

α) \emph{In desire self-consciousness relates to itself as single, solely existing
self-consciousness}; it is a selfless object that is to satisfy the craving; but,
in that the object is consumed, and satisfaction sets in,
% ---- printed p.39 (PDF 47) ----
consciousness has in its self-feeling understood that it, as subjective, must at
the same time essentially be objective.

β) The concept of self-consciousness as a subject that is at the same time an
object finds its complete expression in this, that self-consciousness sees in its
object another self-consciousness. In that the two self-consciousnesses set over
against each other are each according to their concept independent, the tension
arises between being master and being servant. The servant is selfless, insofar as
his self is a moment in the master's consciousness; but the master too is selfless,
insofar as his being master rests on his having a servant; he is master only by
means of the servant. With this reciprocity self-consciousness makes the transition
to articulated universal will, to positive freedom.

γ) \emph{In universal self-consciousness the individual knows himself} as serving
the will of the many, in that he essentially serves himself. Here there is not a
sum of wills, but the many wills are moments in the universal will. This
self-consciousness is the foundation for all virtues: those of love, honor,
friendship, bravery, sacrifice, distinction, and so on.

c) \emph{Reason.} Reason is the highest union of consciousness and
self-consciousness in knowledge of the object and knowledge of itself. Reason is
the certainty that its determinations are just as much objective, determinations
of the essence of things, as that they are the subject's own thoughts. Fichte's
``I'' and Herbart's ``real'' have hereby become moments; reason's knowledge is not
merely subjective certainty, but truth; but truth is unity of certainty and
objectivity, of knowledge and being.

``The goal, absolute knowledge, or the spirit knowing itself as spirit, has for its
path the recollection of the spirits as they are in themselves, and accomplish the
organization of their realm. Their preservation, in respect of their free existence
appearing in the form of contingency, is history, but, in respect of their
conceived organization, the phenomenal
% ---- printed p.40 (PDF 48) ----
knowledge's science; both together, the conceived history, form the recollection
and the Golgotha of absolute spirit, the actuality, truth, and certainty of its
throne, without which it would be the lifeless, solitary; only out of the chalice
of this realm of spirits foams up for it its infinity.''

\begin{center}\rule{2.5cm}{0.4pt}\end{center}

\begin{center}
  {\large\bfseries C.}\par\vspace{2pt}
  {\bfseries Factual Knowledge.}\par\vspace{4pt}
  \rule{2cm}{0.4pt}
\end{center}
\phantomsection
\addcontentsline{toc}{subsection}{\texorpdfstring{C.\quad Factual Knowledge}{C. Factual Knowledge}}

\begin{center}{\bfseries \S\,8.}\end{center}

\noindent If we now look back on the theories of knowledge dealt with here, the
sensualistic as well as the idealistic, then it is quite striking to us that not a
single one among them all satisfies ``sound understanding'' in its consequences,
while they nevertheless in their cooperation --- insofar as each on its own most
forcefully asserts its side of the matter --- as it were unite in giving ``sound
human understanding'' its right.

From this it cannot now be concluded that ``sound understanding'' is wiser than
philosophy, but only that it, with the sureness of instinct, holds fast to the
members whose mutual relation philosophy wishes to explain. When philosophy, then,
instead of clearing up the relation of the members, explains away or bungles an
essential member, ``sound understanding'' senses mischief, and protests.

``Sound understanding'' readily grants that one cannot become conscious of a
surrounding world without taking up its impressions and corresponding images into
one's circle of representation; but it maintains at the same time that the world we
become conscious of must, as truly as it is a real world, have a real existence
independent of our consciousness. The idealist may now exert himself to prove the
indissoluble unity of world and consciousness as much as he will: ``sound
understanding'' pro-
% ---- printed p.41 (PDF 49) ----
tests. That our sensations and representations proceed from immediate
sense-impressions, the understanding will readily grant; but when the sceptic
draws from this the consequence that cause and effect exist only in our
imagination, that the laws are not laws, ``sound understanding'' stands still, and
begins to protest.

A philosophy that is in open contradiction with ``sound understanding'' is false;
a philosophy that immediately conforms to ``sound understanding'' is superficial.
Philosophy will not arbitrarily contradict the natural understanding, but only
show that the so-called sound sense, without having any inkling of it, incessantly
contradicts itself.

``Sound understanding'' demands that I should immediately become conscious that
there is something outside my consciousness; but in becoming conscious of
something, it is of course in my consciousness; insofar as it is outside, I cannot
become conscious of it: thus I am to become conscious that what is in my
consciousness is not there; what a contradiction! When the sun shines, and the
trees are mirrored in the stream, ``sound understanding'' means that the light is
brightness, whether or not there is anyone who sees it, and that the trees really
are mirrored in the stream, whether or not this mirroring is reflected in a seeing
eye. But the question is how one is to convince oneself of this. By sight? That is
a contradiction, since here it is a matter of the phenomenon, not insofar as it is
seen, but insofar as it is not seen. By thought? That is a contradiction; the
phenomenon does not let itself, as phenomenon, be thought, for in that case the
one born blind could also think it. By the testimony of others? An empty evasion;
since here it must go for the one as for the other. Only for the seeing one is
light brightness; only in that the brook and the tree are mirrored in the eye can
the tree itself be mirrored in the brook.

It falls naturally that one thinks back to the time when he did not yet exist; and
just as naturally that he thinks of the moment when he will again have ceased to
be.
% ---- printed p.42 (PDF 50) ----
Nonetheless there is here in both a strange contradiction. That consciousness can
be conscious of its own existence is evident; but how can the existing
consciousness become conscious of itself as non-existing? In reality? No, that is
a contradiction, for when consciousness really is, its non-being is really
excluded; and, when consciousness is not, it cannot become conscious of anything.
In fantasy? But when consciousness sees its non-being only as a fantasy, how can
it then convince itself that the real either has corresponded or will correspond
to this fantasy? ``Sound understanding'' may readily advance with its explanation,
but the matter is that the explanation contains just as many contradictions as
that which it is to explain.

What is truth? Agreement between what consciousness takes to be real, and this
real itself. And untruth? Naturally, conflict between what is posited in
consciousness and what is in reality. If now all reality were at once gone, and
the empty consciousness with its fantasies and thoughts stood alone behind, it
would be able to carry on the most arbitrary play with figments. That truth is
gone when the real world-content is gone, sound understanding will surely readily
% [errata applied (Rettelser p.42): print ``the arbitrary''; read ``the real'']
grant; but the real as merely real is, when consciousness itself is gone, not the
True either. A stone, a star, any real object whatever, is in and for itself
neither something true nor something untrue. In order that the difference between
true and untrue can arise, a judgment must be passed on the real object; but in
order that a judgment can be passed, there must absolutely be a judging
consciousness. But if consciousness is an indispensable requisite for truth, the
real for itself cannot possibly be the True; how then can ``sound understanding''
suppose it can think a time when there were in \emph{truth} real atoms, real
elements, a real earth, but yet
% ---- printed p.43 (PDF 51) ----
no, no consciousness at all? Real objects that could not possibly be objectified;
what a contradiction!

The outer, real world is, in its pure unconscious reality, a nothing; if anyone
will maintain the opposite, he must maintain it by force of consciousness, and
thus contradict himself.

What is an Eternal? When a universally valid proposition is known, it does not
look at all as though the True first came into being in the moment of cognition;
the True shows itself on the contrary to consciousness as that which has always
been. That every thing is what it is; that two magnitudes, each of which for
itself is equal to one and the same third, are mutually equal; that a ground must
necessarily have its consequence, an effect its cause — all these so-called
``eternal truths'' present themselves to us as those that have been prior to all
finite knowledge; and yet they are in themselves nothing when one takes knowledge
away: how can the consciousness that comes to be in time become conscious of the
non-becoming, the eternal?

What does contradiction mean? An existing nothing is a contradiction; a round
square is a contradiction; a flying horse is also a contradiction. Nonetheless
these contradictions are of different kinds. A flying horse one can represent to
oneself, but not think. The unthinkable lies not straightforwardly in the logical,
that the predicate ``flying'' contradicts the concept ``horse,'' for this concept
could of course be extended; but in this, that the whole concept thus extended
stands in conflict with the conditions and laws of experience. It is otherwise
with the round square; here the contradiction is seen immediately in the
representation itself. To represent the square as round is impossible: to think a
round square is, on the other hand, possible, insofar as it is possible to think
the predicate's \emph{conflict} with the subject. The contradiction is thus not
merely that the predicate's conflict with the subject makes the judgment
meaningless; but the contradiction is in a higher sense this, that the
contradiction itself, this unthinkable, nevertheless lets itself be thought. This
% ---- printed p.44 (PDF 52) ----
becomes especially clear when the task is to think the existing nothing. It is
indeed nothing, but in that consciousness posits this negation, nothing has
precisely existence in thought, thus a thought existence.

By pondering a little on these many contradictions ``sound understanding'' will
best learn to understand that the eager strife the philosophers carry on about
``the principle of contradiction'' is more than a strife about words.

Herbart and Hegel are both agreed that the concepts of experience contain a
contradiction; they are furthermore agreed that the contradiction is to be
suspended; but about the relation of the contradiction to the essence of things,
indeed about the very essence of the contradiction, they are in the highest degree
at odds. In that Herbart holds fast to the principle of contradiction, he sees in
the contradiction an error that one is to avoid, a delusion that one is to remove.
If the concepts of experience contain contradictions, then the fault lies not in
the objective nature of things, but in the knowing subjectivity. The thinking
consciousness must therefore, by reworking the concepts of experience and doing
away with the contradiction, see to remove the false apprehension. Zealous against
the contradiction, Herbart even accuses the Hegelian philosophy of Empiricism,
since it ``takes up the concepts from experience unchanged, and, despite the
insight into their contradicting nature, precisely because they are empirically
given, regards them as justified, indeed even for their sake remodels logic.''
Hegel, on the contrary, with his ``\textit{Alles ist ein Widerspruch} [Everything
is a contradiction],'' lays the contradiction into the essence of things, so that
a being in finitude is necessarily a being in contradiction (solution of the
Kantian antinomies). On contradiction rests not only the pain of life, but also
the drive that drives existence forward: contradiction calls forth motion, and
motion is in turn the solution of the contradiction. As little as existence can
exclude contradiction, so little can thought exclude it: the pulse of thought is
dialectic, the contradiction the nerve of dialectic.
% ---- printed p.45 (PDF 53) ----
\begin{center}{\bfseries \S\,9.}\end{center}

How intimately contradiction is united with thought is seen in the logical
judgments. The word ``leaf'' can, as a word in language, point to a sensibly given
in its immediate singularity, but also, as a conceptual expression, designate
something universal. An exclamation such as ``this leaf!'' is in linguistic form
imperfect, and gives a definite meaning only when someone at the same time points
to a real leaf. If the eye is now fixed on ``this leaf,'' thought can be united
with the sensible, in that a judgment is uttered: ``this leaf'' is green. But
``this leaf'' is not green, for it is yellow; ``this leaf'' is thus green and yet
not green; ``this leaf'' is both green and yellow; here is a swarm of
contradictions. In the first place the predicate ``green'' is, in relation to
``this leaf,'' something universal, whereas ``this leaf'' itself is something
single. The judgment thus states that the single is a universal; but the single is
nevertheless not a universal, but a single: here is a contradiction. What then is
``this''? ``This'' is a leaf, ``this'' is a table, ``this'' is a horse; ``this''
is thus not ``this,'' but something other than ``this'': sheer contradictions.

Through the history of philosophy runs a series of controversies that all point to
the contradiction present in the matter. In Stilpo's utterance: ``the cabbage that
is offered for sale here, is not; for cabbage existed already 1000 years ago,''
lies the meaning that the individual in its sensible immediacy cannot be thought.
The Scholastics' controversy about the relation between \textit{res} [things] and
\textit{universalia} [universals] points to the same contradiction, which also, as
regards the more recent philosophy, repeats itself in the opposition between
Sensualism and Idealism.

Locke and Hume are both attentive to the delusion that can arise in this, that the
universal is held fast for itself in the word; and, since the objective Idealism
had advanced so far that Hegel's Dialectic was assumed to have transformed
% ---- printed p.46 (PDF 54) ----
the sensible into a moment vanishing in the concept, Feuerbach (\textit{Philosophie
der Zukunft}) came forward with a critique that demanded unconditional Sensualism.

For the understanding of Idealism's one-sidedness we bring out in this connection
the relation between the sensible and the thinkable, between the individual and
the universal, between the real and the essential.

\emph{a) The sensible and the thinkable.}

That the sensible is an other of thought, a barrier of reflection, Idealism must
of course grant; but, in that reflection shows that thought precisely needs this
barrier, the otherness, the barrier itself, is nevertheless thought after all, and
the sensible, the Negative, which thought thinks in thinking itself, is thereby
melted down into a moment in the concept. This ambiguity can for a part be veiled
in the pure, objective Logic, where ``matter,'' this illogical, nonetheless by the
help of Dialectic at last becomes just as logical as ``form''; but in ``natural
philosophy,'' where reflection strikes against the fact, the error is laid bare.

``By the old, general thought that every body consists of four elements, or the
more recent Paracelsian one that it consists of mercury or fluidity, sulphur or
oil, and salt (Jacob Böhme called it the great Trinity), and by many other
thoughts of this kind, the refutation has in the first place been made easy, in
that under those names one would understand the individual, empirical materials
that are most nearly designated by these names. But it is not to be mistaken that
they should much more essentially contain and express the conceptual
determinations. One must therefore much rather admire
% [print: ``algemeine'' for ``allgemeine'' — reproduced as printed]
\textit{die Gewaltsamkeit, mit welcher der Gedanke, der noch nicht frei war, in
solchen sinnlichen besondern Existenzen nur seine eigene Bestimmung und die
algemeine Bedeutung erkannte und festhielt; er darf darum auch nicht auf
experimentirende Weise widerlegt werden:} [the violence with which thought, which
was not yet free, recognized and held fast, in such sensible particular
existences, only its own determination and the universal meaning; it may therefore
also not be refuted in an experimental way:]'' (Hegel, Naturphil. S. 271).
% ---- printed p.47 (PDF 55) ----
This apprehension is quite in the spirit of Idealism, and that Hegel too will not
have it restricted to an assessment of that thought ``\textit{der noch nicht frei
war} [that was not yet free],'' but even expressly applied in favor of his own
interpretation of the phenomena, is evident from the way in which he himself, e.g.,
introduces the ``lunar'' and the ``cometary'' principle; for precisely in this
``lunar'' and this ``cometary'' is seen ``\textit{die Gewaltsamkeit, mit welcher
der Gedanke in solchen sinnlichen besondern Existenzen} [the violence with which
thought, in such sensible particular existences],'' recognizing and holding fast
only ``\textit{seine eigene Bestimmung und die allgemeine Bedeutung} [its own
determination and the universal meaning].''

The fundamental error (πρῶτον ψεῦδος) lies in the misprision of the sensible. The
sensible is by no means suspended because it is grasped as thought's other; the
element of sensuousness is more than a logical negation, it is for thought the
real barrier, the in-itself-unthinkable, which is therefore to be grasped and
determined in relation to the immediate sense.

``According to the moment of its `being-in-itself,' matter would be the single
point; according to the moment of its `being-outside-itself,' it would most nearly
be a multitude of mutually excluding atoms. But, in that these, by excluding one
another, relate to one another, the atom has no reality, and the atomistic as well
as absolute continuity or infinite divisibility is only according to possibility
present in it.'' In this Hegelian apprehension it does indeed at first glance look
as though thought penetrated matter. By the logical opposition of
``\textit{Insichseyn} [being-in-itself]'' and ``\textit{Außersichseyn}
[being-outside-itself]'' and its dialectical treatment one does indeed escape the
difficulties that Atomism, this stumbling-stone of Speculation, otherwise brings
with it when it is to be clarified in a way satisfying to those versed in the
matter; but the difficulty is only got around, not suspended: in that the element
of sensuousness asserts its right, the relation between matter's discretion and
continuity takes up an unthinkable, a for
% ---- printed p.48 (PDF 56) ----
reflection impenetrable; for the sensible is thought's barrier.

To convince ourselves from yet another side how unthinkable the sensible is, one
need only cast a glance at the Herbartian view. With objective Idealism, Herbart's
Intellectual-Realism namely has this in common, that it, though in its way, will
make the sensible thinkable. Immediately apprehended, the sensible is only a
semblance; but this semblance presupposes an in-itself-being. This being, an
independent reality resting on itself, can be known only by thought. The otherwise
material atom is here transformed into an atom of thought. Of such peculiar,
positive, unchangeable thought-atoms, qualia, there is a great multiplicity.
``Just as well as a mathematical point can'' — in spite of \textit{principium
indiscernibilium}! — ``coincide with another at the same place, just as well can
indeed these absolutely positive thought-atoms penetrate one another, and coincide
in one and the same point;''
% [errata applied (Rettelser p.48): quotation closes at ``…same point;''; the
% following clause is Nielsen's own, standing outside the quotation]
but such thought-atoms are then round squares: the sensible does not let itself be
dissolved.

\emph{b) The individual and the universal.}

In ``\textit{Phänomenologie des Geistes},'' ``this Here'' and ``this Now,'' which
involuntarily designated the punctual, were transformed by reflection into ``the
universal,'' and the punctual suspended. But, in that by this ``Now!'' this
``Here!'' this ``There!'' one points to a sensible, the universal is again
suspended, and the punctual comes into force. ``This here!'' is of course not the
same as ``this there!''; the universal is in the word, but the word points, by
this pointing, out beyond itself to the individual. It is with respect to this of
significance that Locke and Hume have investigated the origin of concepts.

In the universal concept the representation vanishes; with the determination of
the individual it returns again:
% ---- printed p.49 (PDF 57) ----
to mistake the representation is to disdain the fact, and to mistake the
individual.

In the working-up of the element of representation two factors are active: the
productive imagination and the reflecting thought. The involuntary representing
forms the image in agreement with the object; hereby is conditioned the factual
agreement between the circle of representation and the given surrounding world.
The free imagination, on the contrary, brings about disagreement between reality
and consciousness, in that it in various ways transforms the representations. If
we now ask whence the thought, restlessly working on the dialectic of its
concepts, gets its most pliable, most spiritual, and thus most convenient
material: from the factual individual-determinations, or from the representations
transformed by imagination and the freely hovering universal-images? the answer
cannot be doubtful.

To what degree the objective Idealism has mistaken the element of representation
by exchanging facticity's images for hovering intuitions can be seen from many
examples. ``The first qualified matter is matter as pure identity with itself, as
unity of reflection into itself; to that extent it is only the first, itself still
abstract manifestation. Existing in nature, it is the relation to itself as
independent against the other determinations of the totality. This existing,
universal \emph{Self} of matter is \emph{Light}, as individuality the
\emph{Star}, and this as moment in a totality, the \emph{Sun}.'' (Hegel Naturphil.
S. 130).

So long as this intuition of the nature of light keeps itself in pure generality
far from all finite facticity-determinations, one may readily grant it an
ingenious application of logical categories; but if, on the contrary, a connection
to definite phenomena is demanded, the whole apprehension will show itself so
hollow and hovering that here at no point is there even a possibility of a
penetration into the factual.
% ---- printed p.50 (PDF 58) ----
Where Speculation enters more closely into a characterization of given objects, and
thus engages with the individual, the misunderstanding becomes still more
palpable. ``Clay-earth is, like silica, the immediate, simple, unopened concept,
thus the first different earthy — the possibility of combustibility\ldots, talc or
bitter-earth is the subject of salt; hence comes the bitterness of the sea. It is
an intermediate taste that has become fire-principle\ldots.'' (Naturphilosophie S.
416).

Instead of developing the concepts by sharp analysis of the individual, and
reducing the particular in such a way that the content-rich universal can find its
determinations confirmed in the given, Idealism commits the blunder of letting the
universal-forms move at such a distance from the facts that these, like phenomena
in the distant horizon, take on hovering outlines, and, according as ``the
concept'' finds it serviceable, become either rocks, or animals, or clouds.

\emph{c) The real and the essential.}

That one can be right in something essential and yet in reality wrong is just as
comprehensible as that something can be true in general, and yet in its
application to the single become untrue. That thinking is identical with being is
essentially true; but that our conscious thinking and the finite existence of
things should be identical is really untrue; that what exists must express the
category of existence is essentially true; but that it is only the category of
existence that exists is really untrue. Those who speak of essence, and those who
speak of the real, speak a different language.

``The universe is to be thought under the image of a line, at whose midpoint falls
$A = A$, at whose ends falls on the one side $\overset{+}{A} = B$, i.e. an
overreaching of the subjective, on the other side $A = \overset{+}{B}$, i.e. an
overreaching of the objective, yet in such a way that relative identity takes place
% ---- printed p.51 (PDF 59) ----
even in these extremes. The one side is the Real or Nature, the other the Ideal.
The real side develops itself according to three potences, in that potence
designates a definite qualitative difference of subjectivity and objectivity: the
first potence, $A$, is matter and the force of gravity — the greatest preponderance
of objectivity; the second potence, $A^2$, is light, an inner — as gravity an
outer — intuiting of nature. Light is the higher expression of the subjective. It
is the absolute identity itself. The third potence is the common product of light
and the force of gravity, the organism, $A^3$. The organism is just as original as
matter.''

If now a strict empiricist, a man who knows of no speculation, suddenly came to
hear this Schellingian explanation of nature, he would undoubtedly regard the whole
as a figment of the brain. Hereby, however, it is not proved that what is stated in
those speculative propositions is in itself meaningless; hereby it is only given
that the empiricist does not understand the language. If someone who would explain
the world's maintenance suddenly seized Lagrange's
% [print: a small dot precedes ``da'' in the numerator — possibly a differential dot]
formula: $\int \dfrac{da}{dt}\,dt = \Sigma\,\dfrac{p}{in-i'n'}\,\sin(in-i'n'+q)$,
with the assurance that the solar system is maintained solely on the ground that
$i$ and $i'$ can never simultaneously be equal to $0$, so that the denominator
$in-i'n'$ can never become $0$, because $n$ and $n'$ are incommensurable; then even
those who otherwise think rationally about real things, but have no concept of
mathematics, would judge this mathematical explanation just as severely as the
empiricist that idealistic one.

That ignorance and lack of philosophical sense unite to tear down Idealism means
nothing; in order to understand the philosophical language, one must just as surely
study philosophy as one must, in order to understand the mathematician, study
mathematics.
% ---- printed p.52 (PDF 60) ----
But even when we seek to understand philosophy philosophically, it cannot be denied
that the objective Idealism has, with its profound cognition of essence, sinned
against reality in a twofold way, partly by abstracting from real points of
departure, partly by wishing to force a real content into unreal concepts. ``The
Newtonian theory, according to which light is to spread in lines, or the
wave-theory, according to which it is to spread wave-like, as the Eulerian ether or
as the trembling of sound, are,'' says Hegel (Naturphil. S. 141), ``material
representations, ``\textit{die für die Erkenntniß des Lichts nichts nutzen} [that
are of no use for the knowledge of light].''\,'' Hereby, then, there is abstraction
from all real points of departure, and every controlling investigation made
impossible.

Instead of entering into the course of the real investigations, and by the
concept's dialectical superiority leading the empiricist, by force of his own
acknowledged presuppositions, to a cognition of the essence of the matter, this
Idealism has not only, with an arbitrary apprehension of such phenomena as
refraction and colors, but from first to last, wished to correct the expressions
and procedure of Empiricism, without, however, rightly grasping, let alone seeing
through, the problems of Empiricism.

\begin{center}{\bfseries \S\,10.}\end{center}

When it is a matter of a cognition of the real, the theory of experience must have
the decisive voice. Humean Scepticism makes, in its zeal for experience, true
experience impossible. By proving the — though only subjective — universal validity
of the a priori forms, Kant now doubtless brings philosophy into closer
understanding with Empiricism. But, although the theory of experience to a certain
degree readily grants Kant the impossibility of knowing the interior of things, it
is far from grasping time and space as merely being in our consciousness. When the
astronomer speaks of the moon's distance from the earth, he surely does intuit this
distance from the earth, not as a merely subjective mode of apprehension,
% ---- printed p.53 (PDF 61) ----
but as a real, objective spatial relation; when it is said that Neptune perturbs
Uranus, the reciprocity between the planets is surely not sought in our
understanding alone, but in the nature of things. The separation Kant makes between
the phenomena of experience and the things themselves is a schoolroom theory that
has not the least significance in the theory of experience. By asserting the
objective validity of the a priori forms, Hegel, on the contrary, has — despite all
his offenses against reality — in this respect pointed out the essential agreement
of philosophy and Empiricism.

If this agreement is now also to be carried through in the form of the concept from
the standpoint of reality, the rational and yet empirical connecting links between
the a priori forms and the immediate facts must let themselves be pointed out. It
is not enough that a universally valid relation between ground and consequence,
conditions and the conditioned, cause and effect, and so forth, is acknowledged a
priori; the task is to assure oneself of the real relations.

The relation between cause and effect is a relation of necessity that has its
abiding, universal expression in law. The cause for itself is not a law, though it
is an a priori law that everything must have a cause; the effect for itself is
naturally not the law either; but the law is that the effect must set in when the
cause takes place. On these a priori presuppositions rest observation and
experiment.

When a real cause, a real ground-relation, a real law, is to be found out, one must
begin with an assumption. That the earth's attraction is the cause of the stone's
fall, that a liquid's evaporation is the cause of cold, that every body has its
definite melting point, that ebb and flow are owed to the influence of sun and
moon, all such assumptions point originally back to a subjective estimate, to
impressions of experience, which, however often they may be repeated, cannot
possibly prove the necessary, or establish the universal.
% ---- printed p.54 (PDF 62) ----
In order then to see how far the theory of experience is, by its ``observations''
and ``experiments,'' able to secure the universally valid, we consider especially:
abstraction, induction, and the real law.

\emph{a) Abstraction.}

The abstracted reminds of that from which it is abstracted; thus pure space reminds
of matter, and demands, so to speak, its lost content back. To space's discretion
correspond matter's atoms, to space's continuity matter's mass. The
universal-determinations that lie in geometry's lines, surfaces, and bodies have
validity too insofar as they are stated of the material. That the straight line is
the shortest path between two points, that two straight lines can cut each other
only in one point, and so on, are propositions that can be abstracted from the
material, and applied to the material. Insofar as matter is grasped as the
abstract, the moment of experience here coincides with the a priori.

But even where the moment of experience steps forth more strongly — as when the
space-filling matter is ascribed a force of resistance, a force of expansion,
gravity, and so on — we are by no means, as Hume supposes, merely referred to the
association of ideas and habit. ``All bodies have gravity, all bodies have force of
cohesion, all bodies can be divided'': a simple denial of these propositions'
universal validity has nothing to signify; they are supported by millions of
impressions of experience. These propositions are indeed in their generality
abstract, but this abstract is a reflex of the concrete; he who therefore will deny
the universal validity must expressly adduce facts, and take upon himself the
\textit{onus probandi} [burden of proof].

What the incessant reciprocal determination of the element of experience and the a
priori signifies is shown by physics. ``It is important that the application of the
most simple mathematical truths in natural science be carefully justified, so that
% ---- printed p.55 (PDF 63) ----
no doubt is left about nature's agreement with thought'' (H. C. Ørsted); and it is
likewise important that the propositions of natural science be grasped
mathematically; for that which especially furthers the strict investigations, in
which the consequence becomes recognizable on the presupposition, and this in turn
on that, is precisely the exact form in which mathematics (Cf. \emph{Philosophie og
Mathematik} pp. 12--14) grasps the functions in question.

The difference between the abstract, philosophical and the mathematical-physical
atomic theory then also consists essentially in this, that the latter at once
grasps reality in its punctuality, whereas the former never gets beyond the
one-sidedly universal. Mathematics too doubtless begins with abstractions, but it
shows at the same time how the abstractions can be integrated. The mathematician
speaks not merely of atoms and the atoms' union in molecules, but formulates
besides the expression in such a way that it precisely suits an investigation of
the real. Let the number of molecules be, e.g., $m$, i.e. the mass, $v$ the volume
of the mass, which encloses a number of molecules; one then has in the formula
$\dfrac{m}{v} = \varrho$ an exact expression for density, and herein a
presupposition for the determination of the specific gravities of the various
bodies.

If now $v$ encloses a number of $n$ molecules, one gets $\dfrac{v}{n} = e^3$, and
$e^3$ thus as the calculating expression for the molecule. If further, in a body,
$m$, the finite part of the mass, becomes infinitely small, thus: $dm$, then one
obtains $\dfrac{dm}{dv} = \varrho$, whereby then also the finite sum of the parts,
$\Sigma$, passes over to $\int$, the infinite sum; one thus has for the
\emph{homogeneous} body, in that $\varrho$ is a function of the various places in
the mass, the formula: $\varrho\int dv$, and the formula: $\int \varrho\, dv$ for
the \emph{heterogeneous},
% ---- printed p.56 (PDF 64) ----
and herein a fundamental concept, from which countless specifications agreeing with
experience let themselves be derived.

The mathematical punctuality makes possible an analysis of reality that in
sharpness can measure itself with Dialectic's analysis of essence.

\emph{b) Induction.}

In abstraction the individual case is disregarded; in induction the individual gets
its significance again. How the individual members lead over to the universal,
insofar as this has existence in them; how the parts by their peculiarity determine
the whole, so that the essence of the whole finds its complete expression in the
reality of the parts — this is shown by induction. ``It is the remarkable thing
about this kind of inference that the conclusion is never a singular or particular,
but always a \emph{universal} judgment. The categorical inference infers from a
rule to a case; the disjunctive, on the contrary, from the cases to the rule. There
the rule is what is presupposed, here what is attained by inference. Categorically
one can only partially, hypothetically only negatively (\textit{modo tollente} [by
the mode of denying]), infer from the particular to the universal. The disjunctive
mode of inference is the only one that permits an affirmative and at the same time
universal inference from the particular to the universal. Therein lies the magic
power of induction: it lets us, from the juxtaposition of observations and facts,
know the law.'' (\textit{Die Theorie der Induction von E. F. Apelt. Leipzig 1854. S.
18}).

With respect to the development and application of induction, Socrates, Aristotle,
and Kepler can be said to designate the turning-points.

``Two things,'' says Aristotle, ``can no one with fairness dispute to Socrates,
induction and the determination of the universal concepts.'' How Socrates
proceeded, in that by consideration of a multitude of individual cases and
% ---- printed p.57 (PDF 65) ----
examples he sought to clarify the ethical concepts, is well known. But, although
this Socratic ἐπαγωγή is doubtless a progression from the individual and particular
to the universal, it is nonetheless — as Apelt too remarks — not what we now
understand by induction, but rather an analyzing art of abstraction.

As logician and natural philosopher, Aristotle is called upon to give induction a
more definite form, and then sets ἐπαγωγή partly against συλλογισμός, partly against
ἐπιστήμη. In the first case ἐπαγωγή is the inference from the particular to the
universal, from the parts to the whole, from the cases to the rule; in the second
case ἐπαγωγή also designates the beginning of the universal (ἀρχή ἐστι καὶ τοῦ
καθόλου); whereas ἐπιστήμη only draws out rational inferences from the universal.
Ἐπαγωγή leads to the principle, ἐπιστήμη develops its content from the principle. To
this Apelt now remarks: ``Induction (ἐπαγωγή in opposition to συλλογισμός) never
leads to the principle, one presupposes it already; it leads rather to a law that
ties a definite circle of phenomena to a principle. The regressive, epagogic train
of thought that leads to the principle (ἐπαγωγή in opposition to ἐπιστήμη) is no
inferential connection, but an abstraction. This confusion of induction and
abstraction has in the history of philosophy remained standing until recent
times.''

As the one who in recent times has rightly grasped and scientifically applied
induction, Kepler is named; for Bacon of Verulam only emphasizes its significance,
but does not know its application. Kepler's discovery of the Mars orbit happened by
application of induction; the solution of this task Apelt treats with great acumen
and expertise. ``The art consists in finding out, from a series of numerical values
for two connected variable magnitudes, the law of their mutual dependence on one
another, i.e. the form of the function. Herein lies the great difficulty of the
task.''
% ---- printed p.58 (PDF 66) ----
The logical form of the Keplerian induction is this: ``The Mars orbit contains in
itself the Mars positions:
\[ O_1,\ O_2,\ O_3,\ \ldots O_n. \]
At the position $O_1$, i.e. at the time when Mars's true anomaly is $v_1$, the
magnitude of the radius vector is $r_1$
\[
\begin{array}{ccccccccc}
\text{,,}&\text{,,}&\text{,,}& v_2 &\text{,,}&\text{,,}&\text{,,}&\text{,,}& r_2\\[2pt]
\text{,,}&\text{,,}&\text{,,}& v_3 &\text{,,}&\text{,,}&\text{,,}&\text{,,}& r_3\\[2pt]
\multicolumn{9}{c}{\vdots}\\[2pt]
\text{,,}&\text{,,}&\text{,,}& v_n &\text{,,}&\text{,,}&\text{,,}&\text{,,}& r_n\\
\end{array}
\]
The Mars orbit is a curve whose polar equation is: $r = \dfrac{a(1-e^2)}{1-e\cos
v}$. But $r = \dfrac{a(1-e^2)}{1-e\cos v}$ is the ellipse's polar equation. The Mars
orbit is thus an ellipse, in one focus of which the sun stands.''

\emph{c) The real law.}

How induction, beginning with the observation of the individual, leads from the
individual to the total, in that it leads through the facts to the universal, and
thereby to the cognition of the real law, can be seen from the adduced example.

That the Mars orbit really is an ellipse is not proved by observing all the
individual points in the orbit; such an observation is impossible, since the orbit
contains infinitely many points. Nonetheless induction gives sufficient certainty.
This certainty is not superficially supported on a probability-calculation that
only wishes to approach the elimination of chance: no, here is an essential
inference that grounds a doctrine.

To grasp the probative force of induction, one must well note that it is not to
decide whether the phenomenon in question, here the planet's motion, is subject
% ---- printed p.59 (PDF 67) ----
to a law or not (for that is decided a priori); but rather to establish
\emph{which} law has here in reality found its expression.

Between the a priori presupposition that here is a law, and the empirical question:
which law? the distance is so great that Apelt with reason recommends ``heuristic
maxims'' to lead the observer onto the right track. The ``maxims'' themselves are in
this connection without special interest; whereas the following instructive
direction (Anf. Skr. S. 50) ought not to be passed over.

``In the categorical inference the three termini are each time given; in induction,
on the contrary, the major concept must often first be sought. When one is to seek
something, one must first know \emph{what}, and \emph{how}, and \emph{where} one is
to seek it. Induction seeks the law on which a definite circle of phenomena
depends. One must then evidently know beforehand that the law's necessity, and not
chance (lawlessness), rules in this circle. Thus, e.g., the course of the stars is
bound to a law, the number of the earthly elements, on the contrary, not. But the
mere presupposition of lawfulness is not yet sufficient; one must also have a
certain representation of the character of the law one will seek. When Kepler
sought the figure of the Mars orbit, he knew already prior to his induction, i.e. a
priori, that the planets do not go their way through the heavens according to whim
and chance, but that they describe around the sun an orbit that lets itself be
designated according to the laws of geometry. One must, if one will not miss the way
to truth, know whether the postulated lawfulness is merely geometrical, or
mechanical, or teleological. Thus some seek even today teleological grounds of
explanation for the phenomena of life in physiology, others physical (mechanical
and chemical) ones, because one is not yet agreed about the laws prevailing in this
domain.''
% ---- printed p.60 (PDF 68) ----
That Galileo, who emphasizes the principle of matter's inertia, represents
abstraction, as Kepler induction, and that both directions of mind are united in
Newton (\textit{Philosophiæ naturalis principia mathematica}), is an ingenious
remark of Apelt; whereas the speculative interpretation Hegel (Naturphil. S.
94--124) gives of the fall-motion and the Keplerian laws is just as unstable as his
judgment of Newton's theory of gravitation is unmotivated.

Neither the sensualist nor the idealist is able to know the real law; the former
not, because he denies the a priori; the latter not, because he overlooks the fact.

\begin{center}{\bfseries \S\,11.}\end{center}

That a cognition of the real laws is an inviolable condition for knowing the
existing things, scarcely anyone will draw into doubt. But the question how a
cognition of the real relates to a cognition of the essence indicates in this
connection a new point of view, from which the opposition between Kant and Hegel
can be more closely assessed. According to Kant it would then be possible to carry
it even very far in cognition of the real, without thereby in the least coming
closer to the essence of things; according to Hegel it would conversely be possible
to know the eternal essence, without troubling much about the difference of the real
things. How this is to be understood will first become clear to us as we more
closely investigate the relation between things and laws, things and their
properties, essential knowledge and real objects.

\emph{a) Things and laws.}

Even when one grants what is universally valid in the laws, it nevertheless makes an
essential difference whether one grasps these laws as essenceless rules without
exceptions, or sees in their fulfillment an expression of the nature of things.
% ---- printed p.61 (PDF 69) ----
In the theory of experience a double language is spoken. On the one side it is
said: ``We do not know what matter is in itself; force is to us an unknown, we know
only its expressions; the essence of light is unknown to us, we know only the
phenomenon; we do not know the cause of heat, we know only its effects.'' On the
other side it is said again: ``it is not by speculations and theories, but by
observations and experiments, that we penetrate into the essence of things, and get
to know the hidden forces. We know nature's laws; these laws constitute the
universal in the manifold, the eternal in the changes; these natural laws are true
thoughts of nature, so that what we know in nature is nothing foreign, but the same
as that which lives in ourselves.''

These two views stand in open conflict with each other. If it is true that we know
only the phenomenon, not the essence, only the expression, not the force, only the
effect, not the cause; then it is untrue that the laws of things express the essence
of things. If it is true that the natural laws, as revelations of an eternally
active reason, express the essence of things, then the impossibility of knowing the
essence, the force, the cause, is only a fiction.

Such questions as: what is space in and for itself? what is matter in and for
itself? what is force in and for itself? can betray much unclarity in the a priori.
Matter in and for itself, pure matter, is just as surely an abstraction as pure
space. One cannot ask \emph{whence} or \emph{how} matter has come into space, for
the being of matter and of space is in principle the same. Here is an essential
identity in a real distinction; as the sensible content, matter forms a real
opposition to the outer form. It is so far from the case that the intuition of an
empty space should be merely subjective, that it on the contrary shows space's
objective validity as a real distance between two bodies removed from each other.
% ---- printed p.62 (PDF 70) ----
As little as matter can be in itself without being in space, just as little can
force be in itself without being in matter. Name whatever force one will, it will
always be determined by a peculiar relation in which the one matter can find itself
to the other. What is force of cohesion? An expression for the resistance a body
makes against the separation of its parts. Here, then, there is in the body not
first a multitude of parts in and for themselves, and besides a force in and for
itself; no, the force of cohesion is an expression of essence of the material
parts; the strong cohesion expresses that the parts are something in themselves
only in this, that they are something for one another mutually. As it goes with the
force of cohesion, so it goes also with the force of gravity, with the magnetic,
electrical, chemical forces, and so on.

If all force is thought away, matter is only an empty abstraction; if matter is
thought away, force is an abstraction. If one cannot say that it is matter that
originally gives force existence, since matter, in order itself to be, must
presuppose force; so one cannot, on the other side, assume either that force gives
matter existence, since every expression of force must presuppose an existing
matter from which it can proceed. One need not brood over the question how force
originally came into matter, since the being of force and of matter is in principle
the same.

As for chemistry's so-called elements, it is just as contradictory to wish to
determine their nature in and for itself as to wish to speculate out what the
chemical attraction is solely in and for itself. Oxygen, hydrogen, sulphur, such
elements are only in and for themselves what they express in their mutual relation
and combinations. It is not a defect in our cognition that we cannot see what the
chemical force is unconditionally in and for itself; for it is, as an unconditioned
in and for itself, a nothing, an abstraction, a figment.
% ---- printed p.63 (PDF 71) ----
That Kant in his critique of reason had to strike upon an inexplicable ``thing in
itself'' is comprehensible; but to wish, after all that the Hegelian Dialectic has
clarified, still to hold fast to such a ``thing in itself,'' and thereby to suppose
one avoids the dialectical contradictions, is only to deceive oneself: a thing in
itself, a thing behind all conditions, is a lawless thing; for things are under the
law only insofar as they are with inner necessity for one another mutually, and
condition one another reciprocally; a thing in itself, an unconditioned thing, is
not merely a contradiction, but, when it is held fast, a hardening of abstraction,
an absurdity.

If it is now evident that things lay bare their essence only in their mutual
reciprocity, and if the law is the formal expression of necessity that for this
reciprocity indicates the definite consequences of definite presuppositions; then
it is doubtless also incontrovertible that the law is more than a universal rule
without essence, that a real cognition of law is at the same time a cognition of
essence, that the different laws indicate different sides of the essence of things.

\emph{b) Things and their properties.}

As the essence is known by its phenomenon, so the thing must be known by its
properties. The thing is the existing unity of the properties, and the properties
are in turn the inner determinations of things in existing plurality. But, although
the property is what is proper to the thing, the property of the one thing can in
reality become recognizable only by its finite relation to the property of the
other thing; for the illumination of the matter it is required that the reciprocity
of the properties be seen from the real side.

The difference that in reality is between the immediate and that mediated through
many intermediate links, between the partial, which is apprehended singly, and the
total, whose inner connection demands all-sided, accurate, persevering research,
brings it about that the distance between a preparatory
% ---- printed p.64 (PDF 72) ----
cognition of the phenomena and a concluding cognition of essence steps clearly
forth.

By merely holding to the optical phenomena, one still gets no representation of the
chemical effects of light; from the mere sensation of heat one cannot explain to
oneself that heat expands bodies. How does it now come about that light, the same
cause, can have so many and highly different effects? Can the same thing that is
luminous also be warming, or is there perhaps a peculiar heat-substance? Can light
be transformed into heat, and heat into light, and wherein then consists the
transformation? When such questions are set in motion, it is seen that he who
really wills a true cognition of the essence of light and heat can scarcely, like
Hegel, overlook ``the Newtonian theory'' and ``the Eulerian ether.''

What significance the hypotheses disdained by Hegel must have even for the
cognition of essence is seen already from this, that it is a matter of an
investigation of matter's proper discretion (\textit{Fechner: Atomenlehre. Leipzig
1855. S. 18--32}) when it is a matter of an investigation of heat's identity with
light.

That the emanation-doctrine presupposes atoms of finite size goes without saying;
but neither can the undulation-theory be carried through except under a similar
presupposition. In refraction the broken white ray spreads itself into a narrow
fan of colors; this phenomenon can be brought into mathematical agreement with the
vibration-hypothesis only when the ether-parts are assumed to be separated by
finite discretion. The narrow fan of colors rests, namely, on the fact that the
different rays have different refrangibility, and that the distance of the
ether-parts is not vanishing in comparison with the breadth of a light-wave. The
same view — that, namely, the ether-parts are in finite separation — is further
confirmed when the usual light-phenomenon is compared with the polarization-
phenomenon. In a usual light-ray
% ---- printed p.65 (PDF 73) ----
the vibrations have all possible directions, in a polarized light-ray they are
parallel, in both they are transversal. But the continuation of transversal
vibrations is possible only under the presupposition of the ether's discretion.

As now the light-phenomenon can in all respects be interpreted according to the
vibration-hypothesis, so can the heat-activity too. The propagation of heat by
contact is indeed so different from radiation that at first glance it should seem
impossible to explain such diverse phenomena from the same general presupposition;
but Fourier has shown that the phenomena let themselves be reconciled when one
assumes, namely, that the bodies are not continua, but consist of separated small
parts. The conductivity of the bodies then becomes dependent on the vibration-speed
of the conducted heat, thus on the temperature: heat and the heated characterize
each other mutually.

That the heat-rays act most strongly when they fall perpendicularly on the surface
of the body, and on the contrary more weakly when they fall in an oblique
direction, is inexplicable so long as one continues to start from matter's
continuity; if, on the contrary, one presupposes a finite discretion, the
phenomenon is not merely explicable, but even calculable, in that the ray-effect is
weakened according to the law of the sine.

In chemistry and crystallography not only does the atomic theory here pointed out
find peculiar confirmation, but the doctrines of light and heat at the same time
undergo a new testing. If we then look at the matter in its wholeness, the essence
of light can be known only from the totality of the phenomena; but the totality has
an infinite compass: the one totality reaches into the other. To investigate light
as the pure light in and for itself is impossible; the properties of light are known
only by the properties of the things that are affected by light; the same holds of
heat, of magnetism, of electricity, and so forth. These ``potences'' are not merely
what they are, each for itself in the reciprocity of their properties with the
properties of the things in question;
% ---- printed p.66 (PDF 74) ----
but they are at the same time for themselves what they are in their reciprocity with
one another mutually.

\emph{c) Essential knowledge and real objects.}

In relation to the universal essence Hegel is right against Kant; but reality is
right against Hegel. Instead of knowing the real relation between thing and
properties, Hegel has already in ``\textit{Phänomenologie des Geistes}'' laid out the
development of concepts in such a way that in reflection one proceeds from the thing's
lower to objectivity's higher category. But by this logical progression it is so far
from the case that the problem of reality becomes recognizable, that it is rather
veiled. The universal is the objective, but the universal is identical with thought;
the objective is, in its peculiarly determined objectivity, an object, but
objectivity is identical with its concept: the object must be objectified by the
subject. Along this path one never gets beyond the universal identity of thinking and
being.

If, on the contrary, we let the existing, the thing with its properties, set itself
in sensibly determined opposition to the subjective reflection, the problem of
reality comes into force. In this point Locke sees more clearly than Hegel.

Strictly speaking, the ``substratum's'' identity with itself in its many properties
and accidents cannot be asserted for any real object with a priori necessity. What
is this? ``Silver! one sees it at once by the peculiar luster.'' — But the luster can
of course dazzle; another body could possibly have a similar whiteness. By one or two
properties an object cannot be known; we must then, in the case of silver, test a
multitude of properties: malleability, ductility, specific gravity, crystal form,
melting point, attraction to chlorine, to sulphur, to oxygen, and so on. Each property
must be tested sufficiently, and a sufficient number of properties must be tested:
here is a task of reality in its infinity.
% ---- printed p.67 (PDF 75) ----
It can be granted to the Hegelian analysis of concepts that it, against Hume as well
as against Kant and Herbart, clarifies the dialectical identity of the thing and its
properties, of substance and its accidents, of the concept and the object in general,
and thereby a priori secures for the theory of experience an objective foundation:
that the silver-substance is in essential identity with the properties of silver
cannot be doubted without the law of reality being violated, and Empiricism
contradicting itself; but \emph{which} properties are here united, and \emph{how}
they are united in this real substance, must be tested by observation and experiment;
the real concept fetches every moment from Empiricism.

By force of the empirical research's inner original connection with the a priori,
there is then an essential knowledge; but according to the heterogeneity of thinking
and sensation, real cognition is subject to finite conditions, and all our knowledge
bounded. The bound of knowledge is twofold: a \emph{relative} one, insofar as here
there is always something we do not yet know, but which we can nonetheless in time get
to know; an \emph{absolute} one, insofar as a finite separation of the content of
consciousness and of reality is grounded in our organization: here, then, is something
we in this existence never get to know.

\begin{center}\rule{0.35\linewidth}{0.4pt}\end{center}

% ---- printed p.68 (PDF 76): II. Videnskabslære / Theory of Science ----
\begin{center}
  {\large\bfseries II.}\par\vspace{4pt}
  {\Large\bfseries Theory of Science.}\par\vspace{4pt}
  \rule{2.5cm}{0.4pt}
\end{center}
\phantomsection
\addcontentsline{toc}{section}{\texorpdfstring{II.\quad Theory of Science}{II. Theory of Science}}
\markboth{Theory of Science}{Theory of Science}

\begin{center}{\bfseries \S\,12.}\end{center}

Where the theory of knowledge ends in thought-confusing or merely negative results
(Sophistry, Scepticism), there can naturally be no talk of true science. As the theory
of knowledge is, so must the theory of science become. If the theory of knowledge draws
a critical bound between phenomenon and essence (\textit{Kritik der reinen Vernunft}),
the theory of science will on all its domains remind of the bound, and as in a
\textit{reservatio mentalis} [mental reservation] always tacitly understand the
impossibility of knowing ``the thing in itself.'' If the theory of knowledge culminates
in subjective Idealism (\textit{Wissenschaftslehre}), the scientific doctrinal
structures, insofar as they are the work of the principle, will transform themselves
into mirrors of reflection for ``the pure I.'' If the theory of knowledge raises itself
to the absolute knowledge of objective Idealism (\textit{Phänomenologie des Geistes}),
the sense-things will fall, but the concepts of essence stand as dialectical realities.

On true insight into the conditions of cognition rests the true insight into the
essence of the sciences. Since now the conditions of cognition are not merely inwardly
subjective, but at the same time outwardly objective, the essence of science, in its
different breakings-through and peculiar appearances, is besides at the same time so
closely grown together with the conditions of given ages, with the aim of life, with
the events of history, that in order to understand the theory one must look at the age.

What philosophy was for the aristocratic antiquity, what theology was for the
hierarchical Middle Ages, that is physics for the bourgeois present: cultural history
and the history of the sciences illuminate each other.

How movements in the outer relate to movements in the inner is shown by the discoveries
in the 15th, and the Reformation in the 16th century. To assess the intellectual
advances one must, however, distinguish be-
% ---- printed p.69 (PDF 77) ----
tween the impulses and awakenings by which great leading thoughts are called forth, and
the quiet, persevering activity of spirit that tests the points of view and gives the
content scientific elaboration.

Though science may at times slumber, because minds are dulled in the handed-down, and
the age lacks inspiring elements; so on the other side thought can also be overstrained
and the forces go astray, when the ferment is too strong and the impressions too
excessive (Cardano, Telesio, Campanella, Bruno).

As for the scientific interest itself, the enthusiasm can either exclusively hold fast to
theoretical ideals (Copernicus, Le Verrier), or keep closer to practical results (J. F.
Böttger, Daguerre); both directions proceed from the spirit of science, for the
theoretical will prove its reality in the practical.

But although science lives a vigorous life only when it finds itself in sound and true
understanding with the powers of life, it cannot be drawn down into the series of means,
and subordinated to the self-interested regards of finitude. Only that science is
cultivated which is cultivated for its own sake. To will the scientific merely for the
sake of utility is precisely not to will the scientific. Self-interest deceives itself;
it wills only the means; but the means fails to appear where spirit is lacking: the mere
craftsman's eye does not see so deeply that it sees the new possibilities.

Where, on the contrary, scientific insight is in every direction recognized as a
spiritual power, and, without regard to self-interest's petty Wherefore?, is embraced
with the interest of infinity, of ideality, science will not remain standing as a pale
theoretical shadow, but, like a luxuriant tree, bear good fruits beneficial to life. The
discoveries of natural science in the present century (Galvani and Volta; Ørsted and
Humboldt) furnish a series of grand examples of true reciprocity of theoretical and
practical interests.

\begin{center}\rule{0.35\linewidth}{0.4pt}\end{center}

% ---- printed p.70 (PDF 78): A. Videnskabernes Videnskab / The Science of the Sciences ----
\begin{center}
  {\large\bfseries A.}\par\vspace{2pt}
  {\bfseries The Science of the Sciences.}\par\vspace{4pt}
  \rule{2cm}{0.4pt}
\end{center}
\phantomsection
\addcontentsline{toc}{subsection}{\texorpdfstring{A.\quad The Science of the Sciences}{A. The Science of the Sciences}}

\begin{center}{\bfseries \S\,13.}\end{center}

Since science has its root in human life, where the ideals of fantasy go before those
of thought, it is natural that the different sciences come into the world at first in
the swaddling-clothes of imagination, and are burdened with unscientific admixtures.
Thus astronomy was long grown together with astrology, chemistry with alchemy,
philosophy with mythology.

It doubtless lies in the drive of science to lift the pupation, break its husk, and
purge out the foreign constituents. But cultivation costs time, the purging exertion;
the hindrances that the human spirit has here to overcome are of different kinds.

a) An \emph{outer hindrance} proceeds from the natural strife between egoism and truth.
The intellectual love of truth has at all times waged combat, not only against the
prejudices of the ages and the thoughtlessness of the multitude, but above all against
their egoism, whose power and influence saw itself threatened by the rising light. So
long as a given condition is of such a character that the spiritual guardians of the
population have an interest in holding to the untrue, they will not easily convince
themselves of the true. Science is doubtless objective; but the subjective does come
into the conviction: insight is conditioned by prospect. Were there still in the
countries considerable positions for astrologers, would there then not also still be
those who felt themselves called to study astrology — of course not for the sake of
prestige, but because they had really convinced themselves of the deeper speculative in
astrology? Those who will the truth must be zealous for the truth; science too has its
martyrs (Vanini, Galileo, J. G. Fichte).

b) Of a nobler, though not less dangerous, nature are the \emph{inner hindrances} that
the richly gifted, overreaching
% ---- printed p.71 (PDF 79) ----
subjectivity can prepare for itself. In enthusiastic striving after deeper insight it
easily happens that the thinker confounds the subjective representation with the
objective object, and then supposes he finds that in reality which nevertheless is
found only in his own fantasy. In the greatest minds, the noblest natures — a Plato
(\textit{Timæus}), a Kepler (\textit{Mysterium cosmographicum}), a Schelling, a Hegel —
mixtures can be pointed out, whereby partly in themselves unscientific elements are
worked into scientific form, partly in themselves scientific constituents are found
placed in an unscientific, fantastic connection.

c) To this comes still the tension between science and view of life, which is increased
by an ever growing \emph{conflict between the real and the ideal}.

The more science in the finite sense holds fast to the real, the more easily will it
undeniably avoid the fantastic, but the more difficult will it on the other side be to
satisfy the ideal demands of personality. The stronger the scientific upswing is toward
the ideal, the more easily will science doubtless satisfy the demand of spirit, but the
more easily will it then also expose itself to overrunning the conditions of reality,
and be caught in illusions.

What then is science? By what does one everywhere and under all shapes know the
scientific from the unscientific? How ought the personal view of life to relate to the
objective, impersonal results of science? For the answering of these questions is
required a comprehensive theory of science, a science of the sciences.

\begin{center}{\bfseries \S\,14.}\end{center}

From the relation between the theory of knowledge and the theory of science it follows
that the true science of the sciences can be none other than philosophy. With this it
agrees exactly that the true comprehensive theory of science can come about only in
that the fundamental interests of the diverse sciences are represented in the
deliberations from which
% ---- printed p.72 (PDF 80) ----
it proceeds. Under ``the comprehensive discussion and debate of all against all''
(Sibbern) each science then has its vote, and each its contribution. That philosophy
works up, clarifies, and organizes these contributions into an intellectual whole
proves that it in this respect does not decide the matter on its own, but works as the
sciences' common organ, as their own universal consciousness, the essence that is their
kinship, the light in which they know and cognize one another mutually.

It is therefore not presumptuous arbitrariness when philosophy sets the different
doctrinal structures before the tribunal of the theory of science, and demands that
with respect to each one of them in particular an account be given both of the
foundation of the structure and of the manner of building; it is not self-taken
authority when, after carefully held examination, it permits itself to assign to those
theories which in foundation and manner of building have more likeness to aesthetic
castles in the air than to scientific doctrinal structures ``a peculiar place'' outside
science.

Philosophy, with all this, does only its duty, in that it, bound by the laws of
cognition, does not act according to whim, but as the lawful organ of the sciences, the
submissive handmaid of scientificity.

If we set together the great representatives of the development of spirit — philosophy,
theology, and physics — with one another mutually; then not only can the first, insofar
as it is set in opposition to the two others, by no means wish to exempt itself from
scientific justification, but reasonably neither can the latter.

Were physics, dazzled by its present popularity, to cast contempt on the investigating
theory of science, and, instead of heeding the light of thought that will shine through
its ground and essence, only to rely on the immediately given, and invoke the
incontrovertible authority of the senses, it would only thereby set itself lower than
the medieval
% ---- printed p.73 (PDF 81) ----
theology. Theology too can speak of the given, use the popular, and invoke an
incontrovertible authority; and so long as there is in human nature a spark of a Higher,
the authoritative pronouncement of faith will always be able to outweigh the
authoritative pronouncement of sensuousness.

Where, then, thus the authoritative pronouncements of Supranaturalism and Naturalism —
now alternately, now partially, now in open enmity, now in mystical agreement — step into
the place of the theory of science and thereby give themselves scientific importance, it
is doubtless natural that philosophy is either straightforwardly despised, or at least
regarded as an intellectual plaything; but then scientificity too is in danger.

But does theology then perhaps want no theory of science? If one could at once hold to
Anselm and yet profit from Descartes, could get \textit{credo ut intelligam} [I believe
in order that I may understand] to say \textit{de omnibus dubitandum est} [everything is
to be doubted], and conversely, \textit{de omnibus dubitandum est} to proclaim itself as
the right \textit{credo, ut intelligam}, then theology would surely readily unite with
philosophy in commending a thorough and sound theory of science; but on these conditions
it is an impossibility for philosophy to unite with theology, since the matter itself is
a contradiction.

Although philology, criticism, and history — which in and for themselves belong under
science — constitute constituents of theology, they by no means constitute the
theological in theology. Insofar as the named disciplines are treated scientifically,
they are treated untheologically; insofar as, on the contrary, they are treated
theologically, they are treated unscientifically. The theological in theology is known,
namely, by this, that the cause from which all effects in the realm of nature as in that
of spirit are partly immediately, partly mediately derived and explained, is not sought
in the nature of things or the essence of reason, but in the absolute will, in the
Almighty, who creates out of nothing. But even the purest, most acute, and most
consistent application of such a
% ---- printed p.74 (PDF 82) ----
principle (Occasionalism; Malebranche) makes science unscientific.

When in the theory of science it is a matter of a principle, the matter cannot be
decided by a More or Less, as though one, by following ``the middle way,'' would suppose
one could still get through scientifically by assuming a little of the supernatural, if
only one did not make too much of it.

The question of a principle is a question of an unconditioned, all-conditioning
presupposition; such a question must therefore always be answered with Yes or No. Either
theology must once and for all acknowledge the supernatural causality, and has then at
the same time repelled the objectively rational foundation, or else once and for all
acknowledge this foundation, and has then at the same time renounced every support in
the supernatural. In the first case theology will thus cease to be science, in the last
case to be theology.

\begin{center}{\bfseries \S\,15.}\end{center}

That life and science cannot relate outwardly and indifferently to each other follows
immediately from both's intimate connection with culture and humanity; that life and
science are nonetheless independent rests on a principled separation of the subjective
and the objective.

In science subjective moods — inclination, disinclination, fear, hope, and the like —
have nothing to signify. Insofar as the question here is of the real, we will not know
how this real, according to our opinion, ought preferably to be, but how it is in itself,
and necessarily must be; insofar as the question here is of the possible, the question
here is not how far this possible best satisfies our hopes, but how far it agrees with the
conditions of existence.

Even when science works in the service of the interests of life, it works by force of the
objective. It is not subjective feelings, not pious wishes, but insight into the
% ---- printed p.75 (PDF 83) ----
essence of things, that leads off the lightning-flash, and sets the steam-carriage in
motion.

Science is doubtless impossible without subjectivity; but the subjectivity is serving.
Here one senses, here one feels, here one thinks; but the merely subjective in sensation
and thought must everywhere be purged out as the contingent, the disturbing, the untrue;
every presupposition is objective. Even there where the content touches subjectivity as
immediately as in the aesthetic, science sees only the true, insofar as it sees in the
matter the incontrovertible, that which is independent of arbitrary predilection, the
universally valid.

From the standpoint of life, on the contrary, subjectivity rules; here all the objective
is serving. In infinite reciprocity with the surrounding world personality works, even in
devotion and sacrifice, for subjective interests, subjective aims, subjective ideals.
Everywhere where subjectivity here is in relation to another, it is through this other in
a peculiar relation to itself. On selfhood, individuality, character, personality life
lays infinite emphasis. That subjectivity now, among several possible tasks, also chooses
science by no means suspends the opposition of life and science. One thing it is to live
for a science, another to lay science at the foundation of life.

Insofar as subjectivity, in rising conflict with the objective world, develops the
personal forces, we can here distinguish the following chief stages.

a) \emph{The relative subjectivity.} At this stage the objective is set equally as means
and as end. By setting it as end, subjectivity knows its bound and relates in serving
fashion; by setting it as means, it satisfies the self, and asserts freedom. Out of
interest of life for the self and the work, for the worker and the labor, subjectivity
will then here make manifold agreements with the objective, venture into many a strife,
suffer many defeats and win many victories, in that it runs
% ---- printed p.76 (PDF 84) ----
through the scale of development from the useful to the exalted, from bourgeois industry
to heroic exploit.

b) \emph{The infinite subjectivity.} By breaking through the objective barrier and
overstepping the bound of the finite, the consciousness of man strives to swing itself up
to infinite freedom. Since reality cannot be denied, and a material surrounding world not
annihilated, the self-liberation must be ideal. In vigorous one-sidedness a choice is
made between two mutually opposed expedients. Partly the objective is under all shapes so
reduced to a mere means that the enjoying individuality can in each moment make itself the
lord of things (Hedonism; Aristippus); partly the objective, as the universally rational,
is so breathed into consciousness that all finitude is transformed into something
indifferent; purified and hardened in infinite resignation, the acting subjectivity will
determine itself, in that it is determined by fate (Stoicism; Zeno — Epictetus).

An inner negative infinity, a dialectical unity of ethical earnestness and carried-through
irony, is the deeper presupposition and unattained ideal of these one-sidednesses.
(Socrates).

c) \emph{The absolute subjectivity.} The formal self-liberation ends in the negative;
subjectivity seeks life, but finds death. Only that life which overcomes death is the
true, the eternal life. That the individual's innate egoism is its own original guilt, a
sin that deserves death; that the pure, guiltless subjectivity has voluntarily borne the
sin, and saved the guilty from death; that the natural being is sheer impotence, only the
will of spirit unconditioned power; that the present is only the means, the future the
end: such presuppositions demand faith, but exclude knowledge.

In open conflict with all that scientific insight can in any way acknowledge as objective,
faith is so far from obtaining any support from science that it precisely by collision
with the objective
% ---- printed p.77 (PDF 85) ----
knowledge must be thrust back into itself, and confess that its only ground is a need of
subjectivity, an inner incontestable demand. At this stage all the objective is
transformed into a point of transition for the subjective. Here it is not a matter of
life's relation to the lifeless, but of the living's relation to the absolutely living;
here the question is not of the will's relation to the will-less, but of the willing's
relation to the absolutely willing.

Life is subjective, the eternal, true life absolutely subjective; the will is subjective,
the eternal, unconditioned will absolutely subjective; the objective is thus, as objective,
not the True: ``Subjectivity is the truth'' (S. Kierkegaard).

\begin{center}{\bfseries \S\,16.}\end{center}

In that the strife between faith and knowledge is thus driven to its extreme, the
intellectual reconciliation is at the same time introduced: the absolutely heterogeneous
principles cannot, in the essential sense, do each other harm. The faith that hates
knowledge, because it must strengthen itself by ignorance, is a false faith. Reality will
be reckoned with; it does not help to close one's eyes: if the earth really goes round,
then it goes round just as well with the believing as with the non-believing. In that
faith declares that subjectivity is the truth, it doubtless loses all significance for
science; but, in that science declares that objectivity is the truth, it at the same time
renounces all influence on faith.

Faith and knowledge can then without contradiction be united in one consciousness; their
mutual strife does not proceed from the principled. The strife is psychological; it is
human frailty, which has neither resignation enough to reconcile itself with science's
objective, nor inwardness enough to hold fast to faith.

That faith hates science is untrue; that it makes subjectivity independent of science is,
on the contrary, true. To set faith highest and yet keep science is possible; to set
science highest, and yet preserve faith, is impossible.
% ---- printed p.78 (PDF 86) ----
But although science itself must see that it can neither acknowledge the absolute
subjectivity nor prescribe laws to it, it is nevertheless in its interest to have this
subjectivity illuminated as all-sidedly as possible.

The separation of scientific and unscientific is made difficult by this, that what
according to its nature must be rejected by science nonetheless becomes an object of
scientific treatment, and so in a way does go along in science. Madness, e.g., is
incompatible with science; nonetheless the symptoms of madness and the treatment of the
mad become an important object for psychology and medical science. As a fetus of
superstitious fantasy the myth is in conflict with the scientific; nonetheless there is a
science of myth, a scientific mythology.

If the most diverse can be set under scientific illumination; if even the most subjective
can become objective, insofar as it is objectified in reflection; then from this point of
view nothing can be objected either against a science of subjectivity forming itself, a
science of the subjective.

That such a science of subjectivity (psychology, aesthetics, ethics, philosophy of
religion) must not, however, be confounded with what one otherwise is wont to call a
science of faith lets itself be seen when we clarify to ourselves the difference between
the scientific in the wider and in the narrower sense.

Scientific in the wider sense is everything that becomes an object of science's treatment:
the scientific lies in the treatment; scientific in the narrower sense is, on the contrary,
only that whose validity science sees, whose truth it can assess, test, and acknowledge.

By using the ambiguity that lies in the category ``scientific,'' the doctrine of faith
sometimes succeeds in giving itself the semblance of science; it is then the task of the
theory of science to point out the most important sophisms belonging here.

\emph{First sophism:} ``If faith is not to rest on subjective notions, and be given over
to the whims of individuals,
% ---- printed p.79 (PDF 87) ----
it must be possible expressly to point out what one is to believe, and in whom one is to
believe: faith must have its object. To investigate the objects of faith thoroughly is to
investigate them scientifically. The doctrine of faith rests on such an investigation;
thus the doctrine of faith is a science.''

This sophism gives the matter the appearance that the objects of faith are secured by
science, because they are tested in science. The matter is the reverse. The objects of
faith dissolve into vapor as soon as science is first permitted to apply its criticism. It
is just as impossible for scientific criticism to acknowledge the Gospel of the Risen One
as it is impossible for science to acknowledge the myth of Hercules. To acknowledge a
miracle by force of scientific criticism is a contradiction, since the possibility of a
miracle makes scientific criticism impossible: where the wonder speaks, the understanding
must be silent; where the miracle begins, science must end.

\emph{Second sophism:} ``Every peculiar science must necessarily have its presuppositions.
Just as logic presupposes and justifies the validity of thought, physics the validity of
experience, and aesthetics the validity of the beautiful, so the doctrine of faith must for
its part presuppose and justify the validity of faith and the content of faith. He who
proves that the doctrine of faith is not a science because it has its peculiar
presuppositions proves at the same time that there is no science at all; but here it holds:
\textit{qui nimium probat, nihil probat} [he who proves too much proves nothing]. Thus the
doctrine of faith is a science.''

This sophism dazzles by the ambiguity that lies in bringing along a presupposition. What
makes the doctrine of faith unscientific is by no means the circumstance that it brings
along presuppositions, but precisely the circumstance that the presuppositions it brings
along conflict with science. It presupposes the miracle, it presupposes the absolute
validity of subjectivity and of subjective conviction; with these presuppositions it has,
so to speak, forced
% ---- printed p.80 (PDF 88) ----
its way in upon science, and broken the scientific measure. If the consequences are to be
scientific, the presuppositions must also be so; on unscientific ground no scientific
structure can be raised: natural science is possible, but miracle-science is impossible.

\emph{Third sophism:} ``The doctrine of faith's supernatural presuppositions have, each
for itself, doubtless only validity for the believer; but between these propositions
mutually there is a logical and to that extent objective connection. Since the connection
is objective, and the propositions of faith precisely find their true interpretation and
deeper grounding in this connection, the grounding itself is objective, and the doctrine of
faith thus, by force of this grounding, a science.''

The dazzling in this sophism rests on this, that a logical connection without more becomes
an objective connection, and this in turn an objective grounding. ``All men are triangles;
Caius is a man, thus Caius is a triangle. In every triangle the sum of the angles is equal
to two right ones; Caius is a triangle, thus the sum of the angles in Caius is equal to two
right ones.'' Between these propositions there is here undeniably a logical connection, but
is there therefore here an objective scientific grounding? From the standpoint of life
there can be an eternal difference between the Holy and the Profane, the Evangelical and the
Mythical; but, when science once cannot assess this Holy, cannot master this Evangelical,
the difference is in the eyes of science as though it were not at all.

Since the propositions of faith are each for itself unconditionally subjective, they are in
science each for itself just as arbitrary as the most fanciful fantasies. If there is then a
scientific grounding of doctrines of faith, then there is also a scientific grounding of
doctrines of myth, a scientific grounding of nonsense.

\begin{center}\rule{0.35\linewidth}{0.4pt}\end{center}

% ---- printed p.81 (PDF 89): B. De særskilte Videnskaber / The Particular Sciences ----
\begin{center}
  {\large\bfseries B.}\par\vspace{2pt}
  {\bfseries The Particular Sciences.}\par\vspace{4pt}
  \rule{2cm}{0.4pt}
\end{center}
\phantomsection
\addcontentsline{toc}{subsection}{\texorpdfstring{B.\quad The Particular Sciences}{B. The Particular Sciences}}

\begin{center}{\bfseries \S\,17.}\end{center}

That philosophy calls itself the science of the sciences must not be misunderstood as
though it hereby would give itself out as an oracle of the sciences, the only, proper, and
true science. How such a misunderstanding arises, what confusions it brings with it, ought
to be heeded.

Every knowledge is determined by its object, the peculiar material that cognition works up.
A comprehensive apparatus of knowledge, a methodically acquired wealth of knowledge, gives
learning, but learning is not yet science; only with the development of principles and the
rational working-through of the material comes the scientific. But, although to a certain
bound there can just as well be science without learning as learning without science, every
penetrating cognition of the connection of real facts will always require a peculiar union
of learning and science. All so-called positive sciences are learned sciences.

The special is always subordinate to the universal; but in order to enrich our insight into
the universal, we must begin with the special. The more a chief science expands, the more
subsidiary subjects it will secrete off. It goes with science as with Alexander's conquests;
the greater the empire becomes, the nearer is the turning-point when the ring is changed and
the division takes place. If now such a division, in the case of science, is not to signify
the empire's downfall, but a progress of the constitution, the smaller divisions must be
strictly subordinated to the greater, and these to the organized whole: the totality is the
truth.

If now the totality of truth is grasped as an a priori-metaphysical essence, then it is
understood from this how the
% ---- printed p.82 (PDF 90) ----
philosophy of absolute knowledge has been able to set itself as the one knower upon the
throne and look down on the learned sciences as on its subordinates. Learning is means,
science is philosophy. The particular sciences must now either work as day-laborers, or see
to learn philosophy, and explain Speculation to themselves.

``Speculation is All, is to behold in God, to contemplate that which is. Science itself has
worth only insofar as it is speculative, that is, contemplation of God as he is. But the
time will come when the sciences more and more will cease, and the immediate cognition set
in. Only in the highest science does the mortal eye close, in that no longer man sees, but
the eternal seeing itself is seeing in man.'' (Schelling).

But as much as these and similar proclamations testify to a fusion of religious and
scientific enthusiasm, still such an apocalyptic reference to a turning-point, when ``the
eternal seeing itself'' will become ``seeing in man,'' cannot replace for science the lack
of composure and methodical sobriety that is an inviolable condition for advance in
knowledge.

One-sidedness has, as the theory of knowledge shows, its deeper ground in the misprision of
the element of representation that characterizes the objective Idealism.

Just as light in given cases must presuppose a flame, the flame a combustion, and the
combustion combustible materials; so absolute knowledge must doubtless presuppose diverse
materials of representation, but only in order to let them burn up in the Dialectic from
which the concept shines. If now the material of representation is in this way treated, not
as an intellectual building-material, still less as the flesh and blood of thoughts, but
only as fuel; then it follows of itself that one does not take it so exactly with the
individual pieces of the fact.
% ---- printed p.83 (PDF 91) ----
The consequence of this is naturally mistakes in the individual and arbitrarinesses in the
apprehension of the total; but, even when these blunders are corrected, the standpoint will
be no less false.

To value the element of representation, one should remember that the intellectual light, just
as little as the physical, can furnish colors and visual images unless it is refracted in the
different media, thrown back by the different objects, caught up and clarified in definite
impressions. Even when the presentation of the categories corresponding to a positive content
keeps nearest to the universal, one notices an essential difference between the emptiness of
the barren abstractions bereft of factual knowledge and the play of colors, the natural
freshness, with which concepts spring forth that are worked through by intuition, and testify
to autopsy.

\begin{center}{\bfseries \S\,18.}\end{center}

In his treatise ``\textit{Ueber den Begriff der Wissenschaftslehre}'' (\textit{Sämmtliche
Werke, Berlin 1845}) Fichte has raised the question how far the theory of science can be
assured that it on the whole has exhausted human knowledge. The answering of this question
reduces to the following circle: ``When the proposition $X$ is the first, highest, and
absolute principle of human knowledge, then human knowledge is a single system; for the last
follows from the proposition $X$. Since now there is to be a single system in human knowledge,
the proposition $X$, which really (according to the science set up) grounds a system, is the
principle of human knowledge as such; the system grounded on it is that single system of human
knowledge. Beyond this circle one cannot get.'' Further it is said (S. 65) of the relation
between the theory of science and the particular sciences: ``One sees why only the theory of
science has absolute totality, whereas all particular sciences will be
% ---- printed p.84 (PDF 92) ----
infinite. The theory of science contains merely the necessary; all the other sciences bear on
freedom, both with respect to our own spirit and to the nature independent of us.''

From the standpoint of absolute knowledge Hegel has given a presentation of the different
sciences' mutual connection.

``In a usual encyclopædia the sciences are taken empirically, as they are found. In it they
are to be completely enumerated, and ordered in that the like, that which coincides under
common determinations, is set together. The philosophical encyclopædia is, on the contrary,
the science of the necessary connection determined by the concept, and of the philosophical
coming-to-be of the sciences' fundamental concepts and principles.''

``The sciences are, according to their mode of cognition, either empirical or purely rational.
Absolutely regarded, both are to have the same content. It is the aim of scientific striving
to raise that which is known empirically to the always True, to the concept, to make it
rational, and thereby to incorporate it into the rational science.''

``The rational expansion of the sciences is at the same time an expansion of philosophy
itself. The whole of science then divides into 3 chief parts: 1) \emph{Logic}. 2) \emph{Natural
science}. 3) \emph{Science of spirit}. The sciences of nature and of spirit can then be regarded
as the system of the \emph{real} or particular sciences, in distinction from the pure science
or logic, since they are the pure science's system in the shape of nature and spirit.''

As transition from logic to natural science, mathematics is set up with respect to space and
time. Natural science itself falls into: a) \emph{Mechanics} with respect to matter and motion.
b) \emph{Physics of the inorganic} with respect to light, magnetism, electricity, the physical
bodies, and the chemical process. c) \emph{Physics of the organic}, under which then again
belong geology, which
% ---- printed p.85 (PDF 93) ----
treats the formation of the earth; botany and plant physiology, whose object is vegetable
nature; physiology, \emph{comparative} anatomy, zoology, which treat healthy, animal nature, as
medicine the sick. The science of spirit too falls into three sections: a) Spirit in its concept,
\emph{Anthropology}, \emph{Psychology} (language — \emph{Philology}). b) The practical spirit,
which expresses itself in deed and action; under this \emph{the doctrine of right}, \emph{political
science}, \emph{history}. c) Spirit's completion in art, religion, and science, to which correspond
\emph{Aesthetics}, \emph{doctrine of religion}, \emph{Philosophy}.

But although the particular sciences are thus set up against one another mutually and against
philosophy, their ideal independence is nevertheless only a semblance. The encyclopædia ``is
properly a presentation of philosophy's universal content; for what in the sciences is grounded
on reason depends on philosophy, whereas what is in them of that which rests on arbitrary and
external determinations, or, as it is said, is positive and statutory, lies outside it.''

From his Kantian standpoint Apelt (\textit{die Theorie der Induction}, S. 65--70) has indicated
the principled in the separation of the sciences' diverse fundamental character.

``There are,'' it is said, ``two different sources from which our cognitions spring, sense and
pure reason. From the former comes sense-cognition, from the latter the pure rational cognitions.
The latter fall into mathematics and philosophy. There are as a consequence three different kinds
of cognitions: \emph{empirical}, \emph{mathematical}, \emph{philosophical}. Each of these three
modes of cognition has its peculiar character: contingency and intuitability belong to the
\emph{empirical}, intuitability and necessity to the \emph{mathematical}, necessity without
intuitability to the \emph{philosophical}.''

``In mathematical cognitions we become conscious of the forms of composition of things, in
philosophical ones of their forms of
% ---- printed p.86 (PDF 94) ----
connection \textit{in abstracto} [in the abstract]. Kant called the former: forms of the figurative
synthesis, i.e. intuitable forms of combination, the latter: forms of the intellectual synthesis,
i.e. only thinkable forms of combination, laws.''

``The consideration of the different system-forms' relation to one another and to the different
kinds of cognition leads to a scientific architectonic, i.e. to a doctrine of a system of all
human sciences.''

``The three different modes of cognition correspond to the three different system-forms. The pure
philosophical form of cognition belongs under the categorical system; the pure mathematical form
under the hypothetical system, and the pure empirical content under the conjunctive system.
Philosophy is a cognition of mere concepts. But with such a mode of cognition no other systematic
derivation is possible than by a categorical subordination of peculiar cases under the generality
of the rule. Mathematics is a cognition by concepts of pure intuition. In this intuition the least
composite is the most original and the ground from which new consequences are derived by ever new
compositions. The derivation of the composite from the uncomposite happens here by hypothetical
connection of thought. Finally the pure empirical has only to do with the apprehension of the
individual given by sense-intuition. Each fact rests on itself, and the one stands beside the other
in a conjunctive system that co-orders parts into a whole.''

The critically apt in this deduction is to be acknowledged; but it shows itself here, as everywhere,
that the merely critical, without true peculiarity with the dialectical, is in philosophy unfruitful,
and leads to schematism. The separation of sense and pure reason, on which this whole architectonic
is grounded, has in reality no significance. The separation rests on an abstraction, but the
abstracted will again be integrated by that from which it is abstracted.
% ---- printed p.87 (PDF 95) ----
That sense-intuition is heterogeneous with thought is incontrovertible; but these heterogeneous are
for reflection dialectical moments in the same real cognition. The sensible and the thinkable, the
individual and the universal are in the thing not the same and yet the same. In reality all is
concrete. The cognition of reality is conditioned by concretions; but the intellectual concretions
are in turn conditioned by the integration of the abstract forms.

No one who has even merely a general concept of differential and integral calculus will, when he
reflects philosophically, be able to deny the dialectical relation between the thought line, the
surface's negative bound, the determinate negation that cannot be integrated, and the intuited line,
the rectangle's differential, the positive bound that can be integrated into a finite rectangle (Cf.
\emph{Philosophie og Mathem.} S. 65--67; 75--76).

The immanent Dialectic will dissolve the particular sciences' independence, and transform them into
moments in philosophy's system; the delimiting criticism will restrict philosophy to a peculiar
rubric, and supposes on the whole to help by schemata and rubrics. Both efforts are one-sided.
Though a science for itself, philosophy cannot by any critical line of demarcation be separated off
from the other sciences; and although the other sciences have each in itself the universally
scientific, they nevertheless do not have it in such a way that ``what in the sciences is grounded on
reason'' should therefore be able to be said to be ``dependent on philosophy.''

Mathematics's difference from philosophy is not determined straightforwardly by the intuition and
non-intuition of quantity, but by the peculiar analysis to which mathematics subjects the quantitative.
(Phil. og Mathem. S. 9--14).

As for the positive sciences, their independence is immediately given with the material peculiar to
each in particular, without its being thereby given that philoso-
% ---- printed p.88 (PDF 96) ----
phy should not also in its way be able to appropriate the same material. The harder the opposition is
between the object and reflection, the more remote will the positive science lie for ``the absolute
knowledge's'' philosophy. It then makes here an essential difference whether the given object of
thought is aesthetic, historical, or physical.

The aesthetic object, the given work of art, is doubtless always from one side or another related to
something sensible, and must to that extent be apprehended empirically. But the empirical is here so
shone-through by consciousness that the spiritual rules. The objective Idealism has then also without
more vindicated for itself aesthetics as part of philosophy. Meanwhile aesthetics, with its special
analyses, its autopsy, its history of art, and so on, has another side, from which it, without giving
up its scientific character, may well be said to occupy an independent position over against philosophy.

That history, with its facts and its apparatus of learning, stands still more independent, is evident;
nonetheless the historical object, as belonging to the past, the life of spirit, and memory, still has
so ideal a tinge that Idealism itself can move within these domains as though it were only on its own.
Only physics is hard, stands unconditionally against, and permits no idealistic encroachments. Against
the physical object, against reality's polygon, ``the absolute knowledge'' strikes, and betrays its
weakness.

\begin{center}{\bfseries \S\,19.}\end{center}

\emph{History narrates}, \emph{philosophy reflects}: here is a bound that cannot be leveled by a More
or Less. The narration can doubtless take up reflection, and conversely; but the former is nevertheless
a mark of history, the latter of philosophy.

To narrate is to determine the universal by the individual in such a way that thought is dramatized in
events, and the inner goes up in the outer: here is the element of representa-
% ---- printed p.89 (PDF 97) ----
tion predominant. To reflect is to determine the universal in such a way that events furnish concepts,
that intuition is deepened into the inner, and representation dissolves into thought.

From this difference in the manner of presentation a difference in the treatment of the whole and its
parts is a natural consequence. The narrating history moves through the particulars of the parts step by
step forward to a mediate making-visible of the whole; the reflecting philosophy begins at once with
insight into the total, apprehends the particular parts as essential moments in the content of the whole,
and cannot give itself over to the detail. Even when philosophy interests itself most of all in the
individual, this individual is always shone-through by the total. Philosophy can from different sides
illuminate the total's self-mediation through the special; but it does not tolerate the special's
immediacy, and cannot follow the narrating history along mediacy's long detour. Philosophy has, as the
philosophy of history, another aim than learned history; but for the same reason learned history then also
has another aim than philosophy.

If here now, with respect to the aim and the means, the task and its treatment, there is a principled
difference between the historical and the philosophical, then history and philosophy can neither advance
in converging lines, so that they approach a fusion, according as they approach a completion. Presentations
in which the historical and the philosophical equally counterbalance, and then also in a way halve each
other, are only intermediate forms, neutral domains, which with their temporary fruitfulness do not concern
the principles.

Learned history and the philosophy of history can mediately further each other, correct and enrich each
other, and thereby come into an ever more intimate understanding; but precisely this understanding will
confirm the difference.
% ---- printed p.90 (PDF 98) ----
To separate true from false, history must be critical; to point out the causal connection of events,
pragmatic; to find and present the spirit whose interests and plan the events serve, philosophical. In all
three respects the learned historian may well need philosophical cultivation; but his activity of spirit
in no respect goes under philosophy. Everywhere referred to facts and grown together with the given, the
historian's thought, far from losing itself in the philosophically universal, will on the contrary be
articulated by the intuitable, and completed in artistic individuation: historiography is an art.

If a historian has made some chief part the object of an all-sidedly penetrating study of sources, he has
thereby come so near to a past reality that he as it were sees it face to face; will then not his fantasy
be so filled with images, and his mind so crisscrossed by individual impressions, that he can scarcely,
without preceding ferment, order his material and begin his presentation of what has been intuited?

Essentially seen, the tension rests on a peculiar reciprocal relation between the real and its depicted.
In the depiction the real is to be rendered; with all separation of essential and inessential there can
here not, as otherwise in an art, be thought of aesthetic idealization. It is a real personality, a real
state, a real age, whose physiognomy is to be seen in given features, unadorned as it is, with wrinkles and
freckles. And yet spirit is to be present, and imprint itself in the depiction, that it not be overloaded
with the trivial, and run out into the real's endless prolixity.

With such a task in view the historian will see that there is no peculiar form that completely exhausts
reality's content, but that in the presentation there is incessant strife between loss and gain. What here
in one respect is gain, is in another respect
% ---- printed p.91 (PDF 99) ----
loss; this holds not only of the plan on the whole, but repeats itself even in the execution of the
individual. That the party spirit does not always need to do violence to the facts in order to get history
on its side, it owes doubtless nearest to its own sophistical suppleness, but yet also essentially to the
circumstance that the objective in history, however strictly it is held fast, necessarily obtains a
coloring from the objectifying consciousness. Different presentations of the same subject can, according
to the individuation, be equally justified.

But if the historical form is so artistically individual that one and the same universal task (e.g., the
history of the French Revolution) first finds its complete solution when it is solved in several peculiar
ways; then it is seen from this to what degree the element of representation must protest against a fusion
of learned history and the philosophy of history.

\begin{center}{\bfseries \S\,20.}\end{center}

It is not the whim of chance, and not human arbitrariness, that may draw the bound between the particular
sciences. Each science is in itself an essential totality; its bounds, divisions, and elaboration are
conditioned by the material, governed by thought, and determined by the comprehensive task.

In the method of every science the universally scientific is seen: everywhere a reciprocal determination of
the immediate and the mediate in cognition, of detour and shortcut, of necessary movement of the matter and
free movement of thought (Philosophie og Mathematik. S. 44--46); likewise every science's systematizing
striving will require definitions, divisions, subordination of lower and higher unities under one another
mutually and under a highest unity, and so on. But the peculiar material and the peculiar aim in which it is
worked up will nonetheless in the method as in the formation of system
% ---- printed p.92 (PDF 100) ----
bring with it such decisive peculiarities that each science thereby not only becomes a doctrinal structure
in a peculiar style, but is even in an inner, infinite way organized into a relative monad, an intellectual
individuality.

To what degree the peculiar material acts upon the formation of system will especially become manifest from
those sciences whose content is expressed in typical natural forms. Here is, in the mineral kingdom, the dead
form: the crystal; in the plant kingdom the living form: the organism; in the animal kingdom the ensouled
form: the individual.

a) \emph{The dead form.} In its typical purity the substantial form of bounding is mathematical. By laying a
mathematical body (e.g., the cube) at the foundation for the determination of the axes and the corresponding
section, pure crystallography can with the sureness of an exact science round itself off into a whole
articulated into particular systems — the spheroidal, the rhombohedral, the pyramidal, the prismatic. Hereby
it becomes not only possible to lead an infinite manifold of derived special forms back to their fundamental
form valid in each system; but even the combinations, alterations, and distortions that proceed from formations
in free nature are so far from bringing the whole into confusion that they on the contrary become intelligible
by the systematics of the pure types (Applied crystallography).

It is otherwise when the question is not merely of the substantial forms of bounding, but at the same time of
the corresponding substances' chemical and physical properties, of the mineral's characterization. Here the
preponderant regard to the chemical substantiality (Berzelius) comes into open conflict with an equally
preponderant regard to the individuality expressed in the physical properties (Mohs), the strict demand for
systematic grounding in conflict with the schema and its many rubrics.

A more equal regard to chemical and physical distinguishing marks can now doubtless lead to ingenious
combinations of
% ---- printed p.93 (PDF 101) ----
system and schema; but only gradually, as science penetrates into the essential connection of the substances'
chemical and physical properties, will the schematic go up into the systematic.

b) \emph{The living form.} Although the real bound that separates the organic from the inorganic does not
immediately let itself be pointed out, the essential difference between the mineral and the plant is
nonetheless clear enough. Whereas the crystal, the dead form, goes up in the mathematical, the mathematical
(three-numbered plant, five-numbered plant, the difference-angle of the stem-leaves, and so on) is, on the
contrary, as mathematical of subordinate significance in the living forms. In the different species as well
as in the individual specimens of each species the organic runs through such a series of metamorphoses and
corresponding articulations that here one can rightly speak not merely of incessant metabolism, but also of
lower and higher stages.

The sexual system, the so-called artificial plant-system, with its classes and orders (Linnaeus), is a grand,
comprehensive, on the whole practically applicable schema; but only when a plurality of essential properties
and organ-relations is taken into the grounding of the division does the natural, i.e. the proper, system begin
to form itself. Strict holding-fast to a metaphysical principle and arbitrariness in the apprehension of the
given (Reichenbach) has, however, instead of bringing the system to completion, on the contrary shown what
difficulties are here still to be overcome.

c) \emph{The ensouled form.} Whereas the individual unity in the case of the plant either is not at all, or is
at least so protean that it vanishes into phytons and cells, the animal is, on the contrary, by force of the
concept's necessary demands (sensation, will), originally set as a definite individual. The stricter
articulation, whereby the organism's individual members are developed for themselves, advances on the whole
with a rising centralization, which gathers the organs'
% ---- printed p.94 (PDF 102) ----
different activities into the sensing, self-determining unity of life.

It is then a matter of course that the distinguishing marks of the species and the typical separation into
lower and higher orders — \emph{molluscs}, \emph{articulates}, \emph{vertebrates} — must in the animal kingdom
appear with a quite different clarity and orienting validity than analogous divisions in the plant kingdom.
Strict Empiricism (Cuvier) has here too stepped into opposition to natural-philosophical Speculation (Geoffroy
St.-Hilaire). But, in order that the system can be brought to completion, the reality of the genus-concept must
first be decided, and the connection of the higher classes with the lower determined in such a way that chief
forms and intermediate forms show themselves as naturally necessary members in the ascending chain of ensouled
unities.

If now the universally scientific is everywhere by material and task grown together with the peculiar, then each
science must also exclusively look to its ideal, and cannot possibly take its scientificity in fief from
philosophy.

How little philosophy is entitled to ascribe to itself the thought and the developments of concepts that belong
to the learned sciences is expressly seen from this, that categories such as substance, individual, species,
genus, and the like are developed in the real sciences in a quite different way than in an objective logic or
speculative natural philosophy.

The laws of thought are doubtless everywhere the same; but to hold fast to the laws of thought for themselves in
their pure generality is precisely to keep them in a questionable one-sidedness. The strict thinking has in the
one science a quite different character than in the other: the particular sciences presuppose particular talents.

\begin{center}{\bfseries \S\,21.}\end{center}

Insofar as the particular sciences have in material and aim a finite delimitation, they are subject-sciences
[Fagvidenskaber]. How peculiarly the different subject-sciences' view is
% ---- printed p.95 (PDF 103) ----
formed becomes manifest when they, by treating one and the same subject, approach mutual contact. In the
determination of soul, life, and life-principle a thoroughgoing difference shows itself between the strictly
empirical physiology, for which the organism is ``a living machinery'' that does not even need a life-force for
``machine-master'' (G. Valentin), and the psychology grounded on self-observation, which grasps the soul as
constituting essence, and lets life
well up from ``an inner source.''

The strife that can thus arise between real subject-sciences must, however, above all not
be confounded with the strife between the theory of science and a miracle-science. One
thing is science's strife with the unscientific; another is the strife between scientific
one-sidednesses. One thing is a strife against confusing appendage of an unwarranted one;
another a strife for true equilibrium between warranted ones. One thing it is that the
archaeologist and the geognost, each equipped with his apparatus of factual proofs, meet
in strife at a burial-mound; quite another that a medieval theologian, without acquaintance
with the natural marks, invokes a supernatural, and mixes into the investigation about the
human races and ``our first parents.''

In order that the individual science can advance, it must be supported by those most nearly
related to it. In that mathematics is a chief science for itself, it is at the same time an
auxiliary science for mechanics and astronomy. Chemistry, mineralogy, and geognosy unite as
auxiliary sciences in geology. The helping sciences are in turn helped by those they have
themselves helped. To what degree botany and zoology further geology, and are in turn
furthered by it, is conspicuous. In the great circle of coordinate tasks no subject has an
exclusive title to be the chief subject: in the organism of the sciences the means is aim
and the aim means.
% ---- printed p.96 (PDF 104) ----
As easy as it is to grasp the particular sciences' mutual reciprocal relation in general,
just as difficult can it in turn be to determine this relation with respect to the
individual science in particular. Physiology furnishes a telling example. The doctrine of
digestion presupposes insight into the chemical processes; how many constituents ought now
to be taken up from chemistry? The doctrine of sight presupposes optics, the doctrine of
motion mechanics; how much ought now to be borrowed from these sciences? Thus one could go
on.

Traditions and conventions have in such cases no decisive voice. Every expansion and
restriction must proceed from the nature of the matter. But precisely in the development of
the matter lies the difficulty. The one science can doubtless refer to the other; but in the
mere reference the relation is external. If the point to which reference is made lies outside
the matter, then the reference is, strictly speaking, superfluous. If, on the contrary, the
point in question is an essential part of the matter itself, then it ought also to that
extent to be taken up into the development of the matter.

But in what form is it to be taken up? As a borrowed proposition? But the borrowed
proposition is like a dead, mechanical admixture, a torn-off branch that is attached to a
foreign trunk. If the attachment is to be organic, if a grafting is to take place, then the
borrowed proposition must bring along so much of the grounding connection that thereby a part
of the one science is carried over into the other. But in this way the one science will come
to take up a multitude of extracts from other sciences; and yet every \textit{juxtapositio}
[juxtaposition] is in conflict with the essence of science.

The diverse contributions a science receives from others it must assimilate; references,
borrowed propositions, extracts, all constituents conveyed from without must be worked up in
such a way that the externality vanishes, and the patching-together is suspended.
% ---- printed p.97 (PDF 105) ----
To suspend the patchwork and work together the heterogeneous contributions requires now
doubtless a peculiar organizing talent, but then also above all a clear and distinct insight
into the relation between subsidiary subject and chief subject. This relation becomes
especially involved when sciences that in and for themselves are chief subjects partly
furnish a comprehensive foundation, partly, for a definite purpose, are treated as subsidiary
subjects.

Mathematics, which in itself is a chief subject, and for astronomy and mechanics a
comprehensive foundation, must, insofar as it serves chemistry, be treated as a subsidiary
subject. Chemistry, which in itself is a chief subject, and for mineralogy and physiology a
comprehensive foundation, must in many peculiar cases expressly be treated as a subsidiary
subject. The category foundation indicates in this connection an intermediate determination
between chief subject and subsidiary subject.

To treat a science as a chief subject presupposes, however, only one decisive regard; but to
treat a chief science as a subsidiary subject demands a plurality of regards, a doubling of
reflection that makes the task difficult.

The bound between the different sciences is in its essence dialectical; but in the
subject-sciences this dialectical cannot become the object of any peculiar exhaustive
reflection; instead of such a reflection there steps a living leading thought, a peculiar
decisive purpose, a definite plan. The aim expressed in the plan is a fundamental concept
determining all peculiar demands, penetrating all parts of the presentation. Where this
fundamental concept is not clear, where the definite aim is not at all points present to
consciousness, the presentation will become now too prolix and overloaded with unnecessary
detail, now again — of course ``for lack of time!'' — too superficial and survey-like.

The principles according to which the branches of science, with their chief and subsidiary
disciplines, are grouped under a university's
% ---- printed p.98 (PDF 106) ----
different ``faculties'' cannot possibly act solely according to abstract-theoretical regards;
for that science is in an all-too-many-sided relation to real life.

The reciprocal relation between science and life is mediated by the state and its
institutions. In the different faculties several often very diverse sciences are worked
together into peculiar wholes calculated for definite spheres of activity. If these circles
of knowledge are rounded off according to all-too-meager practical regards, science will be
violated, and the scientific spirit lost. If these wholes are organized according to
all-too-high-flown theoretical ideals, the interests of life will step into opposition to the
scientific, and science lose.

The reciprocal relation between theory and practice must stand clear to consciousness as the
plan is laid. Once the plan is laid, and the means assigned, the carrying-through ought to
happen with strictly persevering consistency. The interest now falls decisively on the
theoretical, precisely because this theoretical has taken up the regard to the practical as a
moment integral to its plan. If the study of theory can really ground a station in life, then
it also deserves that fresh vital forces be applied to bring it to its best possible
completion within fixed bounds. The student who in this respect, without doing his duty,
would filch the final advantage, is false toward science and false toward the state. For the
state will not merely have a number of businessmen as practically usable organs; it will in
these organs see the scientific spirit, the higher cultivation, represented.

\begin{center}\rule{0.35\linewidth}{0.4pt}\end{center}

% ---- printed p.99 (PDF 107): C. Philosophien og de særskilte Videnskaber / Philosophy and the Particular Sciences ----
\begin{center}
  {\large\bfseries C.}\par\vspace{2pt}
  {\bfseries Philosophy and the Particular Sciences.}\par\vspace{4pt}
  \rule{2cm}{0.4pt}
\end{center}
\phantomsection
\addcontentsline{toc}{subsection}{\texorpdfstring{C.\quad Philosophy and the Particular Sciences}{C. Philosophy and the Particular Sciences}}

\begin{center}{\bfseries \S\,22.}\end{center}

When the ancient problem how a science is attained is set forth thus: ``The knower can,
insofar as he already is knowing, get nothing to know; but neither can the ignorant one get
anything to know, since he, in order to come to a knowledge, would first have to see his
ignorance, and, in order to see his ignorance, already beforehand be a knower,'' the problem
must doubtless at first glance look like a sophism. The matter nonetheless has, in the
unconditioned sense, its correctness. One can, as regards finite expertise, get much to know
that one did not know beforehand; but from an absolute not-knowing to knowing there is no
finite transition.

As for the a priori, Plato (Phaedrus) has already pointed out the essential (ἀνάμνησις).
Every advance in knowledge is point for point conditioned by an immediately evident; since
now nothing can be evident except in that it goes together with our own consciousness, and
reminds us of what we already know, the advancing development of concepts rests on an
intellectual recollection, a reflecting recall, in which we incessantly find the True as in
ourselves, and, so to speak, take our own knowledge into possession.

He who learns mathematics gets essentially nothing else to know than what he originally
knows, and yet by learning mathematics he precisely makes advances in knowledge; the same
holds of the a priori in aesthetics, ethics, and dialectic.

With respect to the a posteriori, the advance of knowledge will doubtless be conditioned by
this, that here new knowledge is constantly gathered, that new facts are tested, new
discoveries made, and so on; but is the matter herewith decided?
% ---- printed p.100 (PDF 108) ----
If the advance of knowledge in relation to the a posteriori is grasped solely from this side,
it is easy to dismiss philosophy. If it cannot even exhaust the a priori — since mathematics
owes it no thanks — how much less can it then engage with positive facts! If the historian
sets it a task: it cannot answer, for a historical task is answered only historically. If a
chemist turns to it: it has no knack for experimenting, makes no discoveries in chemistry. If
the astronomer wants help with an observation: it scarcely understands him. If the physician
asks it for advice: it is ignorant. Even what it has on its own speculated out concerning the
relation between soul and body, the physiologist perhaps cannot with difficulty even use.

But in these affairs too it is a matter not merely of getting something new to know, but also
of knowing what one already knows, of getting one's knowledge thoroughly, inwardly,
all-sidedly brought to consciousness.

Although every science that begins with the first principles, and goes systematically forward,
doubtless develops its content in such a way that the otherwise prepared beginner can follow;
so it is nonetheless unmistakable that one first fully understands the beginning and the
earlier when one, after having come to the conclusion, looks back. The recall whereby the
foregoing is illuminated by the following is an inwardizing of the matter, a doubling of
reflection, which, with respect to finite points, does not indeed enrich us with anything new,
but nonetheless in infinite degree heightens our knowledge.

If there is now, also in relation to facts, an advance of knowledge that consists not in a
multiplication of the things that are illuminated, but in the intensity of the light that
illuminates the things; then it depends in this respect likewise on reproduction, on
inwardizing, on recol-
% ---- printed p.101 (PDF 109) ----
lection (ἀνάμνησις), and thus philosophy has to that extent the same validity in the a
posteriori as in the a priori.

In positive matters, in the affairs of experience, philosophy does not go before, but follows
after. Everywhere where Speculation with its ``intellectual intuition'' has wished to
anticipate Empiricism, lift the mystery, and unveil the hidden, it has only made itself
ridiculous. But if anyone now concludes from this that philosophy is just as superfluous when
it follows behind as it is unwarranted when it will go before; he can, by dwelling on what he
already knows, without difficulty convince himself of his mistake.

To dismiss philosophy one could adduce that the higher knowledge owed to afterthought's
reflecting retrospect is attained not by philosophizing, but by repeated study of the
subject-sciences in question. A deeper insight into the essence of mathematics is won not by
philosophizing, but by incessant repetition of mathematical analyses; a higher consciousness
of the chemical is won not by speculating, but by again and again revising what is won, and
making the recognized truths quite living, penetrating, and clear to oneself. By thus going
through science after science, everywhere with sharpened eye for presuppositions, principles,
and the inner infinity of a systematic connection, afterthought will then bring the heightened
light, and enrich us with an insight that can disdain philosophical abstractions.

If now this yield of consciousness is brought out for itself, in such a way that it can take
on the shape of the concept and be clarified in categorical form; then we have full
compensation for the vanished philosophy.

But in this compensation for the vanished philosophy, philosophy is precisely present. Every
relation to an object is in science an objective relation of consciousness. One can by
consciousness set thought above sensation, the
% ---- printed p.102 (PDF 110) ----
ideal above the real, but not deny consciousness's barrier. One can by consciousness set
sensation above thought, the real above the ideal; but one cannot by consciousness set
anything in consciousness's place. To become thoroughly, all-sidedly, inwardly conscious of
all real and all possible relations of consciousness is philosophy's task. In this its unity
of essence with the knowing consciousness, philosophy is precisely present when one proves its
non-being, precisely in exaltation when one sees through its degradation, precisely vigorous
when one considers its weaknesses — a spiritual power, a superior potence, an infinite Proteus
invincible to the subject-sciences taken together.

That positive knowledge does not without more replace philosophy becomes not so seldom evident
when learned researchers, especially natural researchers, give up the strict method, write
popularly, and impart to the cultivated an improvised philosophy.

\begin{center}{\bfseries \S\,23.}\end{center}

Over against the other sciences philosophy is doubtless a particular science, but nonetheless
no proper subject-science. Although such ideal objects as thought, truth, life, the good must
be said to be philosophical; thought is nevertheless not an object for philosophy in the same
sense as the starry heaven is an object for astronomy, the plants for botany, the animals for
zoology.

The subject-sciences are solid, always have fixed work, ample livelihood, and secure
recognition in the learned republic; philosophy is in comparison with this like an itinerant
without finite provision, a hovering shape that sinks and rises, dies away and lives up again,
according as the spirit wills it. The subject-sciences are the rich landowners who have the
objects, and live in the manifold of knowledge; philosophy is like that Eros, born of Poros
and Penia (Symposion, Plato), always in want, owns nothing, brings nothing but the
% ---- printed p.103 (PDF 111) ----
gift of penetrating everywhere, and being sated with a glance at the abundance of the rich.

He who cultivates a subject-science has, with all his exertion of spirit, nonetheless always
an unmistakable advantage in that his material has a finite bound. It is otherwise with him
who will cultivate philosophy. Instead of enjoying the impression of independence, the
intellectual αὐτάρκεια, which persevering occupation with even merely a single branch can
grant the capable specialist, a continued study will on the contrary convince the philosophizer
that he cannot take a step on the path of real cognition without feeling his knowledge's
dependence on the learned sciences; whereas the specialist will soon learn to help himself as
an independent researcher, it can be the philosophizer's humble position incessantly to knock
on a foreign door, incessantly to begin anew, and so be an apprentice his whole life.

With all its idealistic mobility philosophy is nonetheless a science determined in the concept.
If each of the subject-sciences can teach us how it will be studied, then philosophy too can for
its part offer a similar guidance.

It is natural that the study of a subject-science is united with the study of the science's
history; but insofar as here there is an immediate material, the treatment of the material will
step into the foreground. One then readily begins by studying first a part of natural science
itself, and next the history of natural science; first historical events to a certain extent,
and then the history of historiography. Science itself will in such cases stand above its own
history.

In philosophy the relation is reversed; here the history of the science is an essential content
for the science itself. The history of philosophy in union with the philosophical chief works
then forms the foundation with which here one must begin.

The next step is the study of the real sciences. To behold the works of art, live with the
poets, and make oneself
% ---- printed p.104 (PDF 112) ----
familiar with the historians requires general cultivation, but no peculiar guidance. This is,
on the contrary, the case when the talk is of the proper real sciences. That insight into the
nature of things is attained only by real investigation of the natural things is shown by the
theory of knowledge; but the mass of knowledge that must here be presupposed is so great, the
all-sidedness and the compass in which the studies in this respect would have to be driven so
considerable, that one, without looking to how much a single person can accomplish, may only
look at the science itself.

That a philosophizer should be perfectly at home in all possible sciences, and, like Aristotle
in his day, possess a whole age's apparatus of knowledge in all its particulars, is not only
impossible, but also in the matter itself superfluous. The philosophical analyses of reality
are doubtless in so intimate connection with the facts of experience and corresponding
observations that handbooks, surveys, popular lectures, and the like do not suffice; but it is
nonetheless also a law that philosophy must not take up more material than the Dialectic can
dissolve, and the concept assimilate. It then depends not so much on the greatness of the mass
of knowledge as on choosing among the mass's parts, sifting the material, and, so to speak,
finding the ore from which a knowledge can be smelted out.

As an inviolable presupposition the theory of knowledge must demand that the philosophizer,
even if he is not a mathematician, nevertheless learn to understand the language of mathematics,
and, although he will not become a chemist, make himself familiar with the chemical laws,
designations, and formulas. Without such a preparation the philosophizer will never penetrate
into the kernel of the matter, never see through the exact adducing of proof, never understand
what physics really means when it speaks of force and matter. Mathematics and chemistry are, so
to speak, the keys to the secret recesses of real science; only he who has these keys is an
initiate, and only for the initiate does the treasure show itself.
% ---- printed p.105 (PDF 113) ----
In order that a study of so many and different branches of science shall now not either lose
itself in bottomless thoroughness, or pass over into loose dilettantism, it is naturally
necessary that the philosophizer be conscious of a task, that he in all this see a definite
goal, and have the endurance to pursue it.

As for the task in its generality, it will be determined by the reciprocal relation between the
real and the ideal, between analysis of representation and analysis of concepts, between the
principles and the finite presuppositions.

As for the task in its peculiarity, partly an individual question arising under peculiar
circumstances, partly a comprehensive regard to philosophy's organic chief parts can determine
the compass and direction of the real studies.

But even when the task is held fast, even most strictly, the real studies still bring
considerable difficulties peculiar to the relation with philosophy. For an analysis of concepts
it is namely not enough that the philosophizer knows a number of the points essentially attractive
to thought; he must have such points in such a manifold, possess them with such freedom, and move
with such sureness, that he, without striking against the real, can move his moments of concept as
though they were not piecemeal-gathered contributions, but inventions of his own fantasy.

How difficult it is to acquire such sureness can easily be seen. From many sides Empiricism
illuminates the inner, mysterious connection that is between things' otherwise so different
properties. The universal is here known in the special, but this special in turn interrupts the
cognition of the universal. Let one, e.g., order the elements in the so-called trilogies; let one
compare barium ($A$), strontium ($B$), calcium ($C$), and note the relation between the chemical
numbers: $\frac{A+C}{2}=B$. If one now looks at the solubility of
% ---- printed p.106 (PDF 114) ----
baryta, strontia, and lime, then one finds that the baryta has the easiest, the lime the heaviest
solubility; whereas the strontia stands midway between. Let one test another trilogy: sulphur
($A$), selenium ($B$), and tellurium ($C$), whose chemical numbers will likewise give:
$\frac{A+C}{2}=B$. If one now looks at the specific gravities, these relate thus: the sulphur has
the greatest, the tellurium the least, and the selenium stands midway between. Let one think of
the relation between the specific heats and the chemical numbers; let one think of the crystals'
relation to light and their expansions under heating; let one think of the law according to which
a rise and fall of the boiling point of the polymeric hydrocarbons ($\mathrm{CH_2}$,
$\mathrm{C_2H_4}$, $\mathrm{C_4H_8}$, $\mathrm{C_{16}H_{32}}$, \ldots) takes place, in that the
point rises for each carbon atom that is added, and sinks for each added hydrogen atom, and so on.

But while here, under this penetration into the interior of things, a multitude of parallels, a
series of presuppositions and consequences, come to light, there can nevertheless in this connection
scarcely be made a single a priori inference independent of the test of experience. That the
specific gravity of iron oxide is less than that of iron is comprehensible; but, if one drew from
this the general inference that the higher
% [errata applied (Rettelser p.106): print ``oxides''; read ``metal oxides'']
metal oxides on the whole must have a smaller specific gravity than the lower, one would err: the
specific gravity of stannic oxide is, e.g., greater than that of stannous oxide; that of plumbic
oxide greater than that of plumbous oxide. If we look at the relation between the two properties,
hardness and specific gravity, then it seems natural that on the whole they should correspond to each
other: the diamond, which is far harder than graphite, really also has greater specific gravity;
yet let one guard against a hasty inference.

When isomorphic substances vicariate, it is seen that hardness and specific gravity by no means
straightforwardly correspond to each other. If iron steps into the garnet instead of lime, or lime
instead of iron; then it holds: the more iron, the greater specific gravity, but conversely: the more
lime, the greater hardness. Every
% ---- printed p.107 (PDF 115) ----
essential connection doubtless contributes to the answering of a Wherefore; but here comes every
moment a Wherefore that cannot yet be answered.

For the philosophizer the subject-knowledges cannot in and for themselves have significance, and
yet it is not so seldom that precisely those points which are quite properly fruitful for the
analysis of concepts lie hidden under a detail of facts and special investigations, which thus
must nonetheless be gone through.

Finally a deeper appropriation of the intellectual that offers itself in a definite circle of
facts can first set in when this intellectual is illuminated from other and corresponding
circles: \textit{Opposita juxta se posita magis illucescunt} [opposites placed side by side
illuminate each other more].

If the total can never be fully developed except by means of the individual, then the reverse
also holds, and first and last in philosophy, that a penetrating study of the individual cannot
possibly be carried through without a corresponding insight into the total.

\begin{center}{\bfseries \S\,24.}\end{center}

But even if the philosophizer has now also conquered the difficulties of real study, drawn
sublimates from the positive sciences, and sees himself able to work together the constituents:
in what form is he then to impart what has been won? to whom ought he to turn, by whom will he
be understood? For the illumination of the difficulty here pointed out we grasp the matter,
first from an objective side, in that we look at the manner of presentation, next from a
subjective side, in that we take account of understanding.

\emph{The manner of presentation.} As for facts and corresponding laws, philosophy must treat
the given as a given, and not by a kind of sophistical turn let the matter get the appearance
that things which are due solely to experience nonetheless proceed with inner a priori necessity
from the pure concept. The untruth
% ---- printed p.108 (PDF 116) ----
that in this respect crept into the Hegelian method has doubtless its ground not in any conscious
dishonesty, but in overstrained zeal of reflection and a consequent disproportion to the immediate.

Immediacy ought doubtless to be lifted; but immediacy must not be suspended in such a way that the
given changes its nature, and is transformed from a positive fact into a metaphysical essence;
immediacy must be suspended in such a way that the given, as given, furnishes an ideal contribution
to the comprehensive development of concepts that the principle itself sets going. Insofar as the
comprehensive development of concepts is and remains the main thing, the given then lays off its
contingency, does not remain in thoughtlessness, but shows its ideal, obeys the concept, and comes
when thought calls.

The point at which the conceptual side of the given each time coincides with the leading concept of
the development is a point of insight. At this point there will always gather so much light that the
proposition which is taken up as a given from another science neither presents itself as a dead
borrowed proposition, nor brings with it, for its final explanation, a circumstantiality that halts
the pace of development and burdens thought with the all-too-mediate.

To conduct separate proof for each peculiar proposition is impossible and unnecessary; to illuminate
each proposition in such a way that the side that is essential for the connection is also clarified
in the connection, that the point on whose understanding it each time depends can also each time be
understood, is possible and necessary.

Every scientific presentation of deeper content always has a multitude of luminous points in the
foreground, a multitude of obscure points in the background, and between both manifold transitions.
Where several sciences treat the same thing, each will make its own distinct; but they will never
make the same distinct in the same way. Likewise with respect to clarity. What one development
according to its
% ---- printed p.109 (PDF 117) ----
purpose must lift to clarity, another may according to its purpose let remain in obscurity. The
logical categories obscurity, clarity, and distinctness are of dialectical application.

Another difficulty with respect to the given is peculiar to philosophy. This science may not, as
a positive subject-science does, give itself over to the manifold of the finite, enumerating,
narrating, describing, and so on. Nonetheless it must in its way take up the outwardly manifold,
insofar as it will give itself a real content.

To remedy this difficulty one has sometimes seized on the desperate expedient of laying out a
schema of concepts, holding to the formal order, letting the reasoning within each of the rubrics
run free, and then filling the emptiness with a multitude of heterogeneous examples fetched from
all positive sciences; this is perhaps the worst method anyone can apply.

The task is to let the form take up the content in such a way that the development of concepts
becomes a development of the matter. In this development of the matter the positive elements must
then be chosen, distributed, ordered, and treated not merely in such a way that the individual
facts are placed where they are singly needed, but also in such a way that the whole aggregate of
real moments is worked through by reflection, that each concluding totality can express the content
of the matter in the form of the concept. Through concluding totalities in different — partly
coordinate, partly mutually enclosing — circles the development will then go to the last conclusion.

\emph{Understanding.} Where tradition and habit once give a certain manner of teaching prestige
and support, it is often striking to what degree one can suppose one explains the inexplicable and
understands the incomprehensible (half-scientific religious books); where, on the contrary, the
principle of cognition demands that in something handed down several changes necessary for
understanding must be made, it is again just
% ---- printed p.110 (PDF 118) ----
as striking to what degree obstinacy can strive against, and find the in-itself-understandable
incomprehensible. That opinions about the understandable and the incomprehensible can be so divided,
that one can make such highly different demands on what is to be understandable, follows from this,
that the category understanding is so infinitely dialectical.

It is an easy matter to reproach a philosophy with incomprehensibility. General cultivation does not
as a rule go deep enough to follow a stricter analysis of concepts: the understanding of the
individual presupposes expertise. But if special expertise is an inviolable demand, then the
consequence of this must be that each specialist holds to his piece, and lets the rest stand as an
incomprehensible. That such a parceling-out of understanding must make philosophy meaningless is then
understandable.

As for the presentation, it can become incomprehensible just as well by holding to the general as by
going into the special. If special analyses are predominant, then one doubtless notices the involved,
the thorny, the prickly; but if the universal-forms are predominant, then one gets perhaps still more
troubles with ``the abstract,'' ``the nebulous,'' ``the hovering,'' ``the ambiguous'' (turns of
phrase, word-philosophy, phraseology). It depends in such things not on the finite What, but on the
infinite How.

To clarify to oneself the relation between understandable and incomprehensible, one must, on occasion
of the demands that can with right be made on the presentation of what is to be understood, then above
all also not forget the demands that can with right be made on him who will understand, such as:
necessary preparation, persevering attention, scientific sense.

Insofar as each science has its peculiar structure of thought, it consequently also has its peculiar
kind of understandability. In the mechanical, understanding comes otherwise than
% ---- printed p.111 (PDF 119) ----
in the chemical, in the exact otherwise than in the dialectical, in mathematics otherwise than in
philosophy.

\begin{center}{\bfseries \S\,25.}\end{center}

If we call investigations of things' first presuppositions and knowledge's principles philosophical,
there can doubtless be no doubt that philosophical investigations are necessary in the different
sciences; but from this it does not yet follow that philosophy for itself is a peculiar science. One
must, e.g., in every branch of science apply criticism, and also has in the different sciences a
multitude of critical treatises, even treatises on criticism; but for that reason surely no one will
maintain that criticism for itself is a peculiar science.

If there were now a circle of critical-philosophical treatises in which the diverse sciences had each
in particular given an account of its principles and presuppositions; then one could doubtless, by
working them all up according to one plan, unite them in a whole; but this whole would, by the finite
connection of the parts, nonetheless just as little become a peculiar philosophical science as the
collected treatises would become a peculiar science by being bound in one volume.

The objective Idealism would now at once lift philosophy's scientific independence by referring a) to
the Dialectic in which the finite contributions are dissolved, melted down, and deposited as the
mobile moments of an inner infinity; b) to the immanent method whereby the absolute content gets its
moments laid out, separated, and gathered in the absolute form; finally c) to the deeper view of life,
the solution of the riddle, which is attained only in Speculation.

But the presuppositions on which philosophy not long ago wished to ground its divine superiority unite
here as flaws that threaten its scientific independence. Thus it holds with
% ---- printed p.112 (PDF 120) ----
respect to a), that the dialectical dissolution may nonetheless not loosen principles and
presuppositions from the soil of reality in which they have their root; with respect to b), that ``the
immanent method'' must go the same way as ``the absolute knowledge''; and finally with respect to c),
that science does not concern itself with ``the deeper view of life,'' since it on the whole does not
sit in judgment over views of life, and that ``Speculation'' cannot solve ``the riddle of existence.''

How far philosophy is a science independent in itself is not dependent on arbitrary opinions or passing
blunders, but on the nature of the matter. If philosophy is a science, then the peculiar form in which
it develops its content must be of such decisive significance that scientificity would suffer an
irreparable loss if there were no work on bringing this form too to systematic completion. We must then
here expressly bring out the influence philosophy exercises on the scientific life of consciousness,
insofar as it is a matter of the appropriation of the material, the liberation of the concepts, and the
inwardizing of the matter.

1) \emph{Appropriation: an assimilating activity.} The outwardly given does not immediately go up in
consciousness. In every subject-science there will always be a mass of constituents so heterogeneous
with the life of consciousness that not even the strongest memory can gather and bear their quantity;
such constituents are then preserved as materials in magazines, repertoires, and the like learned
storerooms.

In relation to the real, scientific insight is conditioned by memory; in relation to the essential, by
recollection; the real is remembered, the essential recollected. The appropriation under which the
scientific material passes over into the flesh and blood of consciousness (\textit{in succum et
sanguinem} [into sap and blood]) is an intellectual process of digestion that works up the material,
and takes on a different character according as it is the
% ---- printed p.113 (PDF 121) ----
real in its particulars or the essential with its universal-determinations that becomes the main thing.

The more the side of reality predominates, the more closely thought joins itself to memory, and the
more also the appropriation becomes restricted to a peculiar rubric of consciousness, determined for
the peculiar subject. The more, on the contrary, the side of essence is brought out, the more intimately
thought goes together with recollection; the finite constituents of the material then become real
moments in the peculiar concept, and the appropriation of the peculiar concept sends out a light that
shines through the whole consciousness.

When now these two modes of appropriation are referred each to its aim and its principle, a difference
shows itself between subject-science and philosophy.

The carried-through assimilation relates to Speculation as real-philosophy to Idealism. In idealistic
Speculation the given is volatilized into a priori essence; in real-philosophical assimilation the a
priori essence is peculiarly determined by the given.

2) \emph{Liberation: a humanizing activity.} Restricted to the different rubrics of consciousness, the
won concepts are still so grown together with finite constituents, still so burdened with particular
presuppositions, that they find themselves as in a state of unfreedom; they are here held as the
material's prisoners within a narrow circle of vision; they are, so to speak, the subject's serfs, who
may not leave the definite workplace. This intellectual slavery of the concepts forms only a one-sided
counterweight to the equally one-sided unboundedness that proceeds from superficial abstractions, and
serves the ingenious play of the association of ideas. Every concept that is had completely is had with
freedom; it then depends on seeing wherein the freedom consists.

That the concept of the mechanical belongs just as originally under mechanics as the concept of the
chemical under chem-
% ---- printed p.114 (PDF 122) ----
istry is incontrovertible; when now Hegel in his Logic nonetheless deals with the concept's own
mechanism and its logical chemism, then it is neither mechanics's nor chemistry's, but rather
philosophy's task to investigate with what right such a transfer has happened. Word-explanations, acute
interpretations of homonyms, synonyms, metaphors, and so on, do not mean much in this case. Here the
question is not of word-explanations; here the question is expressly of the changes in the categories'
validity that are brought about by a passage through different spheres. Such investigations demand
expertise and Dialectic. In that a concept has passed over from one sphere to another, it is no longer
the same, and yet still the same: ``function'' in mathematics and ``function'' in physiology is not the
same, and yet the same; ``approximation'' in morals and ``approximation'' in mathematics is not the same,
and yet the same; ``reflection'' in optics and ``reflection'' in logic is not the same, and yet the same.
Here one can continue to infinity.

Philosophy's humanizing activity does not immediately coincide with a striving toward the popular. The
dialectical liberation of the concept, which at every point must be exactly conditioned by what is given
by assimilation, belongs exclusively to science. This humanizing activity, a necessary condition for
consciousness's being able in all things to be itself, and to master its many rubrics, is possible only
in that philosophy steps forth independently, and lifts its own method.

3) \emph{Inwardizing: a totalizing activity.} What difference there is between knowing something when
consciousness especially enters into the learned rubric, and knowing something in such a way that on a
given occasion it instantly presents itself, and stands clear to thought, is generally known. The same
content of knowledge can in the one form fall so heavy that each time a peculiar exertion of
% ---- printed p.115 (PDF 123) ----
force is required when the knower is to lift it; in another form, when it, without losing its
intellectual specific gravity, has by assimilation been separated from the all-too-circumstantial
admixtures, it can on the contrary be so light that consciousness, freed from the material's burden,
bears its wealth as though it bore only itself: \textit{omnia mea porto mecum} [all that is mine I
carry with me].

The totalizing activity is conditioned by a corresponding method. To the finite separation of thinking
and being corresponds a \emph{critical} method; reflection moves about its object, points out the
dogmatic proposition's untenability, and makes a subjective transition from the one to the other.
Thought moves, but the thing stands still; here one thinks about the object, but the object does not
come to unfold its moments: here the immanent Dialectic of the development of the matter is lacking. To
Idealism's absolute knowledge corresponds a \emph{dialectical} method. Here it is too little to think
about the matter, for the matter is thought through in such a way that it looks as though it thought
itself. The objective then does not properly express the real; it is an abstraction in which the
subjective comes objectively to itself. The objective is not an existing object, it is a concept; it goes
by ``der Begriff'' [the concept] and ``der Begriff''; the pure concept entertains itself: here criticism
is lacking. To a synthetic unity of thinking and being, a unity of essence that cannot suspend the real
difference, corresponds a method in which objective development of the matter and subjective reflection
determine each other mutually; such a method is \emph{critical-dialectical}: criticism clarifies the
finite conditions, Dialectic opens up the inner infinity.

The methodical movement is in philosophy a double one. On the one side the special in the content, the
finite with its particulars, is worked up and assimilated in such a way that the moments go back into the
universal, that the fundamental cognition is enriched, and great totalities are set in
% ---- printed p.116 (PDF 124) ----
motion. In this respect the Hegelian method has an imposing superiority. But since the real, with its
finite members, its angular forms, and existential punctuality, necessarily belongs to the development of
the matter, it is on the other side an equally essential task to analyze the individual, separate the
members, penetrate the entanglements, and bring the individual to dialectical vibration. In this respect
the Kierkegaardian analyses furnish apt examples. On the heights of totality the telescope is applied,
for immersion in the analysis of punctuality the microscope is used.

In the outer world two such opposed modes of viewing cannot possibly be applied simultaneously; to look
at once through a telescope and a microscope is doubtless an impossibility. In knowledge's inner world the
matter stands somewhat otherwise. If already the subject-science, with its whole detail and its rubrics,
can work itself up to higher reciprocal relations of insight into the individual and survey of the whole,
then philosophy, with its assimilating power and liberating activity, must be able to carry it still
further. The critical conditions the dialectical. The more dialectical insight and survey — insight into
the special and survey of the total, insight into the total and survey of the special — alternate with
each other, and in the dialectical refraction awaken each other, sharpen each other, and penetrate each
other, the greater is the intellectual tension, the more the concept's force grows, the higher the waves
of the movement beat, the stronger the stream of cognition goes.

Method relates to system as vital activity to the living organism. In the subject-science, where it
always depends on getting the material placed, the schema can play a role. In philosophy, which is not
burdened by material, schematism is intolerable. Every principled thought is at bottom always systematic;
if the systematic form will not at once present itself, the thought ought to be developed with freedom; it
% ---- printed p.117 (PDF 125) ----
will then afterward be seen from the development in what comprehensive connection it belongs.

The finite knowledge has many particular points of departure and many particular systems; the
``absolute'' knowledge will develop all from one principle in one continuous system. The synthetic
knowledge, which can neither be satisfied by relative points of departure, nor reach the absolute
principle, gets no absolute system, but within its bound an essential relation of relative systems.

The more the individual elements of knowledge are worked together, organized, and clarified, the more
intimate the connection between the whole and its individual parts becomes, the more intimately also the
content of knowledge goes together with the knowing consciousness. The assimilating and humanizing
activity that is here demanded shows clearly that the complete inwardizing is dependent on an independent
philosophy.

If we grasp the particular sciences as definite organs in the great organism of scientificity, then we
must evidently grasp philosophy as the soul of the organism. Here too it will then repeat itself, that a
rising separation and more definite development of the individual organs goes together with a
corresponding concentration. The principled unity of the functions will here too become more independent,
according as it becomes more intimate, and, though in the finite sense a product of the organs'
cooperation, nonetheless with inner infinity assert a subsistence in itself over against the particular
organs.

\begin{center}\rule{0.35\linewidth}{0.4pt}\end{center}
% ---- printed p.118 (PDF 126): III. Ideelære / Doctrine of Ideas ----
\begin{center}
  {\large\bfseries III.}\par\vspace{4pt}
  {\Large\bfseries Doctrine of Ideas.}\par\vspace{4pt}
  \rule{2.5cm}{0.4pt}
\end{center}
\phantomsection
\addcontentsline{toc}{section}{\texorpdfstring{III.\quad Doctrine of Ideas}{III. Doctrine of Ideas}}
\markboth{Doctrine of Ideas}{Doctrine of Ideas}

\begin{center}{\bfseries \S\,26.}\end{center}

What is an idea? Every science presents a content of the Idea, a side of the Idea, a peculiar
idea, but mediately; only philosophy sets itself to present the Idea as Idea. One says in daily
speech that something is so unknown to one that one cannot even form an idea of it; one speaks of
giving the idea for a building, the idea for an undertaking, and so on. A certain
general-representation, a plan, a living leading thought are the categories that, according to
this usage, hover before consciousness. Among the Sensualists the word ``idea'' occurs in a loose,
superficial sense.

Although Kant grants the Idea only subjective validity, he has nonetheless most forcefully
asserted its right and pointed to Plato's merit with respect to the doctrine of Ideas
(\textit{Kritik der r. Vern.} S. 289.)

``Plato made use of the expression Idea in such a way that one sees well that he understood under
it something that is not merely not borrowed from the senses, but that even far surpasses the
concepts of the understanding with which Aristotle occupies himself, since in experience one
meets nothing congruent with it\ldots{} Plato remarked very rightly that our power of cognition
feels a far higher need than to spell out phenomena according to synthetic unity in order to be
able to read them as experiences, and that our reason in a natural way soars up to cognitions that
go far beyond anything that any object experience can give could ever congrue with, but which
nonetheless have their reality, and are by no means mere figments of the brain.''
% ---- printed p.119 (PDF 127) ----
``He who would draw the concept of virtue from experience, he who would make (what indeed many
have really done) that into a model and source of cognition which in each case can serve only as
example, he would make of virtue an ambiguous absurdity, changeable according to time and
circumstances, usable for no rule. On the contrary, everyone notices that when someone is presented
to him as a model of virtue, he nonetheless always has the true original, with which he compares
this alleged model, only in his own head, and merely values it accordingly.''

``The Platonic Republic has, as an allegedly striking example of dreamt perfection that can be found
only in the idle thinker's brain, become a proverb, and Brucker finds it laughable that the
philosopher maintains that a prince would never rule well if he did not become a partaker in the
Ideas. But one would do better to trace this thought more closely (where the excellent man leaves us
without help), and by new exertions set it in the light, than, under the very wretched and shameful
pretext that it is not executable, set it aside as useless.''

With Plato, then, Kant agrees that the Idea expresses the unconditioned, that which is exalted above
experience. ``The concept is either an empirical or a pure concept, and the pure concept, insofar
as it then has its origin in the understanding (and is not a pure sense-image), is called
\textit{Notio}. The Idea, or the concept of reason, is a concept of notions that surpasses the
possibility of experience. He who has once accustomed himself to this distinction must find it
intolerable to hear the representation of the red color named as Idea.''

A significant contribution to the doctrine of Ideas Kant has furthermore made in his ``\textit{Kritik
der Urtheilskraft}.'' If the theoretical reason has taught us to know a world in which all goes
according to natural laws, if the practical reason has opened us an insight into a world of ethics,
in which all is determined by free-
% ---- printed p.120 (PDF 128) ----
dom, then there would be an all-too-insurmountable chasm between the realm of nature and the realm
of freedom, if this chasm were not filled in by the judgment's setting up a ground of unity for both
realms. (Anf. Skr. \textit{Einleitung}.)

``The judgment, which is the faculty of thinking the particular as included under the universal, has
for its object the concept of nature's \emph{purposiveness}. The purpose (aim) is, namely, the
supersensible unity that contains the ground of the objective's reality; every realization of a
purpose is bound up with pleasure: the judgment contains the laws for the feeling of pleasure and
displeasure.''

``Nature's purposiveness can be represented either subjectively or objectively. In the first case I
sense pleasure and displeasure immediately at the representation of an object, before I have formed a
concept of it; my joy then refers only to a harmonious relation of purpose between the object's form
and my power of intuition. The judgment in this subjective sense is called the \emph{aesthetic}. In
the second case I form beforehand a concept of the object, and now assess how far the object
corresponds to this concept. To find a flower beautiful for my power of intuition, I need have no
concept of it; but to find the flower purposive requires a concept. The judgment in this objective
sense is called the \emph{teleological}.''

``The objects of nature please us not only by a subjective purposiveness; they are at the same time
in a purposive relation to one another mutually. The objective, material purposiveness is of a
twofold kind: the \emph{outer} and the \emph{inner}.''

``The \emph{outer purposiveness} is, since it designates only a thing's utility for another, something
merely relative. Snow secures the seed in cold lands against frost, and eases men's communication by
sleds. The sand deposited by the seashore is beneficial for fir forests. The means is in this outer
purposiveness not in and for itself, not
% ---- printed p.121 (PDF 129) ----
according to its concept, but only contingently purposive. The deposition of the sand on the seashore
is not grasped by saying it is a means for the fir forests; it is understandable in and for itself,
quite apart from the concept of purpose: the outer, relative purposiveness lets itself also be grasped
out of \emph{natural mechanism alone}.''

``It is otherwise with the \emph{inner purposiveness}, which consists in this, that something is in
itself both means and end, that its end therefore does not fall outside. The organic natural product
is such a one, in which each of its parts is aim and at the same time means or instrument. In
propagation the natural product is produced as species; in growth the natural product is produced as
individual; in formation each part of the individual produces itself. This natural organism does not
let itself be explained by merely mechanical causes; it must be explained by purpose-causes, it must
be explained \emph{teleologically}.''

``The earlier systems have treated the concept of natural purposiveness dogmatically; they have
either ascribed (Hylozoism, Theism) or denied (Fatalism) to nature as thing in and for itself this
purposiveness. But we, who hold that teleology is only a regulative principle, decide nothing
concerning the question whether inner purposiveness belongs to nature in and for itself, or not. Our
understanding thinks discursively, always starts from the parts, and grasps the whole as result of
its parts; it can therefore not grasp the organic natural products — in which, conversely, the whole
is the parts' ground of coming-to-be and prius — in any other way than under the point of view of the
concept of purpose. If, on the contrary, there were an intuitive understanding (\textit{intellectus
archetypus}) that in the universal knew the particular, in the whole the parts as already
co-determined, then such an understanding would grasp the whole of nature out of one principle,
without needing the concept of purpose.''
% ---- printed p.122 (PDF 130) ----
As teleological principle the Idea is essentially characterized; the next step is to get this
subjective Idea to hold also as objective, to transform this in Kant merely regulative principle into
a constitutive one, a principle developing its content out of itself, and ``the Absolute'' will be
reached.

How Schelling has prepared this transition Hegel acknowledges (\textit{Geschichte der Philosophie, 3
Th., Berlin 1844, S. 613}) ``In the first place, in Schelling the Idea itself is to be brought out,
that he has grasped the True as the Concrete, as unity of the subjective and the objective. It is thus
the main thing in the Schellingian philosophy that in it it is a matter of this deep speculative
content, which, as content, is the content of which it has been a matter in the whole history of
philosophy. The thinking that is freely for itself, not abstract, but in itself concrete, grasps
itself in itself as the intellectual-real world; and this is nature's truth, nature in itself. The
second great thing in Schelling is to have demonstrated the forms of spirit in nature: electricity,
magnetism, and so on, are to him only the Idea's outer modes. The defect is that this Idea as such,
its separation into the ideal and the natural world, and thus the totality of its determinations, are
not pointed out and developed as they follow with necessity from the concept. In that Schelling has
not grasped this side, thought has in him been lost; the work of art thus becomes in him the highest
and only mode in which the Idea is for spirit. But the Idea's highest mode is much rather its own
element; thought, the conceived Idea, is thus higher than the work of art. The Idea is the truth and
all that is true is Idea; the Idea's systematizing into world must be proved as necessary unveiling
and revelation.''

If the philosophy of the Absolute had now really mastered its task of proving how the Idea
systematized itself into world, then not only would philosophy in all essentials be concluded, but
with the Idea's ``unveiling'' and ``revelation'' the riddle of existence would at the same time be
solved once
% ---- printed p.123 (PDF 131) ----
and for all. How little this is the case has already been shown in the preceding sections.

In his ``\textit{Kritik der Urtheilskraft}'' Kant sets up an antinomy under the following form:

``\emph{Thesis}: All productions of material things are possible according to merely mechanical laws.''

``\emph{Antithesis}: Some productions of the same are not possible according to merely mechanical
laws.''

This antinomy is still so far from having found its solution in favor of the cognition of an objective
Idea, that an interpretation of reality according to merely mechanical laws rather threatens to rob
the Idea as Idea of all objective reality.

``Sound sense,'' which feels itself drawn to both sides, can decide nothing; we get — herein Kant is
incontrovertibly right — in this respect antinomic impressions of the objective world.

In astronomy, e.g., the motions of the heavenly bodies are indeed explained by mechanical causes;
nonetheless we receive, by looking into the structure of the world, an impression as though reason
itself were the master-builder who used the individual in order to realize its plans in the whole, an
impression as though we saw the Idea.

The physiologist must undeniably strive, as far as possible, to explain the expressions of life from
the standpoint of ``natural mechanism''; nonetheless ``sound sense'' cannot withdraw itself from the
impression of the purpose for whose sake the organs are formed. It looks as though the fundamental
parts with their unchangeable character were the real first, and life itself a blind product of these
parts' law-governed reciprocity; and it looks again as though life itself were the proper prius, and
the parts came into being only to serve the whole; the whole looks like a realized Idea.

We look at the individual members, and we lose the Idea; we look at the inner whole, and the Idea
stands living before us. We
% ---- printed p.124 (PDF 132) ----
dissect a plant, and we lose the Idea; we let the plant, as it lives and grows, make its typical
impression, and we meet the Idea. We analyze the proportions of measure in one of Raphael's pictures
mathematically, and we lose the Idea; we let the mathematical vanish into the aesthetic total
impression, and we win the Idea.

This antinomic is now so much the harder as we, according to what has here already been shown, can
neither with Kant take refuge in the merely subjective, nor with Hegel deny the real in order to
think the objective. The problems belonging here will in this connection require a peculiar treatment
of \emph{the Idea in its Essentiality}, of \emph{Reality against the Idea}, and of \emph{the Validity
of the Idea}.

\begin{center}\rule{0.35\linewidth}{0.4pt}\end{center}

\begin{center}
  {\large\bfseries A.}\par\vspace{2pt}
  {\bfseries The Idea in its Essentiality.}\par\vspace{4pt}
  \rule{2cm}{0.4pt}
\end{center}
\phantomsection
\addcontentsline{toc}{subsection}{\texorpdfstring{A.\quad The Idea in its Essentiality}{A. The Idea in its Essentiality}}

\begin{center}{\bfseries \S\,27.}\end{center}

``The most heterogeneous things Plato seems to gather under the name Idea. He refers thereto ideal
(aesthetic and moral) concepts, e.g., the Idea of the Good, the Idea of Injustice, pure thoughts,
such as the thought of unity and plurality, concepts of living species (e.g., the Idea of a bull), of
individuals, of their states and expressions of life, the Idea of voice, of color. Indeed, he speaks
of the Idea of human manufactures, of the Idea of a bed, of a table. No object of cognition seems
then to be excluded from participation in the Ideas, and thus it really stands with it. All that has
reality, existence, can be thought in such a way that it is referred under an Idea. Therefore
% ---- printed p.125 (PDF 133) ----
Parmenides corrects the young Socrates's opinion when he hesitates to know Ideas in things that are
regarded as low and insignificant, e.g., hair or filth, since such things are low and insignificant
only from a restricted human standpoint.'' (Poul Møller).

That Plato's doctrine of Ideas stands in intimate connection with his doctrine of knowledge and the
essence of things is incontrovertible, and often enough brought out. A reference to the demand of
knowledge and science in general is, on the contrary, insufficient, since there can undeniably be
science and scientific penetration into the essence of things without Plato's Ideas being for that
laid at the foundation. It must expressly be brought out that it is only the Platonic science that is
inseparable from the Platonic Idea, and that Plato himself gave this science (ἡ περὶ τὸ ὂν καὶ τὸ
ὄντως καὶ τὸ κατὰ ταυτὸν ἀεὶ πεφυκὸς ἐπιστήμη [the science concerning being and the truly-existent
and that which is by nature always self-same]) the name: Dialectic. One cannot therefore treat
Dialectic for itself, and the doctrine of Ideas for itself; they relate as form and content: the
Platonic Dialectic and the Platonic Idea are, essentially seen, two sides of the same.

\emph{1) Ideas and concepts.}

Although Dialectic sees the essential, the universally valid, the True only in the concept (οὗ αὐτὸς
ὁ λόγος ἅπτεται τῇ τοῦ διαλέγεσθαι δυνάμει [which reason itself grasps by the power of dialectic]),
one cannot nonetheless say that Dialectic occupies itself only with concepts, and not, on the
contrary, with representations, sense-impressions, or sensible things. If Dialectic held only to the
concepts insofar as these arose by removal of finite representations, i.e. by abstraction, the Idea
in the objective, scientific sense would itself be an abstraction.

Here is a thoroughgoing difference between the concepts of abstraction and the dialectical realities
Plato calls Ideas. The abstracting thought lets the finite, the sensibly individual, remain standing
as the in-itself-inde-
% ---- printed p.126 (PDF 134) ----
pendent valid common mark, the empty, dependent universal-form. In this sense it is then not the
bull's concept, but the bull-individual, that is the real. For the dialecticizing thought the finite
in its finitude has no reality. Instead of abstracting from the representation, and letting the finite
stand as the properly existing, Dialectic engages with the representation, and shows how the finite,
the sensible in its becoming and changeability, is an in-itself-untrue, a fleeting, a self-consuming,
a semblance whose nothingness thought sees through. The more the finite things show themselves in
their untruth, the more the concept shows itself as the True, as the unchangeable in the changes, as
the self-valid that bears and determines all finite existence, i.e. as Idea. Only when Dialectic has
consumed the bull-individual's immediate, finite reality is the concept itself seen as the producing
principle, as the eternal reality, as ``the bull's Idea.''

Of the sensibly becoming there is no knowledge (τῶν ῥεόντων οὐκ ἔστιν ἐπιστήμη [of things in flux
there is no science]), consequently the Idea must show itself as the being; but the being in its
abstract, differenceless unity is in infinite contradiction with itself, thus the Idea's being must
have in itself an essential difference, take on movement, and show itself as a plurality of beings.
The former is established against the Heraclitean (Theaetetus), the latter against the Eleatic
(Sophist) one-sidedness.

Attaching himself to the Heraclitean doctrine of becoming, the Sophist Protagoras had set up the
proposition that man is the measure of all things (πάντων χρημάτων μέτρον ἄνθρωπος, τῶν μὲν ὄντων ὡς
ἔστι, τῶν δὲ οὐκ ὄντων ὡς οὐκ ἔϛιν [man is the measure of all things, of the things that are that
they are, and of the things that are not that they are not]). To becoming corresponds sensation; when
one feels it warm, it is warm for him; if another finds it cold, it is cold for him. Sense-sensation
is the only mark of truth. This view is described by the Platonic Socra-
% ---- printed p.127 (PDF 135) ----
tes in the dialogue Theaetetus. If man merely as sensing individual were the measure of things, then
an animal could just
as fully as the wisest man, indeed as a god, know the True (ὅτι πάντων χρημάτων μέτρον ἐϛὶν ὗς ἢ
κυνοκέφαλος ἤ τι ἄλλο ἀτοπώτερον τῶν ἐχόντων αἴσθησιν [that the measure of all things is a pig or a
dog-faced baboon or some other still more absurd creature that has sensation]). If all men have equal
insight into the True, then it is a contradiction to give oneself out as a teacher of wisdom
(σοφίζειν); Protagoras must indeed grant everyone right who gives him wrong. He who assumes that not
every man is the measure of things, and that sense-sensation is not science, is just as much right as
he who assumes the opposite; thus it is, according to Protagoras's own doctrine, also true that his
doctrine is not true. Protagoras makes every knowledge of the future, and therewith of the useful,
impossible (ἔστι δέ που καὶ περὶ τὸν μέλλοντα χρόνον τὸ ὠφέλιμον [and the useful, surely, also
concerns future time]). Even if one assumed that every judgment about the present were equally
correct, one could nonetheless not assume that every judgment about the future were equally correct,
that what everyone with respect to agriculture or the healing of an illness happened to assume to be
useful would also really show itself to be the useful. How absurd it is to regard sense-sensation as
science is finally seen from this, that the infinite becoming must make every judgment impossible. The
sensing one and the sensed are both in motion. Sensation itself is precisely the product of these two
motions. The becoming is, and is not, just as the sensation is and is not. Becoming is of a twofold
kind: change of place and change of quality (ἀμφοτέρως τὰ πάντα κινεῖται, φερόμενόν τε καὶ ἀλλοιούμενον
[all things move in both ways, both moving in place and altering in quality]). In both respects the
sensed can be said to be and not to be. In saying: it is such, one can at once add: it is not such,
and not not-such (οὔτε ἄρα ὁρᾶν προσρητέον τι μᾶλλον ἢ μὴ ὁρᾶν, οὔτε τιν' ἄλλην αἴσθησιν μᾶλλον ἢ μή,
πάντων γε πάντως κινουμένων [one must therefore call it no more seeing than not-seeing, nor any other
sensation more than not, since all things are in every way in motion]). What is stated of the
sensation is thus, in the same moment as it is stated, untrue, and the alleged science consequently
not science (οὐδὲν ἄρα ἐπιστήμην μᾶλλον ἢ μὴ ἐπιστήμην ἀπεκρινάμεθα ἐρωτώμενοι ὅ τί ἐστιν ἐπιϛτήμη
[so, being asked what knowledge is, we answered nothing more knowledge than not-knowledge]).
% ---- printed p.128 (PDF 136) ----
By combating the sophistical Sensualism Plato approaches the Eleatics; but in the way in which the
Eleatic proposition — \emph{only the being is, the not-being is not} — was accented, the Platonic
Dialectic likewise discovers contradictions, and corrects the Eleatics in the dialogue that is named
the Sophist. What is one to judge of not-being? If the not-being is not at all, then it can neither be
thought nor stated; if the not-being is a semblance, then this semblance must nonetheless be. This is
precisely the difficulty in thinking not-being, that both he who denies it and he who defends it is
forced to contradict himself. If one grants that there exists a false opinion (ὄντος δέ γε ψεύδους
ἔστιν ἀπάτη [since, indeed, there is falsehood, there is deception]), then one at the same time
presupposes the representation of not-being; for only that opinion can be called false which either
declares the not-being to be being, or the being to be not-being. (κινδυνεύει τοιαύτην τινὰ πεπλέχθαι
συμπλοκὴν τὸ μὴ ὂν τῷ ὄντι [it seems that not-being has been interwoven with being in some such
interweaving]). If one holds the One for itself absolutely, it is impossible to say anything about it.
He who says that only the One is has already two names: the One and being; but if there is here a
difference between the One and its being, then here in this difference there is at the same time a
principle of plurality (τό τε δύο ὀνόματα ὁμολογεῖν εἶναι μηδὲν θέμενον πλὴν ἓν καταγέλαστόν που [to
admit two names when one has posited nothing but the One is surely ridiculous]). That becoming is a
semblance cannot be understood as though the being were to exclude all motion. If the essence of
things were wholly immovable, one could ascribe to it neither life nor reason; it could then not be
known either. The known and the knowing must stand in reciprocal relation to each other; the essence
of things that is to be known must consequently be affected by the knower, and thus come into motion
(μανθάνω, τόδε γε, ὡς τὸ γιγνώσκειν εἴπερ ἔσται, ποιεῖν τι, τὸ γιγνωσκόμενον ἀναγκαῖον αὖ ξυμβαίνει
πάσχειν [I understand this at least, that if knowing is to be a doing something, then the known must
correspondingly be affected]). That the being must be able to move is seen at the same time by a
consideration of the nature of speech, since speech consists precisely in this, that one in
% ---- printed p.129 (PDF 137) ----
propositions combines substantives (ὀνόματα [nouns]) with verbs (ῥήματα [verbs]), i.e. the abiding
with the changing (Δηλοῖ γὰρ ἤδη που τότε περὶ τῶν ὄντων ἢ γιγνομένων ἢ γεγονότων ἢ μελλόντων, καὶ
οὐκ ὀνομάζει μόνον, ἀλλά τι περαίνει, συμπλέκων τὰ ῥήματα τοῖς ὀνόμασι [for it then makes clear
something about things that are or are coming to be or have come to be or will be, and does not
merely name, but accomplishes something, weaving the verbs together with the nouns]). The two
concepts being and not-being are both difficult to grasp, but for wholly opposite reasons.
Not-being, with which the Sophist prefers to occupy himself, cannot be seen for sheer darkness, and
being, philosophy's object, is sought (διὰ λογισμῶν [through reasonings]) in those regions of light
where the eye is dazzled. What not-being is can, however, be known from the mutual relation of
concepts. Some concepts, namely, can be combined either wholly or in part with each other, while
others exclude each other. As for the general concepts being, rest, and motion, we see that rest
and motion do indeed mutually exclude each other, but that they nonetheless both, each for itself,
let themselves be united with being. The being of rest is the not-being of motion, and conversely;
what here, then, in relation to itself, is a being, is, in relation to its other, a not-being (τί
ποτ' αὖ νῦν οὕτως εἰρήκαμεν τό τε ταὐτὸν καὶ θάτερον [what, then, have we now thus called the same
and the other]): being and not-being unite and determine each other mutually.''
% [print anomaly, mirrored from transcription: this closing quote has no matching opening
%  (the Sophist paraphrase runs unquoted from p.128); reproduced as printed. The Danish/English
%  quote balance carries a net −1 from here onward because of this.]

\emph{b) The Ideas themselves and the sensible things.}

In that being, according to the concept, presupposes a not-being, the true original essence must
develop itself in a system of essences and essence-determinations. What was here with respect to
this point prepared in ``Theaetetus'' and ``the Sophist'' is now further carried through in
``Parmenides,'' ``Philebus,'' ``the Republic,'' and ``Timaeus.''

In ``Parmenides'' the investigation is introduced with the remark, set forth by Socrates, that the
being, if it is a plurality, must at once be a like and an unlike. But in that case the like would
have to be an unlike, and
% ---- printed p.130 (PDF 138) ----
the unlike a like; which is absurd (οὔτε γὰρ τὰ ἀνόμοια ὅμοια οὔτε τὰ ὅμοια ἀνόμοια οἷόν τε εἶναι
[for neither can the unlike be like nor the like unlike]). There must then be a likeness for
itself, an unlikeness for itself, a greatness for itself, a smallness for itself, and so on. The
Ideas must have a being in itself unconditioned (αὐτὰ καθ' αὑτά [themselves by themselves]), so
that the essence and properties of the finite things can be determined by a participation therein
(οὐ νομίζεις εἶναι αὐτὸ καθ' αὑτὸ εἶδός τι ὁμοιότητος καὶ τῷ τοιούτῳ αὖ ἄλλο τι ἐναντίον, ὃ ἔστιν
ἀνόμοιον; τούτοιν δὲ δυοῖν ὄντοιν καὶ ἐμὲ καὶ σὲ καὶ τὰ ἄλλα, ἃ δὴ πολλὰ καλοῦμεν, μεταλαμβάνειν
[do you not hold that there is a certain form of likeness itself by itself, and again another
contrary to such a one, which is unlikeness; and that, these being two, both I and you and the
other things, which we call many, partake of them]).

By participating in likeness the things must then become like, by participating in unlikeness
unlike. That the things, by participating in both concepts, become mutually like and unlike, is not
astonishing; but if anyone could prove that the like in itself were the unlike, and the unlike the
like, that would be astonishing. That all is one, insofar as it participates in the One, and at the
same time a manifold, insofar as it participates in plurality, is not to be wondered at; but if
anyone could prove that the One precisely in itself is a manifold, and the manifold one, then one
would have to wonder greatly.

But that which participates in the Idea must participate either in the whole Idea, or in a part of
it. That the whole Idea should be whole and undivided in each of the many objects participating
therein is unthinkable. One can spread a sheet over several people, but each single one then comes
into contact only with a part, not with the whole; each object thus gets only a part of the Idea,
not the Idea in its wholeness. But by being distributed among the objects the one Idea would fall
away from itself, and be parceled into plurality, which is likewise absurd. To wish, e.g., to
divide greatness itself, and to suppose that each of the many great things should be great by a
part of greatness which is smaller than the great, is foolish. But when the objects can get neither
a part of the Idea, nor the whole Idea (μήτε κατὰ μέρη
% ---- printed p.131 (PDF 139) ----
μήτε κατὰ ὅλα τῶν εἰδῶν μεταλαμβάνειν δυνάμενα [being able to partake of the forms neither part by
part nor as wholes]), then they can not participate in the Idea at all.

Perhaps the Idea is only a subjective concept, a thought in the soul? If the Ideas are only
subjective concepts (νοήματα [thoughts]), and thought must necessarily have its object, then one
must either assume that each object consists of thoughts and that all things think, or else that
there are indeed thoughts, but no thinking (ἀνάγκη, εἰ τἆλλα φῂς τῶν εἰδῶν μετέχειν, ἢ δοκεῖν σοι ἐκ
νοημάτων ἕκαστον εἶναι καὶ πάντα νοεῖν, ἢ νοήματα ὄντα ἀνόητα εἶναι [it is necessary, if you say the
other things partake of the forms, either that it seem to you each is composed of thoughts and all
things think, or that, being thoughts, they are unthinking]). Since now this assumption is absurd,
the Ideas must be nature's prototypes, and the objects that resemble them their images (τὰ μὲν εἴδη
ταῦτα ὥσπερ παραδείγματα ἑστάναι ἐν τῇ φύσει, τὰ δὲ ἄλλα τούτοις ἐοικέναι καὶ εἶναι ὁμοιώματα [these
forms stand as it were as paradigms in nature, and the other things resemble them and are
likenesses]).

The concept of participation is then hereby transformed into the concept of resembling something,
of taking on likeness with something (ἡ μέθεξις αὕτη τοῖς ἄλλοις γίγνεσθαι τῶν εἰδῶν οὐκ ἄλλη τις ἢ
εἰκασθῆναι αὐτοῖς [this participation of the other things in the forms is nothing other than being
made like them]). Yet neither does this turn satisfy; the like must indeed be like its like and
participate in the same Idea. But that by which the like becomes a like is precisely the Idea
itself. Should something now be like the Idea, or the Idea like something, then beside the Idea
there would have to appear another, likeness-effecting Idea; and, when this has likeness with
something, again another Idea, and so on to infinity. From whatever side the Idea and its other are
touched, the contradiction breaks forth; it is then of importance to investigate the Idea-essence,
namely with respect to the One (τὸ ἕν [the One]) and the manifold (τὰ πολλά [the many]).

α) If it is assumed that the One is (εἰ ἓν ἔστιν [if the One is]) in sharp separation from all
possible plurality, then this One cannot in the first place be a whole; for, if it were a whole, it
would have to have parts, and thereby come to consist in the plurality of the parts (Ἀμφοτέρως ἄρα
τὸ ἓν ἐκ μερῶν ἂν εἴη, ὅλον τε ὂν καὶ μέρη ἔχον. Ἀμφοτέρως ἂν ἄρα οὕτω τὸ ἓν πολλὰ εἴη, ἀλλ' οὐχ ἕν
[in both ways, then, the One would be of parts, both being a whole and having parts. In both ways,
then, the One would thus be many, and not one]). If the One has no parts, then it has
% ---- printed p.132 (PDF 140) ----
neither beginning (ἀρχή [beginning]) nor end (τελευτή [end]) nor middle (μέσον [middle]), for such
moments would then be its parts. But beginning and end are indeed the limits (πέρας [limit]) of
every thing; if the One has neither beginning nor end, it must be unlimited (ἄπειρον [unlimited]),
and without shape (ἄνευ σχήματος [without shape]). It can then be a partaker neither in the round
nor in the straight, can be surrounded neither by another nor by itself, can be neither in motion
nor in rest, can in its infinite self-contradiction tolerate no predicate whatever, can, in one
word, not be at all.

β) If it is nonetheless assumed that the One is, but in such a way that the difference between the
One itself and its being is brought out (ἄλλο τι σημαῖνον τὸ ἔστι τοῦ ἕν [``is'' signifying
something other than ``the One'']); then the being One is the whole, and the One itself and its
being the two parts of the whole. But in neither of these two parts does the One let itself be
torn loose from the being, or the being from the One. Each part then comes again to consist of two
parts, and each of these again of two; so that the being One becomes an unlimited plurality (ἄπειρον
ἂν τὸ πλῆθος οὕτω τὸ ἓν ὂν εἴη [the One that is would thus be unlimited in multitude]). With the
difference between the One and its being there now arises a series of different, mutually conflicting
predicates. The being One now becomes at once both limited and unlimited; with respect to likeness
it becomes both like and unlike itself and its other; with respect to magnitude it becomes both
greater and smaller than itself and its other, and, with respect to time, both older and younger
than itself and its other, and yet again neither older nor younger either than itself or the other.
One can now state everything possible of the One; which then can also signify that this One, by
being all in all, shows itself as the essence of the manifold, as the aggregate of all realities,
as the primal ground of existence.

γ) If, finally, it is assumed that the One is not (εἰ ἓν μὴ ἔστιν; [if the One is not]), yet in such
a way that this not-being of the One can be thought, stated, and determined by opposition to the
Not-One, then
% ---- printed p.133 (PDF 141) ----
the Not-One, the One's other, will in all its indeterminacy at the same time become more closely
determined. Although the One indeed cannot be, insofar as it is a not-being, it can nonetheless
participate in many properties (μετέχειν δὲ πολλῶν οὐδὲν κωλύει, ἀλλὰ καὶ ἀνάγκη [nothing prevents
it from partaking of many things, but it is even necessary]). By reason of the unlikeness with its
other, the One must namely be like itself; but, in that the other is posited as unlike the One, this
at the same time comes to participate in unlikeness, and thereby becomes at the same time unlike
itself. Indeed, the One must even, since it is not, participate in being. What is stated of the One
must indeed be true; now here it is stated: the One is a not-being. If the One therefore gave off
something of this its being, it would to a lesser degree be a not-being, and precisely thereby
become a being (εἰ γὰρ μὴ ἔσται μὴ ὄν, ἀλλά τι τοῦ εἶναι ἀνήσει πρὸς τὸ μὴ εἶναι, εὐθὺς ἔσται ὄν [for
if it will not be not-being, but relaxes something of being toward not-being, it will at once be
being]). Participating both in being and in not-being, the One at the same time participates in
change, motion, and so on. If one now reflects on the being other, this can not possibly be an other
over against the One, since this One is not (τοῦ μὲν γὰρ ἑνὸς οὐκ ἔσται ἄλλα, μὴ ὄντος γε [for there
will be no others of the One, since it is not]); it must then be an other over against another being
other (ἕτερον δὲ γὲ πού φαμεν τὸ ἕτερον εἶναι ἑτέρου, καὶ τὸ ἄλλο δὴ ἄλλο εἶναι ἄλλου [we say,
surely, that the other is other than an other, and the other is other than an other]): the otherness
expresses that the being, in infinite self-separation, relates to itself by incessantly falling away
from itself. The other is then not a being, but a something changing itself to infinity, a mass
unlimited in the quantity of its parts (ὁ ὄγκος αὐτῶν ἄπειρός ἐστι πλήθει [their mass is unlimited in
multitude]). In this mass each smallest part for itself may indeed seem to be a one; but it is as in
a dream-vision (ὥσπερ ὄναρ ἐν ὕπνῳ [as a dream in sleep]): suddenly, instead of the apparent one, the
manifold shows itself again, and, instead of a least, an infinitely great in relation to the
continued parceling.

With the development of otherness in ``Parmenides'' the element of sensuousness is already
intimated. This moment is
% ---- printed p.134 (PDF 142) ----
especially brought out in ``Philebus,'' which distinguishes between three kinds of being, namely the
limit (πέρας [limit]), the unlimited (ἄπειρον [unlimited]), and that composed of both (τὸ ἐξ ἀμφοῖν
τούτοιν ἕν τι ξυμμισγόμενον [something one, mixed together out of both of these]), to which is still
added a fourth principle, namely the cause of the mixture (αἰτία [cause]); likewise in ``Timaeus,''
where a distinction is made between that which truly is, which does not become, and that which always
becomes, which has no true being (τὸ ὂν ἀεί, γένεσιν δὲ οὐκ ἔχον, καὶ τὸ γιγνόμενον μὲν ἀεί, ὂν δὲ
οὐδέποτε [that which always is, but has no becoming, and that which always becomes, but never is]).
The union of both points to a third, as that which takes up all becoming and all forms in its bosom,
an invisible nurse of things (πάσης γενέσεως τιθήνη [the nurse of all becoming]), which in an
inconceivable way has become a partaker in the rational (ἀνόρατον εἶδός τι καὶ ἄμορφον, πανδεχές,
μεταλαμβάνον δὲ ἀπορώτατά πη τοῦ νοητοῦ [a certain invisible and formless kind, all-receptive,
partaking in some most perplexing way of the intelligible]), the mass (ἐκμαγεῖον [matrix]) out of
which all the sense-phenomena are formed. From matter, the not-being, which can neither be sensed nor
thought, but is grasped only by a kind of spurious inference (μετ' ἀναισθησίας ἁπτὸν λογισμῷ τινι
νόθῳ [graspable without sensation by a kind of bastard reasoning]), the things have their sensible
character; from the Ideas they have their fundamental form and reality; as the Ideas' imperfect
images in the sensible, the things are intermediate links between being and not-being.

\emph{c) The mutual relation of the Ideas.}

It is so far from the case that Plato, though he otherwise on the whole stands near the Eleatics,
remains standing at being's unity in the Idea, that he, as has also with right been remarked (Ritter),
seldom names the Idea in the singular (εἶδος or ἰδέα, in whose place he then rather uses other
expressions, such as: τὸ ὄντως ὄν [the truly-existent], τὸ κατὰ ταὐτὰ ἔχον [that which holds according
to the same], ἡ οὐσία [being], and so on), but almost always uses the plural, and speaks of the Ideas.

How the Eleatic One here with inner necessity separates itself into ideal plurality is evident from
the foregoing. In that each concept, so that the contradiction can be avoided, is to be held fast
for itself — likeness for itself, unlikeness for itself, motion for itself, rest for itself, and so
on — there arises an infinite
% ---- printed p.135 (PDF 143) ----
manifold of substantial units (ἑνάδες, μονάδες [henads, monads]) which, although each of them with
respect to the changing phenomena is a universal (τὸ ἐπὶ πᾶσι κοινόν [that which is common to all]),
and has a being independent of the sensible things (χωρίς ἐστι [is separate]), nonetheless mutually
exclude each other, and cannot possibly go up as moments in a single, absolute unity.

Since, however, it is not the special but the universal concepts that are ascribed highest validity,
it is at the same time seen that a system of superordinate and subordinate concepts must be able to
be formed. It is the dialectician's task that he, with respect to each object in question, is to seek
first the most universal concept, then divide it into its most nearly subordinate ones, and these in
turn into their subordinate ones, until he has exhausted the whole number of concepts that lie in the
middle between the One and the Unlimited. ``Since the whole of nature is akin, nothing prevents one
from starting from one thing and finding all, if only he is brave, and does not tire of seeking.''
(Meno). The graded series of the Ideas must then find its conclusion in a highest Idea, an Idea that,
without supporting itself on any higher presupposition, grounds itself and all (τὸ ἱκανόν [the
sufficient]). ``In the sphere of cognition the Idea of the Good (ἡ τοῦ ἀγαθοῦ ἰδέα [the Idea of the
Good]) is the highest, which is beheld only with difficulty, but which one cannot behold without at
the same time having to know that it is for all the cause of everything right and beautiful, that in
the visible world it has created the light and its lord the sun, and in the spiritual world, where
it alone rules, has brought forth truth and reason, and that therefore everyone ought to have his
gaze directed to it.'' (Republic).

The substantial discretion of the Ideas would, however, make a systematic derivation difficult to
such a degree that one even understands from this what Aristotle brings out, that the Platonic
doctrine of Ideas at last ended by passing over into a kind of Pythagorean number-doctrine. In the
dialogues there is indeed nothing that expressly shows this. A distinction is indeed here made
between pure and applied mathematics (ἀριθμοὶ
% ---- printed p.136 (PDF 144) ----
μαθηματικοί og αἰσθητοί [mathematical and sensible numbers]), and likewise (Republic) a distinction
between mathematical magnitudes and original number-concepts, but only in the same sense as when a
distinction is otherwise made between the thing and the Idea, so that the Ideas of numbers also belong
in the circle of the Ideas; on the contrary, there is here no thought that the Ideas on the whole
should be represented by numbers. Nonetheless the transition is (Zeller) explicable.

Since Plato, as especially ``Timaeus'' shows, intuits the Ideas in the corporeal world under the form
of mathematical determinations of relation, and comes further away from the dialectical, according as
he goes more closely into the concrete, there can be contained herein a hint of how he, when the
stricter Dialectic would not suffice, has been able to seize on the expedient of paralleling the Ideas
with the numbers. The Ideas are then grasped as primal numbers or ideal numbers (ἀριθμοὶ νοητοί,
εἰδητικοί, πρῶτοι [intelligible, ideal, first numbers]), and these primal numbers separated from the
mathematical ones (συμβλητοί [combinable]) in that they are principles, and (ἀσύμβλητοι [non-combinable])
cannot, like the ordinary numbers, be added together, or flow together into larger, arbitrary units.

\begin{center}{\bfseries \S\,28.}\end{center}

The weakness of the Platonic doctrine of Ideas becomes manifest from the critique Aristotle applies in
order to ground his standpoint, and carry the development further. That knowledge presupposes the
in-itself universally valid (τὸ ἐπίστασθαι πῶς ἔσται, εἰ μή τι ἔσται ἓν ἐπὶ πάντων [how will there be
knowing, if there is not some one over all]), that the essence, not the sensibly becoming, is the True,
about this Aristotle is agreed with Plato; but by the peculiar emphasis he lays on the single essence,
the existing, the real concretion, he diverges from Plato. As little as the thing, separated from its
concept, has essential form, so little has the concept, separated from the thing, real existence. With
this stands another divergence in close connection. Whereas Plato sets Dialectic as knowledge's highest
form, as expression of the proper science, Aristotle, on the contrary, prefers the
% ---- printed p.137 (PDF 145) ----
apodictic adducing of proof, and assigns the dialectical a rather subordinate place. In the so-called
aporias (ἀπορία [aporia]) the dogmatic decision is introduced by dialectical considerations (ἔστι δὲ
τοῖς εὐπορῆσαι βουλομένοις προὔργου τὸ διαπορῆσαι καλῶς [for those wishing to resolve difficulties it
is advantageous to work through the perplexities well]). Dialectic is hereby reduced to a reasoning for
and against, a supplement to the epagogic adducing of proof, an indication of the probable (διαλεκτικὸς
δὲ συλλογισμὸς ὁ ἐξ ἐνδόξων συλλογιζόμενος [a dialectical syllogism is one that reasons from reputable
opinions]). With such a reduction of the dialectical a deeper understanding of the Idea's essence is
doubtless impossible; but even under the presupposition of some misunderstanding, a critique which with
polemical acuteness lays bare the weak sides of the doctrine of Ideas, and thereby prepares a
penetrating cognition of the reality of the single essences, must be regarded as an indispensable
contribution to the further advancement of the matter (Μετὰ τὰ φυσικά, πρώτη φιλοσοφία [Metaphysics,
first philosophy]).

\emph{a) Aristotle's critique of the doctrine of Ideas.}

Although the Ideas, according to Plato's and the Platonists' assumption, are to be independent, for
themselves subsisting unities, they nonetheless have no peculiar content. The ideal two and the
empirical two have the same content. One speaks of man-itself, of horse-itself, of health-itself (αὐτὸ
γὰρ ἄνθρωπον φασιν εἶναι καὶ ἵππον καὶ ὑγίειαν [for they say there is man-itself and horse and health]);
but by setting ``itself'' before these head-words, nothing is won. The real horse has essentially the
same conceptual determinations as the horse-itself; to define the real man is the same as to define
man-itself (ποιοῦσιν αὐτοάνθρωπον καὶ αὐτοίππον, προστιθέντες τοῖς αἰσθητοῖς τὸ ῥῆμα τὸ αὐτό [they make
man-itself and horse-itself, adding to the sensibles the word ``itself'']). Aristotle compares the Ideas
thus made independent with the folk-gods; just as these are nothing other than mortal men whom one has
in fantasy made into eternal beings, so the Ideas too are nothing other than perishable things that one
has intuited in the form of the eternal and supersensible (οὔτε γὰρ ἐκεῖνοι, i.e. those who assume the
folk-gods,
% ---- printed p.138 (PDF 146) ----
οὐθὲν ἄλλο ἐποίουν, ἢ ἀνθρώπους ἀϊδίους, οὔθ' οὗτοι τα εἴδη ἀλλ' ἢ αἰσθητὰ ἀΐδια [neither did those make
anything other than eternal men, nor did these make the forms other than eternal sensibles]). The
``homonymy'' (ὁμωνυμία [homonymy]) occurring between the Ideas and the single existences in question
occasions the doubling.

When, then, the adherents of the doctrine of Ideas assume an Idea for each thing (εἴδη ἐστὶν ὁπόσα φύσει
[there are forms of as many things as exist by nature]), and thus double the number of substances, they
appear to Aristotle like men ``who supposed they could not reckon so well with few as with many numbers,
and who therefore first multiplied the numbers before they began with their operations of calculation.''
Such a doubling of substance is idle. Why set the same thing twice? To what end, besides the sensible two
and three, a two and three in the Idea? (ἄτοπος σώρευσις [absurd heaping-up]). The substance cannot
possibly be separated from that of which it is the substance (δόξειεν ἂν ἀδύνατον εἶναι χωρὶς τὴν οὐσίαν
καὶ οὗ ἡ οὐσία [it would seem impossible for the substance to be separate from that of which it is the
substance]). If the essence-determination is once and for all made into something for itself, then one
gets not merely an Idea for each substance, but even for each accident; indeed the Ideas, whose meaning
it is to be in the things, as in their other, become precisely thereby something accidental, for only
the accidental can be in an other.

That the Ideas should constitute the essence of the corporeal things, and yet be incorporeal, for
themselves existing substances, is a contradiction; partly, says Aristotle, Plato also speaks of an
``Ideas' matter,'' which straightforwardly conflicts with the assumption that they should be outside
space; partly in all natural objects matter and becoming belong to the objects' essence and concept, and
consequently this concept cannot be thought as a for-itself existing. One does not see how, and why,
matter participates in the Ideas. To make this understandable, one would have to assume, besides the
Ideas, a still higher principle as cause of this ``participation'' of things. The Ideas contain no
principle for becoming and motion; but without a principle for becoming and
% ---- printed p.139 (PDF 147) ----
motion it is not possible to explain to oneself what the things' ``participation'' in the Ideas is to
signify.

The assumption that the Ideas act as prototypes, and that the things in this respect have part in them,
is an empty representation (τὸ δὲ λέγειν παραδείγματα αὐτὰ εἶναι καὶ μετέχειν αὐτῶν τἆλλα, κενολογεῖν
ἐστι καὶ μεταφορὰς λέγειν ποιητικάς [to say that they are paradigms and that the other things partake of
them is to speak emptily and to utter poetic metaphors]). Were there, e.g., a prototype for Socrates,
then in this there would again have to be an infinite multitude of prototypes, a prototype for the
animal, a prototype for the two-footed, and so on. The distance between the universal man-itself and the
empirical single man is besides so great that here, as connecting link, an intermediate man is required
(ὁ τρίτος ἄνθρωπος [the third man]).

Insofar as the Idea is to be the unconditionedly single, it is not possible to define it; for the
absolutely single does not let itself be defined.

If the Ideas are grasped as primal numbers, how can they then be causes of the sensible things (πῶς
αἴτιοι ἔσονται; [how will they be causes])? The sensible things cannot possibly be given out for numbers,
for the eternal ideal numbers cannot indeed produce perishable things. If it is the number-relations on
which the emphasis is laid, then these, like the things themselves, presuppose a substratum whose
relations they determine (δῆλον ὅτι ἔστιν ἕν γέ τι ὂν εἰσὶ λόγοι [it is clear that they are ratios of
some one thing that exists]). Out of several numbers a new number is formed, but how can out of several
Ideas a new Idea be formed? Or is it perhaps not the numbers, but the units of which they consist, that
constitute the essence of things; but then the question again arises how these units relate to each
other. The triad (τριάς [triad]) is, e.g., man's number: what relation is there now between the units in
this triad and the units in any other triad whatever?

It will already from these few traits be clear that the Aristotelian critique, with all the warrant it
may otherwise have, does not properly set itself to lead us into Plato's original train of thought, but
rather — this is precisely the
% ---- printed p.140 (PDF 148) ----
uncritical thing in this critique — takes all that is had, loose and fixed, from Plato himself or
his adherents, as given, and then sets forth the objections, partly in gathered, partly in
scattered remarks. Only insofar as one looks to what Aristotle himself builds up, in tearing down
the doctrine of Ideas, does the critique get its proper significance.

\emph{b) Aristotelian principles.}

Substance, the original being (τὸ πρώτως ὂν καὶ οὐ τὶ ὂν ἀλλ' ὂν ἁπλῶς ἡ οὐσία ἂν εἴη [the
primarily existent — not a something-that-is, but simply being — would be substance]), is always a
definite this (τόδε τι [a this]), a peculiar subject (πᾶσα οὐσία δοκεῖ τόδε τι σημαίνειν [every
substance seems to signify a this]); substance in the proper sense (ἡ πρώτη οὐσία [the primary
substance]) is individual; in a subordinate sense the genus-concept is substance (ἡ δευτέρα οὐσία
[the secondary substance]); the other universal concepts are substance's attributes or accidents
(μὴ οὐσῶν οὖν τῶν πρώτων οὐσιῶν ἀδύνατον τῶν ἄλλων τι εἶναι [if the primary substances did not
exist, it would be impossible for any of the others to exist]).

In the sensible single essence, matter (ὕλη [matter]) and form (μορφή, εἶδος [shape, form]) are
distinguished. The way in which the reciprocal relation of matter and form is determined shows the
difference between Plato and Aristotle. In Plato matter, as the infinite otherness, is the
not-being, and form, as Idea, a for-itself-being; in Aristotle matter and form are two sides of the
same. In its abstraction from form, matter is the indeterminate, the predicateless, that is not to
be known (ἡ ὕλη ἄγνωστος καθ' αὑτήν [matter is unknowable in itself]). As the abiding that lies at
the ground of all becoming, matter takes on the most diverse forms, but is itself in reality
nothing. When there is talk of a determinate, a such (ἐκείνινον [a ``made-of-that'']), every matter
has a character; but universal matter is without any
% ---- printed p.141 (PDF 149) ----
peculiar character. Matter is on the one side opposed to form, on the other side one with form. As
one with form, matter can pass over into form, lie at the ground of form, and give it content; as
opposed to form, matter (στέρησις [privation]) hinders form from coming to complete actuality.

As the wood relates to the bench and the ore to the statue, so the first matter relates to all that
is determinate. The Eleatics denied becoming, the Megarians denied possibility (ὅταν ἐνεργῇ, μόνον
δύνασθαι. ὅταν δὲ μὴ ἐνεργῇ, μὴ δύνασθαι [that a thing can only when it is acting, and cannot when
it is not acting]): with his concept of matter Aristotle has explained becoming, and asserted
possibility. Out of the absolutely not-being nothing can become: out of nothing comes nothing; but
out of the according-to-actuality not-being, i.e. out of the according-to-possibility being, i.e.
out of matter, something can become. As form designates the actual (τὸ ἐνεργείᾳ ὄν [that which is
in actuality]), matter designates the possible (τὸ δυνάμει ὄν [that which is in potentiality]). By
the actual Aristotle then understands the being developed to totality, the essence insofar as it
has worked out all its determinations to existence; by the possible, the essence as the pure
disposition, the in-itself-undeveloped, which can, but therefore not necessarily must, become a
peculiarly determinate being (λέγομεν δὲ δυνάμει οἷον ἐν τῷ ξύλῳ Ἑρμῆν [we say ``in potentiality''
as, e.g., the Hermes in the wood]).

Just as Aristotle has determined the possible, he has also in corresponding manner determined the
contingent (συμβεβηκός, τὸ ἀπὸ τύχης [the accidental, that which is from chance]). The objective
ground of the contingent consists in this, that all that is not the completed actuality has in
itself the possibility both of being and of not-being, or, in other words, that matter, as the
indeterminate, makes possible opposite determinations (ἀόριστα μὲν οὖν τὰ αἴτια ἀνάγκη εἶναι, ἀφ'
ὧν γένοιτο τὸ ἀπὸ τύχης [the causes from which the chance-result comes must be indeterminate]).
What nature produces in such a way that it does not let itself be determined beforehand is a
contingent, which precisely for that reason also approaches the not-being (φαίνεται γὰρ τὸ
συμβεβηκὸς ἐγγύς τι τοῦ μὴ ὄντος [for the accidental appears to be something close to not-being]).

Matter's relation to form shows itself as a relation of potentiality to actuality (the
Scholastics' \textit{esse potentia} [being in potency] and \textit{esse actu} [being in act]).
Form, which makes the indeterminate, differenceless material into a differentiated, a this (τόδε τι
[a this]), an actual, is then the peculiar virtue, the completed activity, the soul in every thing.
``That which Aristotle
% ---- printed p.142 (PDF 150) ----
calls form may therefore not be confounded with what we call shape [Façon]. A severed hand, e.g.,
still has the outer figure of a hand; on the Aristotelian conception, on the contrary, it is a hand
only with respect to the material, not with respect to the form: a real hand, a hand with respect
to the form, is only that which can carry out a hand's peculiar work. Pure form is that which in
truth is without matter (τὸ τί ἦν εἶναι [the essence]), or the concept of the essence, the pure
concept. But such a pure form does not exist within the bound of determinate being. All that is a
this is, on the contrary, a composite of material and form, a σύνολον [composite whole]. The
individual is the harder to know, the more it bears the material in itself. But from what has here
been said it follows at the same time that the opposition between matter and form is a fluid one;
what in one respect is material is in another respect form: building-timber is, in relation to the
finished house, material, in relation to the unhewn wood, form; the soul is, in relation to the
body, form, in relation to reason — the form of form (εἶδος εἴδους [form of form]) — it is
material. From this standpoint the whole of all existence must present itself as a ladder whose
lowest rung is a first matter (πρώτη ὕλη [first matter]), which is not at all form, and whose
highest rung is the last form, which is not at all matter, but pure form (the absolute, divine
spirit). What is between these two extremes is in one respect matter, in another form, i.e. it is a
continual passing-over from the first to the last.'' (Schwegler).

Matter is that out of which the formed proceeds; but there must be something by which it proceeds:
τὸ κινῆσαν ἐκ δυνάμεως εἰς ἐνέργειαν [that which moves from potentiality to actuality]. The
transition from possibility to actuality is motion. The moving and the moved must be alike in
genus, unlike in species. It is a matter that already has the form for which another matter has the
disposition; the former then imparts to the latter its form. The moved and the movable (τὸ κινητόν
og τὸ κινητικόν [the moved and the moving])
% ---- printed p.143 (PDF 151) ----
cannot be before the motion; nor can either of them come to be before the motion, for such a
coming-to-be would precisely be a motion: every motion presupposes a preceding one, this again a
preceding one, and so on to infinity. The same is seen from the concept of time. Time cannot be
thought without a Now, but this is a middle between past and future. Time is without beginning; it
is a determination by motion (πάθος τι κινήσεως [an affection of motion]); an absolute arising and
passing-away, a coming-to-be out of nothing and a becoming-nothing, cannot be known (οὐ γίγνεται
οὔτε ἡ ὕλη οὔτε τὸ εἶδος, λέγω δὲ τὰ ἔσχατα [neither matter nor form comes to be, I mean the
ultimate ones]).

As the efficient cause (τὸ αἴτιον [the cause]) is the beginning of motion, the aim (τὸ τέλος, οὗ
ἕνεκα [the end, that for the sake of which]) is motion's terminus. Not only that which is undertaken
in conscious intention, but also that which is produced by natural causes, has its purpose; the
essence's actuality is motion's purpose, a purpose that lies in the motion itself. In relation to
the purpose, necessity (ἀνάγκη [necessity]), which like chance has its ground in matter, is more
closely determined. Matter is not, in Aristotle, as in Plato, the merely negative; matter is the
means without which form could not possibly realize the finite things, the means without which the
purpose could not possibly be attained. Here, then, with respect to the finite, there is a
development both according to presupposition and consequence (τὸ ἐξ ὑποθέσεως ἀναγκαῖον [that which
is conditionally necessary]) and according to the purpose (εἰσὶν ἄρα δύ' αἰτίαι αὗται, τό θ' οὗ
ἕνεκα καὶ τὸ ἐξ ἀνάγκης [there are then these two causes, that for the sake of which and that from
necessity]). Matter, form, the efficient cause, and the purpose are the four so-called Aristotelian
principles.

\emph{c) The Aristotelian entelechy.}

``The Aristotelian system (Schwegler) is a system of becoming; in this system the principle of
Heraclitus returns in richer and more mediated conception, as the Eleatics' principle in the
Platonic. Aristotle has made a considerable step toward the overcoming of the Platonic dualism. If
% ---- printed p.144 (PDF 152) ----
matter is form's possibility, becoming reason, then the opposition between Idea and the world of
sight is at least in principle potentially overcome, insofar as it is one being that presents itself
in both, in form and material, only at different stages of development.''

``The relation of the potential to the actual Aristotle makes visible by the relation of the
unworked to the worked, the master-builder to the one occupied with the building, the sleeping to
the awakened. The seed is potentially a tree, the full-grown tree is it actually; a potential
philosopher is also he who at this moment does not philosophize; a potential victor is the better
general even before the battle; potentially space is divisible to infinity; potentially, according
to possibility, is on the whole all that has a principle of motion, development, change,
being-other in itself.''

``Actuality or entelechy, on the contrary, designates the perfect action, the attained goal, the
completed actuality (the full-grown tree, e.g., is the seed-corn's entelechy), the activity in which
the action and its completion coincide in one, e.g. seeing, thinking: he sees, and he has seen, he
thinks, and he has thought — is one and the same, whereas in the activities that are bound up with a
becoming, e.g. learning, walking, healing, both — action and completion — fall apart'' (κίνησις
ἐντελής as opposed to κίνησις ἀτελής [complete motion as opposed to incomplete motion]).

The perfect actuality, the absolute entelechy, is then the concept in which the Aristotelian
doctrine of essence comes to conclusion. This concept's reality is in the Metaphysics illuminated
from different sides.

α) The first mover. The finite motions have beginning and end, but motion itself is eternal like
time. If there is an eternal motion, there must also be a moving first. The first mover cannot be
according to possibility; but, if its essence is to be the completed actuality, we must convince
ourselves
% ---- printed p.145 (PDF 153) ----
that the actual is always the presupposition of the possible. In the relative sense the actual does
indeed proceed from the possible: the plant as actuality proceeds from the seed-corn, as from its
possibility; but the possibility is nonetheless not the first: the seed-corn, the possible plant, is
itself a product of a real plant (against the Pythagoreans and Speusippus: τὸ σπέρμα ἐξ ἑτέρων ἐστὶ
προτέρων τελείων, καὶ τὸ πρῶτον οὐ σπέρμα ἐστίν, ἀλλὰ τὸ τέλειον [the seed is from other prior
complete things, and the first is not seed but the complete]). On actuality and activity Aristotle
lays — it is precisely one of his deepest thoughts — an unconditioned emphasis; the actual is the
original (διὸ καὶ τοὔνομα ἐνέργεια λέγεται κατὰ τὸ ἔργον, καὶ συντείνει πρὸς τὴν ἐντελέχειαν
[wherefore the name ``actuality'' (energeia) is said with reference to ``work'' (ergon), and tends
toward entelechy]).

That actuality in essence and substantiality is the first is, with respect to time and motion,
expressed thus, that the one actual overreachingly presupposes the other, until one comes to the
first mover (τοῦ χρόνου ἀεὶ προλαμβάνει ἐνέργεια ἑτέρα πρὸ ἑτέρας ἕως τῆς τοῦ ἀεὶ κινοῦντος πρώτως
[in time one actuality always precedes another, up to that of the primary ever-mover]). The more
something is in actuality, the more energy it has; the more energy something has, the less material
it is, and the less it is according to possibility. Therefore one need not, like some physicists,
fear that the heaven with the sun and the fixed stars should ever grow tired of moving.

If the first mover were according to possibility, it would also have to have been possible that no
motion had come about (ἐνδέχεται γὰρ τὸ δύναμιν ἔχον μὴ ἐνεργεῖν [for that which has potentiality
may not be acting]). There was then not either eternal chaos or eternal night (ὥστε οὐκ ἦν ἄπειρον
χρόνον χάος ἢ νύξ [so there was not chaos or night for infinite time]) when the motion began; but
the same has always been, either in constant alternation, or in another way, as truly as actuality
is prior to possibility (εἴπερ πρότερον ἐνέργεια δυνάμεως [if indeed actuality is prior to
potentiality]). If now the same is always to act in constant alternation, then there must be
something that always acts in the same way (δεῖ τι ἀεὶ μένειν ὡσαύτως ἐνεργοῦν [something must
always remain acting in the same way]); if there is to be arising and passing-away, then an other
must incessantly act
% ---- printed p.146 (PDF 154) ----
otherwise and again otherwise (ἄλλο δεῖ εἶναι ἀεὶ ἐνεργοῦν ἄλλως καὶ ἄλλως [another must always be
acting otherwise and otherwise]). The first mover is thus on the one side, in relation to itself,
always like itself, on the other side, in relation to another, always acting otherwise: τί οὖν
ἄλλας δεῖ ζήτειν ἀρχάς; [why then must one seek other principles?]

The first heaven is eternal (δῆλον, ὥστε ἀΐδιος ἂν εἴη ὁ πρῶτος οὐρανός [clearly, then, the first
heaven would be eternal]); the first heaven is only one. Were there several heavens, as there are
several men, then their principles would indeed be identical in concept, but different in number.
What is thus in number a plurality must have matter; but the first essence (τὸ τί ἦν εἶναι τὸ
πρῶτον [the primary essence]) can have no matter. The first mover, which is itself unmoved, since it
is always like itself, is then one both in concept and in number: ἓν ἄρα καὶ λόγῳ καὶ ἀριθμῷ τὸ
πρῶτον κινοῦν ἀκίνητον ὄν [the first mover, being unmoved, is thus one both in account and in
number].

β) The Good. Every possibility is determined by contradiction; for what is not possible cannot be,
but the possible need not in actuality be. The possible can then be and not be, and this holds of
all that is finite. What is according to possibility can be changed, and in the change become the
opposite of what it was. The healthy can become sick, the sick healthy; the better can become
worse, this again better. The eternal, on the contrary, cannot be according to possibility (οὐθὲν
ἄρα τῶν ἀφθάρτων ἁπλῶς δυνάμει ἐστὶν ὂν ἁπλῶς [nothing imperishable is simply in potentiality
simply]); the eternal is with necessity what it is: the eternal is the unchangeable (ἀπαθὲς καὶ
ἀναλλοίωτον [impassive and unalterable]). In this unchangeableness the entelechy is the Good. The
changeable can fall away from its concept, miss its goal, and therein show itself as the bad, the
evil; but the eternal, the unchangeable, has the terminus in itself, cannot fall away from its
concept, and is precisely thereby the perfect Good. There is then in the principles and in the
eternal no evil, no fault, and no corruptibility; corruptibility belongs to the evil (οὐκ ἄρα οὐδ'
ἐν τοῖς ἐξ ἀρχῆς καὶ τοῖς ἀϊδίοις οὐθέν ἐστιν οὔτε κακὸν οὔτε ἁμάρτημα
% ---- printed p.147 (PDF 155) ----
οὔτε διεφθαρμένον. καὶ γὰρ ἡ διαφθορὰ τῶν κακῶν ἐστιν [so neither in the first principles nor in the
eternal things is there anything evil or faulty or corrupted; for corruption belongs to evil
things]). As the eternal, perfect Good, the entelechy is that which moves in that it is desired; for
all in nature is laid out toward a highest common goal (πρὸς μὲν γὰρ ἓν ἅπαντα συντέτακται [for all
things are ordered together toward one]). As the in-itself perfect Good, the entelechy can move all
without itself being moved; for that which is desired and thought moves without itself being moved.
What shows itself beautiful is desired, what is beautiful is willed. We desire the beautiful because
it shows itself, but it does not show itself because we desire it. The purpose moves because it is
loved, and what is moved by the purpose then moves in turn the rest. The divine then consists in
eternal, infinite energy, thus in eternal living, for life is energy (ὥστε ζωὴ καὶ αἰὼν συνεχὴς καὶ
ἀΐδιος ὑπάρχει τῷ θεῷ [so that life and continuous, eternal duration belong to God]). This life is
blessed; such a joy is granted us only in short moments; but all joy (ἡδονή [pleasure]) is an effect
of this principle.

γ) Reason. Thinking and the thought are in reason one (ὥστε ταὐτὸν νοῦς καὶ νοητόν. τὸ γὰρ δεκτικὸν
τοῦ νοητοῦ καὶ τῆς οὐσίας νοῦς [so that reason and the object of thought are the same; for that
which is receptive of the intelligible and of substance is reason]); but the thought is, as the
thought-content, higher than the thinking: he who thinks the worst is in reality just as much
thinking as he who thinks the best. The true, eternal thinking must then be determined by the true,
eternal content. In the finite, material existence thinking cannot become identical with itself and
its content, but where there is no matter, thinking is one with the thought (οὐχ ἑτέρου οὖν ὄντος
τοῦ νοουμένου καὶ τοῦ νοῦ, ὅσα μὴ ὕλην ἔχει τὸ αὐτὸ ἔσται καὶ ἡ νόησις τοῦ νοουμένου μία [the
thought and its object not being different, in things that have no matter it will be the same, and
the thinking of the object one]). If the thought is composite, then thinking changes itself by
passing over from one part to another; but what has no matter is indivisible like the human
understanding, or as when a composite has reached its completion. The purpose is then not in this or
that particular, but in a whole the Good, as an essence different from the piecework (οὐ γὰρ ἔχει τὸ
εὖ ἐν τῳδὶ ἢ τῳδί, ἀλλ' ἐν ὅλῳ τινὶ τὸ ἄριστον, ὂν ἄλλο τι [for the good is not in this or that
particular, but the best is in a certain whole, being something else]). So it stands
% ---- printed p.148 (PDF 156) ----
with the thought that thinks itself in all eternity. It is the divine understanding that eternally
thinks itself, and this thinking is the thinking of thinking (αὐτὸν ἄρα νοεῖ εἴπερ ἐστὶ τὸ
κράτιστον, καὶ ἔστιν ἡ νόησις νοήσεως νόησις [it thinks itself, then, if it is the most excellent,
and the thinking is a thinking of thinking]).

By impressing energy, pointing out the principle of motion, developing the immanent teleology, and
analyzing the essence of the Godhead, Aristotle has continued Plato's work, and carried the doctrine
of Ideas further.

\begin{center}{\bfseries \S\,29.}\end{center}

When the pulse stops, life ceases; when the development stops, the Idea vanishes. If we now regard
the given traits in the Aristotelian entelechy as given, existing, stationary, the physiognomy is
dead, and the worm of contradiction develops itself in the dead.
% [print: the word rendered ``dead'' (dødt) is set with an ambiguous d/b glyph (``døbt'');
%  context (``in the dead'') makes ``dead'' certain.]

``One does not see (Schwegler) how something can be a moving cause and yet itself unmoved, cause of
all becoming, i.e. of all passing-away and arising, and yet an energy that remains like itself, a
principle of motion without possibility (potentiality): for the moving must after all stand in a
relation of suffering and acting to the moved.''

That here in each of these points a contradiction shows itself cannot be denied; but it depends on
whether such a contradiction is applied critically or dialectically. If one sets up, according to
the logical principle of exclusion (\textit{principium exclusi medii inter duo contradictoria} [the
principle of the excluded middle between two contradictories]), the proposition: ``The first mover
must necessarily be either moved or unmoved,'' one can easily prove that the Aristotelian philosophy
culminates in an unsolved and insoluble contradiction.

Outside the Idea the contradiction is critical; in the development of the Idea it is dialectical. To
grasp the relation of the unmoved and the moved in such a way that they immediately exclude each
other is to grasp it without Idea: in the Idea the unmoved is itself moved. The Idea is the inner
unity of the otherwise mutually excluding determinations, the identity of the otherwise mutually
contradicting moments.
% ---- printed p.149 (PDF 157) ----
This conception of the contradiction is in express conflict with the Aristotelian logic, but in
intimate understanding with the Aristotelian entelechy. When Aristotle determines his πρῶτον κινοῦν
[first mover] as an ἀκίνητον [unmoved], it is everywhere clear from the connection that by this
ἀκίνητον he cannot possibly think an idle resting that should exclude self-motion.

In the imperfect energy rest and motion fall apart: the resting does not move; the self-moving does
not rest; the resting must be set in motion, the moved go to rest. In the perfect energy the motion
coincides with the rest: here is rest in the motion, motion in the rest. By ἀκίνητον, then, according
to the meaning, only the finite in the motion (e.g. local motion) is excluded, not the inner
self-motion. But that the meaning has not found its conceptual expression in the presentation proves
at the same time that the form of presentation is inadequate: the dialectical is lacking.

If a glance at the Idea, immediate intuition of the Idea, could satisfy philosophy's demand, the
demand would already often and in many different ways have been satisfied. Not only Plato and
Aristotle, but also Jacob Böhme would then have satisfied the demand. But if the task is to give the
Idea its pure, clear, articulated conceptual form, we must, after having heard Plato and Aristotle,
especially turn to Hegel as the one who, by his glance at the contradiction and his dialectical
method, occupies the highest place.

\emph{a) The Dialectic of contradiction.}

That black is white, that good is evil, that round is square, are contradictions that are just as
little tolerated in the doctrine of Ideas as in existence; but that with respect to the Idea itself
there are contradictions of another significance can among other things be seen from the Platonic
dialogues. In ``the Sophist'' and in ``Parmenides''
% ---- printed p.150 (PDF 158) ----
it is indeed expressly set out to keep the contradiction out, and yet the contradiction breaks
through, and that with so much the greater force, the more the acuteness exerts itself to withstand
it.

``The not-being is not'' — that is the proposition that is unconditionally to be held fast, and then
it turns out that this not-being nonetheless has a being. The One is not a manifold, and then it
turns out nonetheless that precisely this One makes itself a manifold. In this transition is the
Idea. The Idea is not the contradiction; but, in that the thoroughgoing contradiction dissolves all
finite difference into the differenceless, the Idea is itself the solution of the contradiction. The
manifold is, in its finite and precisely thereby infinite manifoldness, not the One, as little as
the One in and for itself either is or can be the manifold; that the One by its own contradiction
nonetheless makes itself manifold means that it cannot hold itself in the idealess. The idealess
One, the Eleatics' τὸ ἕν [the One], is in its solitude barren, unfruitful, and cannot endure the
manifold. When now the Idea approaches, this in-itself-hardened One is attacked by the contradiction,
compelled to give up its enclosedness, and pass over into the manifold. Seen under the critical angle
of the contradiction, this transition is the One's downfall: the One now is not, for now the manifold
is; but this critical is lifted by the fact that the contradiction serves the Idea. The One then does
not perish, no, it goes up into the Idea, sucks in the Idea's force, produces the manifold, and shows
itself now as the manifold's One, as the Idea's own true, content-rich essence. The more thinking is
sharpened, the more distinctly it is seen that the carried-through contradiction is precisely the
contradiction's solution — in the Idea.

From this point of view the relation between Plato and Hegel can be more closely illuminated. In the
Platonic Dialectic only the critical side of the contradiction is yet brought to consciousness.
% ---- printed p.151 (PDF 159) ----
Concept-analyses distinguish themselves precisely by the fine, microscopic separation of the
moments; the effort is to avoid the contradiction by help of the distinction: \textit{qui bene
distingvit, bene docet} [he who distinguishes well teaches well]. The subtle, apt distinction now
doubtless brings with it a transition from the one conceptual determination to the other; but this
transition is only something that happens, something wonderful that befalls the thinker; it is the
contradiction that is after him, and drives him from the one position to the other, without the
proper essence of the transition being seen through. Only in Hegel is it brought to consciousness
what the transition signifies, and in what relation it stands to the contradiction. If we call the
Platonic Dialectic a dialectic of distinction, we must call the Hegelian a dialectic of identity. In
the dialectic of distinction the antinomic rules; the dialectic of identity is laid out to master
the antinomic.

``\textit{Die wahre und positive Bedeutung der Antinomien besteht nur überhaupt darin, daß alles
Wirkliche entgegengesetzte Bestimmungen in sich enthält.} [The true and positive meaning of the
antinomies consists only, in general, in this, that everything actual contains opposed determinations
in itself.]''

``\textit{Was das Vorkommen der Dialektik in der geistigen Welt, und näher auf dem Gebiet des
Rechtlichen und Sittlichen anbetrifft, so braucht hier nur daran erinnert zu werden, wie, allgemeiner
Erfahrung zufolge, das Äußerste eines Zustandes oder eines Thuns in sein Entgegengesetztes umschlagen
pflegt, welche Dialektik dann auch vielfältig in Sprüchwörtern ihre Anerkennung findet. So heißt es
z. B.} summum jus summa injuria, \textit{womit ausgesprochen ist, daß das abstrakte Recht auf seine
Spitze getrieben in Unrecht umschlägt. Eben so ist es bekannt, wie im Politischen die Extreme der
Anarchie und des Despotismus einander gegenseitig herbeiführen pflegen.} [As for the occurrence of
Dialectic in the spiritual world, and more nearly in the domain of the juridical and ethical, one need
here only be reminded how, according to general experience, the extreme of a condition or of a doing
is wont to turn round into its opposite, which Dialectic then also finds its recognition manifoldly in
proverbs. So it is said, e.g., summum jus summa injuria (the highest right is the highest wrong),
whereby it is expressed that abstract right, driven to its peak, turns round into wrong. Likewise it
is known how in the political the extremes of anarchy and of despotism are wont mutually to bring
about each other.]''

``\textit{Das Bewußtseyn der Dialektik im Gebiet des Sittlichen in seiner individuellen Gestalt finden
wir in jenen allbekannten Sprüchwörtern: Hochmuth kommt vor dem Fall — Allzuscharf macht schartig
u. s. w. — Auch die Empfindung, die leibliche sowohl als die geistige, hat ihre Dialektik.}
% ---- printed p.152 (PDF 160) ----
\textit{Es ist bekannt, wie die Extreme des Schmerzes und der Freude in einander übergehen; das von
Freude erfüllte Herz erleichtert sich in Thränen, und die innigste Wehmuth pflegt unter Umständen
sich durch Lächeln anzukündigen} [The consciousness of Dialectic in the domain of the ethical, in its
individual shape, we find in those well-known proverbs: pride goes before a fall — too sharp makes
notches, and so on. — Feeling too, the bodily as well as the spiritual, has its Dialectic. It is known
how the extremes of pain and of joy pass over into each other; the heart filled with joy relieves
itself in tears, and the most intimate melancholy is wont under circumstances to announce itself
through a smile.]'' (Hegel).

In Plato and Aristotle the contradiction is a being that occasions the dialectical: either there is in
this a contradiction, or there is none; either the contradiction is here or there. In Hegel the
contradiction itself is the dialectical: it is, and yet is not; it is, for the finite is; it is not,
for the Idea is. Here is in this finite a \textit{summa injuria}; it is finitude's ethical
contradiction, and this contradiction is man's torment. But \textit{summa injuria} says precisely that
here is a \textit{summum jus}; for only in that righteousness is, can unrighteousness be. So it is then
not a torment to suffer wrong for righteousness' sake, but a blessedness, for righteousness is in the
Idea, and the Idea is the blessed; it is thus a torment to suffer, and yet blessed to suffer: the
contradiction is, for the finite is; the contradiction is not, for the Idea is.

\emph{b) The infinite negativity.}

One hears not seldom the opinion expressed that negative knowledge does not satisfy, that philosophy
ought to be positive. It now depends on how one understands the positive. If the positive is grasped
so immediately that it has the negative outside itself, it is grasped in finite form, and is then in
this its finitude fallen prey to the contradiction, and thereby to the negative. A philosophy that is
dogmatically positive by keeping the negative out is without Idea, and the idealess positivism
precisely by the Idea destined for dissolution into the negative. Therefore Dogmatism is always
attacked by Scepticism. Dogmatism is idealess, because it lacks
% [d/b glyph flag in the Danish (Rettelser): print ``forbi'' for ``fordi'' (= because), twice here]
the negative; Scepticism is idealess, because it lacks the positive. The positive is in the Idea
infinitely negative, and this negative precisely therefore the true positive. The
% ---- printed p.153 (PDF 161) ----
objection that a philosophy is not positive enough is then, in the concept, to be understood so, that
% [d/b glyph flag in the Danish: print ``venbing'' for ``vending''; the word is ``Indvending'' = objection]
it is not negative enough. Scepticism is, in the same degree as it has the negative, nearer to
reaching the true positive than Dogmatism.

``\textit{Der Skepticismus darf nicht blos als eine Zweifelslehre betrachtet werden, vielmehr ist
derselbe seiner Sache, d. h. der Nichtigkeit alles Endlichen schlechthin, gewiß. Wer nur zweifelt, der
steht noch in der Hoffnung, daß sein Zweifel gelöst werden könne, und daß das eine oder das andre
Bestimmte, wozwischen er hin und herschwankt, sich als ein Festes und Wahrhaftes ergeben werde.
Dahingegen ist der eigentliche Skepticismus die vollkommne Verzweiflung an allem Festen des
Verstandes, und die sich daraus ergebende Gesinnung ist die der Unerschütterlichkeit und des
Insichberuhens. Dieß ist der hohe antike Skepticismus, wie wir ihn namentlich beim Sextus Empirikus
dargestellt finden, und wie derselbe als Komplement zu den dogmatischen Systemen der Stoiker und
Epikuräer, in der spätern Römerzeit seine Ausbildung erhalten hat} [Scepticism may not be regarded
merely as a doctrine of doubt; rather it is certain of its cause, i.e. of the nullity of everything
finite absolutely. He who only doubts still stands in the hope that his doubt can be resolved, and that
the one or the other determinate, between which he sways back and forth, will turn out to be something
fixed and true. On the contrary, the proper Scepticism is the complete despair of everything fixed of
the understanding, and the disposition resulting therefrom is that of imperturbability and of
resting-in-itself. This is the high ancient Scepticism, as we find it presented namely in Sextus
Empiricus, and as it, as complement to the dogmatic systems of the Stoics and Epicureans, received its
development in the later Roman period.]'' (Hegel).

To carry through the negation is to negate it; in the negation of the negation the positive returns.
The ancient Scepticism has the infinite negation, but in a finite way. The art consists in shaking all
assumptions and overthrowing all positive propositions; one goes from denial to denial without halting
at any affirmation whatever. But since each denial for itself is only a denial of this or that finite,
and thereby itself finite, the continuation of the partial negations is an endless finitude that lacks
essentiality. The inner infinity of the negation, on the contrary, is essentiality's own immanent
reflection, is as the negation of the negation the positive.

``\textit{Denn das Wesen ist das Selbstständige, das ist als durch seine Negation, welche es selbst
ist, sich mit sich vermittelnd; es ist also die identische Einheit der absoluten Negativität und der
Unmittelbarkeit. — Die Negativität ist
% ---- printed p.154 (PDF 162) ----
die Negativität an sich; sie ist ihre Beziehung auf sich, so ist sie an sich Unmittelbarkeit; aber sie
ist negative Beziehung auf sich, abstoßendes Negiren ihrer selbst, so ist die an sich seyende
Unmittelbarkeit das Negative oder Bestimmte gegen sie. Aber diese Bestimmtheit ist selbst die absolute
Negativität und dieß Bestimmen, das unmittelbar als Bestimmen das Aufheben seiner selbst, Rückkehr in
sich ist} [For the essence is the independent, which is as mediating itself with itself through its
negation, which it itself is; it is thus the identical unity of absolute negativity and of immediacy. —
Negativity is negativity in itself; it is its relation to itself, so it is in itself immediacy; but it
is negative relation to itself, repelling negating of itself, so the in-itself-being immediacy is the
negative or determinate against it. But this determinacy is itself the absolute negativity, and this
determining, which immediately as determining is the sublating of itself, is return into itself.]''
(Hegel).

And what then is the whole Idea? No thing! And the properties of things? No thing! And the thing in
itself? No thing! By establishing each of these negations we are led into Scepticism, and thereby away
from — i.e. on the way to — the Idea.

When it is here said of the Idea that it is no thing, the meaning is this, that the thing is, and is in
all immediacy the positive. We get by this denial of the Idea a world of things; but the ideal has in
itself no validity: only the sensible is; the non-sensible is not. Hereby is characterized the modern
Scepticism, ``\textit{welcher bloß darin besteht, die Wahrheit und Gewißheit des Uebersinnlichen zu
leugnen, und dagegen das Sinnliche und in der unmittelbaren Empfindung Vorhandene als dasjenige zu
bezeichnen, woran wir uns zu halten haben.} [which consists merely in denying the truth and certainty
of the supersensible, and on the contrary in designating the sensible and that present in immediate
feeling as that to which we have to hold.]''

When it is here said of the thing's properties that they are not things, the meaning is this, that the
property is precisely the thing's negation, in that it is nonetheless a positive determination of the
thing. To be property (predicate) is, as Aristotle too impresses, not the same as to be thing
(subject). Man is white, snow is white: the white is, as property, not a thing. The white has in its
pure whiteness no, absolutely no real being; the whiteness is only insofar as it is in the white
substance. As it now goes with the whiteness, so it goes with the hardness, softness, sourness,
sweetness, and, in short, with all possible properties. The many different properties seem at first
glance to be sheer positive determinations, but dissolve, when
% ---- printed p.155 (PDF 163) ----
thinking tests them, into pure negations. But, in that the properties' reality in this way vanishes,
the reality of things must also vanish; the thing in itself is a phantom: the thing is only what it is,
in and by its properties.

To let the dialecticizing thought transform all properties into figments, all positive determinations
into partial negations, is the distinguishing mark of the ancient Scepticism, which then also lets the
reality of things be consumed in the same negations.

In that the thing is in this way negated, and becomes no thing, the Idea comes forth as Idea. The Idea
is not the thing; but the thing dissolved into negations is restituted in and by the Idea. That the
idealess immediacy bears the negativity in itself is a proof that it is subject to the Idea. To see the
properties' immediate plurality, not as an aggregate of positively different ones, but as negations of
the thing's own essential unity, is to see the properties in the Idea; to see the thing's immediate
unity not as substratum, but as negation of the properties' plurality and positive difference, is to
see the thing in the Idea. The immediately positive is in its essence negative; only in the negation of
the negation is the Idea, the true positive essence.

\emph{c) The speculative form.}

The positive in the Idea is, with another expression, at the same time the speculative. As the positive
proceeds from the negation, the speculative proceeds from the reflection.

How Hegel grasps the relation between immediacy, reflection, and speculation becomes manifest
especially from his ``Wissenschaft der Logik.'' In the form of immediacy being is one, and not-being an
other; likewise essence is one and semblance an other. In this immediate separation it looks as though
one could hold fast to the one of the juxtaposed concepts, and give up the other. If now this attempt
is made in
% ---- printed p.156 (PDF 164) ----
earnest, the reflection begins. In the reflection it shows itself that the separated concepts determine
each other mutually. Being has as positive being a reflex of its negation, not-being; conversely there
is here in not-being, as the expression itself shows, a reflex of being. It is not merely we who make
our reflections about being and not-being, no, the concepts mirror each other mutually: the reflection
is objective. Would one, in order to hold to the pure essence, think all semblance away; then the pure
essence in its infinite emptiness would itself be semblance: by denying the semblance, we would
precisely get the semblance confirmed. Would one conversely, instead of knowing an essence in the
semblance, deny the essence, and hold fast to the semblance alone, then we would thereby get semblance
of semblance, a semblance shining through itself, i.e. the essence: by negating the essence we get the
essence confirmed. The unity of the mutually negating and thereupon affirming concepts is seen in
speculation. Thus becoming is a speculative unity of being and not-being, as the phenomenon is a
speculative unity of essence and semblance; the speculative unity is dialectical.

``\textit{Ein spekulativer Inhalt kann deshalb auch nicht in einem einseitigen Satz ausgesprochen
werden. Sagen wir z. B. das Absolute sey die Einheit des Subjektiven und des Objektiven, so ist dieß
zwar richtig, jedoch insofern einseitig, als hier nur} \emph{die Einheit} \textit{ausgesprochen und
auf diese der Akcent gelegt wird, während doch in der That das Subjektive und das Objektive nicht nur
identisch, sondern auch unterschieden sind.} [A speculative content can therefore also not be expressed
in a one-sided proposition. If we say, e.g., the Absolute is the unity of the subjective and the
objective, then this is indeed correct, yet one-sided insofar as here only \emph{the unity} is
expressed and the accent laid on this, while in fact the subjective and the objective are not only
identical but also distinguished.]''

By setting speculation in intimate, inseparable connection with the dialectical reflection, Hegel has
secured for the speculative a scientific character; whereas Schelling, who remained standing at ``the
intellectual intuition,'' let the speculative go together with the mystical, and approach the
apocalyptic, whereby all manner of fantasticalities could gain entrance.
% ---- printed p.157 (PDF 165) ----
The Idea's speculative form is in Hegel in inseparable unity with the dialectical method. Although
immediacy, reflection, and speculation penetrate each other at all points of the development, the great
tripartition — being, essence, and concept — is nonetheless laid out in such a way that being on the
whole corresponds to immediacy, essence to reflection, and the concept to speculation. The Idea, which
is especially dealt with in the last part of the doctrine of the concept, is therefore not a special
part of the Logic's content; the whole Logic is a doctrine of Ideas: the Idea is the concept's absolute
content, as the concept is the Idea's absolute form.

``\textit{Die Idee ist die Wahrheit; denn die Wahrheit ist dieß, daß die Objektivität dem Begriffe
entspricht, — nicht daß äußerliche Dinge meinen Vorstellungen entsprechen; dieß sind nur} \emph{richtige}
\textit{Vorstellungen, die Ich Dieser habe. In der Idee handelt es sich nicht um Diesen, noch um
Vorstellungen, noch um äußerliche Dinge. — Aber auch} \emph{alles} \textit{Wirkliche, in sofern es ein
Wahres ist, ist die Idee, und hat seine Wahrheit allein durch und kraft der Idee. Das einzelne Seyn ist
irgend eine Seite der Idee, für dieses bedarf es daher noch anderer Wirklichkeiten, die gleichfalls als
besonders für sich bestehende erscheinen; in ihnen zusammen und in ihrer Beziehung ist allein der
Begriff realisirt.} [The Idea is the truth; for the truth is this, that objectivity corresponds to the
concept — not that external things correspond to my representations; these are only \emph{correct}
representations that I, this one, have. In the Idea it is not a matter of this one, nor of
representations, nor of external things. — But also \emph{everything} actual, insofar as it is a true,
is the Idea, and has its truth solely through and by force of the Idea. The single being is some side
of the Idea; for this it therefore needs still other actualities, which likewise appear as particularly
for-themselves subsisting; in them together and in their relation alone is the concept realized.]''

In the Platonic presentation the subjective rules, insofar as now this (εἰ ἓν ἔστι [if the One is]),
now that (εἰ ἓν μὴ ἔστι [if the One is not]) is arbitrarily assumed; so subjective a form cannot
possibly be the Idea's adequate form. The Hegelian method, on the contrary, is objective, and can by
its systematic unfolding of the
% [d/b glyph flag in the Danish: print ``abæqvate'' for ``adæqvate'' (= adequate)]
content's moments furnish an objective, clear Idea-form. With his dissecting dialectic of distinction
Plato could not get the individual Ideas to arise as moments in the Idea's unity; with his immanent
dialectic of identity Hegel can point out the Idea's absolute unity in all the particular Ideas, in
that these are at once independent and yet moments.
% ---- printed p.158 (PDF 166) ----
``\textit{Die Idee selbst ist nicht zu nehmen als eine Idee von irgend Etwas, so wenig als der Begriff
bloß als bestimmter Begriff. Das Absolute ist die allgemeine und Eine Idee, welche als} \emph{urtheilend}
\textit{sich zum System der bestimmten Ideen besondert, die aber nur dieß sind, in die Eine Idee, in ihre
Wahrheit zurückzugehen. Aus diesem Urtheil ist es, daß die Idee} \emph{zunächst} \textit{nur die Eine,
allgemeine Substanz ist, aber ihre entwickelte wahrhafte Wirklichkeit ist, daß sie als Subjekt und so
als Geist ist.} [The Idea itself is not to be taken as an idea of something, as little as the concept
merely as determinate concept. The Absolute is the universal and one Idea, which, as \emph{judging},
particularizes itself into the system of the determinate Ideas, which, however, are only this: to go
back into the one Idea, into their truth. It is from this judgment that the Idea is \emph{at first} only
the one, universal substance, but its developed true actuality is that it is as subject and thus as
spirit.]''

\begin{center}\rule{0.35\linewidth}{0.4pt}\end{center}

\begin{center}
  {\large\bfseries B.}\par\vspace{2pt}
  {\bfseries Reality against the Idea.}\par\vspace{4pt}
  \rule{2cm}{0.4pt}
\end{center}
\phantomsection
\addcontentsline{toc}{subsection}{\texorpdfstring{B.\quad Reality against the Idea}{B. Reality against the Idea}}

\begin{center}{\bfseries \S\,30.}\end{center}

The Idea is then the truth, and the truth identity of essence and reality. He who in the real sees only
an essential, insofar as he therein sees something universal, sees only a half-truth; he who in reality
looks only at the individual, the special, the merely factual, likewise sees only a half-truth. In that
a half-truth is set as the whole truth, untruth arises. So long as Idealism as doctrine of essence and
Empiricism as doctrine of reality keep themselves in their finite opposition, no cognition of truth, no
true cognition of the Idea, will be possible.

That reality in the Idea is to correspond to the concept is, however, easily said; but the question is
whether human knowledge now also really masters the true unity of concept and reality. One thing is my
subjective concept of reality, another is reality's own objective con-
% ---- printed p.159 (PDF 167) ----
cept; one thing is reality's objective concept, another is reality itself. Reality's objective concept
is doubtless in the Idea one with its real content; but precisely because
% [d/b glyph flag in the Danish: print ``forbi'' for ``fordi'' (= because)]
the identity is only in the Idea, it is in reality not for us, when none of us is the Idea. The real, as
it is for us, is known precisely by this, that there is always a finite difference between content and
form. Man's essence is doubtless in every real man; but nonetheless no real man can in reality be the
same as man's Idea. If we think an ideal man, a man who could with truth say of himself: ``I am man's
Idea,'' as existing; then suffering and degradation would have to express that such a man's reality was
nonetheless in this present inadequate to its essence.

If the knower, by knowing the Idea, can really make himself one with the Idea, then there is in
philosophy an adequate cognition of the Idea, ``an absolute knowledge''; but, if every attempt to
identify oneself with the Idea itself must be rejected by science as an attempt in the mystically
fantastic, then ``the absolute knowledge'' is in reality unattainable.

Among the philosophers who have in earnest set themselves to grasp the real, to lift reality's content
to the form of the concept, and thereby to identity with the Idea, Aristotle and Hegel are without
comparison the most considerable; but it is so far from the case that either Aristotle or Hegel should
have solved the task, that the task, literally speaking, has not yet been present for either of them.

``The whole Aristotelian philosophy's riddle,'' says Apelt, ``lies now properly in the abstraction
εἶδος. First Aristotle teaches logically quite rightly that there are three principles (ἀρχαί):
ὑποκείμενον and the two ἀντικείμενα: εἶδος and στέρησις, i.e. the subject and the affirming and
denying predicate. Here, then, τὸ εἶδος is the predicate in the affirm-
% ---- printed p.160 (PDF 168) ----
ing judgment. But metaphysically the matter turns round for him, τὸ εἶδος transforms itself into τὸ τί,
i.e. into subject, it obtains essentiality. The metaphysical question concerning that which is to be τὸ
ὑποκείμενον or οὐσία, i.e. the essence of things, he decides thus, that the mass (ὕλη) merely exists
according to possibility (δυνάμει ὄν), the formed mass, on the contrary, in actuality (ἐντελεχείᾳ ὄν),
the shape (εἶδος or μορφή) is thus the ground of actuality. In this comparison of the grounds of being
with the categories of modality there is lacking the indication of the ground of necessity, i.e. the
law\ldots''

``This philosophy had two roots: ἐμπειρία and νοῦς, experience and pure reason. The Scholastics, who
developed the doctrine from its principles, kept νοῦς as the possibility of principles, and let ἐμπειρία
fall. Bacon, on the contrary, conversely rejects νοῦς, and thus ἐμπειρία becomes to him the source of
principles by means of ἐπαγωγή. `Induction,' so he teaches, `assumes no original concepts of the
understanding and does not rest satisfied with the immediate instruction of the senses.' The first is a
consequence of his rejecting νοῦς; the second is the Aristotelian separation between αἴσθησις and
ἐμπειρία\ldots'' Induction's object is thus τὸ εἶδος, forma\ldots. Τὸ εἶδος is, according to Aristotle,
something single, actual, substantial; according to Bacon, something universal and necessary, a
universal rule without essentiality: natural law.

Insofar as the Baconian proposition about natural law as universal rule without essentiality is to hold
for more than a phrase, it is (Cf. S. 60) only a proof of how far Empiricism in its strife against
Idealism can remove itself from the Idea. Empiricism rests on this, that the for-us real is and remains
inseparable from the finite, from the things and their parceling. If now reality's emphasis falls on
the finite facts, then the laws valid for these facts must also be grasped in their pure facticity,
% ---- printed p.161 (PDF 169) ----
not as metaphysical essentialities, but as universal rules found by induction. That laws can in reality
be discovered factually, and known in factual universality, before there can yet be talk of any
essential grounding, is shown by natural science in manifold examples. Held fast for itself, the
Baconian view is idealess, but as counterweight against the merely metaphysical it has its great
significance.

Reality wills to be known; if now the real belongs to the Idea's essence, then the cognition of reality
too must be an inviolable condition for the cognition of the Idea, and Empiricism, even if it knows it
not, find itself in the Idea's service. The factual relation between Speculation and Empiricism is,
however, not understanding, but strife.

\begin{center}{\bfseries \S\,31.}\end{center}

For the immediate consciousness it looks as though space and time properly did not concern each other,
as though one could well have space without time, and conversely, live in time without living in space.
In reality it is, however, otherwise. What time signifies is seen in motion upon space; what space
signifies is seen in motion upon time: motion expresses an intimate, essential relation of space and
time. This essential relation is specified into real relations, in that motion is separated into
definite motions; the difference between the motions existing in reality is a difference between a
slower and a faster, which is designated by velocity, in that velocity in turn designates the relation
of the traversed space and the time applied thereto, $v=\frac{s}{t}$.

That a thus-beginning analysis of reality does not without more let itself be transformed into a
logical, metaphysical analysis of essence is not merely in general clear from the Aristotelian physics,
but can even in more special traits be seen from the Hegelian blunders. Instead of following the
mathematical direction
% ---- printed p.162 (PDF 170) ----
Hegel will explain not merely velocity, but even the fall-law, from the pure concepts: space, time,
and corporeality.

``\textit{In dem Satze, daß die Geschwindigkeit den Zeiten proportional ist, ist die Geschwindigkeit
zunächst überhaupt gesagt; so wird sie überflüssigerweise mathematisch als} $\frac{s}{t}$ \textit{die
schlechtgleichförmige gesetzt, so die Kraft der Trägheit herein gebracht und ihr dieß Moment
zugeschrieben. Damit aber daß sie den Zeiten proportional sey, ist sie vielmehr als die gleichförmig
beschleunigte} $\frac{s}{t^2}$ \textit{bestimmt, und jene Bestimmung von} $\frac{s}{t}$ \textit{hat
hier keinen Platz, und ist ausgeschlossen.} [In the proposition that velocity is proportional to the
times, velocity is at first stated quite generally; so it is superfluously set mathematically as
$\frac{s}{t}$, the merely-uniform, so the force of inertia is brought in and this moment ascribed to
it. But in that it is proportional to the times, it is rather determined as the uniformly accelerated
$\frac{s}{t^2}$, and that determination $\frac{s}{t}$ has here no place, and is excluded.]''
(Naturphil. S. 87.)

Now follows the speculative explanation:

``\textit{Das Gesetz des} \emph{Falles} \textit{ist gegen die abstracte gleichförmige Geschwindigkeit
des todten von Außen bestimmten Mechanismus ein} \emph{freies} \textit{Naturgesetz; d. h. es hat eine
Seite in ihm, die sich aus dem} \emph{Begriffe} \textit{des Körpers bestimmt. Indem daraus folgt, daß
es aus diesem muß abgeleitet werden können, so ist dieses sich vorzusetzen, und der Weg anzugeben, wie
das Galileische Gesetz,} ``\textit{daß die durchlaufenen Räume sich wie die} \emph{Qvadrate} \textit{der
verflossenen Zeiten verhalten,''} \textit{mit der Begriffsbestimmung zusammenhängt.}

``\textit{Dieser Zusammenhang ist aber als einfach darin liegend anzusehen, daß, weil hier der Begriff
zum Bestimmen kommt, die Begriffsbestimmungen der} \emph{Zeit} \textit{und des Raums gegeneinander}
\emph{frei} \textit{werden, d. i. ihre} \emph{Größebestimmungen} \textit{sich nach denselben verhalten.
Nun ist aber die Zeit das Moment der} \emph{Negation}\textit{, des Fürsichseyns, das Princip des Eins;
und ihre Größe --- irgend eine empirische Zahl --- ist im Verhältnisse zum Raum als die} \emph{Einheit}
\textit{oder als Nenner zu nehmen. Der} \emph{Raum} \textit{dagegen ist das} \emph{Außereinandersehn}
\textit{, und zwar} \emph{keiner andern Größe}\textit{, als eben der Größe der Zeit; denn die
Geschwindigkeit dieser} \emph{freien} \textit{Bewegung ist dieß, daß Zeit und Raum nicht} \emph{äußerlich}
\textit{, nicht zufällig
% ---- printed p.163 (PDF 171) ----
gegen einander sind, sondern Beider} \emph{Eine} \textit{Bestimmung ist. Die der} \emph{Einheit}\textit{,
als der Form der Zeit, entgegengesetzte Form des} \emph{Außereinander} \textit{des Raums, und zwar ohne
daß irgend eine andere Bestimmtheit sich einmischt, ist das} \emph{Qvadrat}\textit{: die Größe außer
sich kommend in eine zweite Dimension sich setzend, sich somit vermehrend, aber nach} \emph{keiner
andern als ihrer eignen} \textit{Bestimmtheit --- diesem Erweitern sich selbst zur Grenze machend, und
in ihrem Anderswerden so sich nur auf sich beziehend.}

\textit{Dieß ist der Beweis des Gesetzes des} \emph{Falls} \textit{aus dem} \emph{Begriffe der Sache}\textit{.}''
% [German multi-paragraph quote (Naturphil. S. 87): the opening quote is repeated at each paragraph
%  start but closed only once at the end — a print convention reproduced verbatim; this adds +1 to the
%  „/“ (``/'') balance. English rendering of the whole quote follows:]
[The law of \emph{fall} is, against the abstract uniform velocity of the dead, from-without-determined
mechanism, a \emph{free} natural law; i.e. it has a side in it that is determined from the \emph{concept}
of the body. Since it follows from this that it must be able to be derived from this, it is to be set
before oneself, and the way indicated how the Galilean law, ``that the traversed spaces relate as the
\emph{squares} of the elapsed times,'' hangs together with the conceptual determination. This connection,
however, is to be regarded as lying simply therein, that, because here the concept comes to determining,
the conceptual determinations of \emph{time} and of space become \emph{free} against each other, i.e.
their \emph{magnitude-determinations} relate accordingly. Now time is the moment of \emph{negation}, of
being-for-self, the principle of the one; and its magnitude --- some empirical number --- is to be taken,
in relation to space, as the \emph{unit} or as denominator. \emph{Space}, on the contrary, is the
\emph{being-outside-one-another}, and indeed of \emph{no other magnitude} than precisely the magnitude of
time; for the velocity of this \emph{free} motion is this, that time and space are not \emph{external},
not contingent against each other, but are the \emph{one} determination of both. The form of the
\emph{being-outside-one-another} of space, opposed to that of the \emph{unit} as the form of time --- and
indeed without any other determinacy mixing in --- is the \emph{square}: magnitude coming outside itself,
setting itself in a second dimension, thus multiplying itself, but according to \emph{no other than its
own} determinacy --- making this expanding itself into its limit, and in its becoming-other thus relating
only to itself. This is the proof of the law of \emph{fall} from the \emph{concept of the matter}.]

In this ``Beweis aus dem Begriffe der Sache'' [proof from the concept of the matter] there shows itself,
as everywhere, the Hegelian eye for essential relations; ``daß Zeit und Raum nicht äußerlich, nicht
zufällig gegen einander sind'' [that time and space are not external, not contingent against each other]
is the midpoint of thought about which the whole turns; but that Hegel just as little here as elsewhere
has an eye for the proper analysis of reality is at the same time conspicuous, and betrays itself
involuntarily by his dismissal of the formula $\frac{s}{t}$.

Already at the beginning (S. 85), where he does concede, concerning the laws of motion: ``es sind
unsterbliche Entdeckungen, die der Analyse des Verstandes die höchste Ehre machen'' [they are immortal
discoveries that do the highest honor to the analysis of the understanding], he comes, by mistaking the
mathematical-physical conception, into a skewed position toward the matter. ``Die Voraussetzung bei
diesem apriorischen Beweise ist, daß die Geschwindigkeit im Fall \emph{gleichförmig} beschleunigt ist;
der Beweis aber besteht in der Verwandlung der Momente der mathematischen Formel in physicalische Kräfte,
in eine \emph{beschleunigende} Kraft, welche in jedem Zeitmoment einen (denselben) Impuls mache, und in
eine Kraft der \emph{Trägheit}, welche die in jedem Zeitmomente erlangte (größere) Geschwindigkeit
fortsetze --- Bestimmungen, die durchaus ohne empirische Beglaubigung sind, so wie der Begriff nichts mit
ihnen zu thun hat.'' [The presupposition in this a priori proof is that velocity in fall is
\emph{uniformly} accelerated; but the proof consists in the transformation of the moments of the
mathematical formula into physical forces, into an \emph{accelerating} force which in each time-moment
produces one (the same) impulse, and into a force of \emph{inertia} which continues the (greater) velocity
attained in each time-moment --- determinations that are utterly without empirical accreditation, just as
the concept has nothing to do with them.]
% ---- printed p.164 (PDF 172) ----
With this dismissal Hegel has broken the staff over every analysis of reality; for every such analysis is
conditioned by this, that the factors are set apart in such a way that their cooperation lets itself be
calculated.

The phenomenon dealt with is, as is known, this, that a body which under the fall-motion traverses in the
first second the space g, traverses in the second second 3 g, in the third sec. 5 g, and so on. The
emphasis thus falls here, not, as Hegel supposes, on a squaring of the time by means of the space, but on
the evenly increasing velocity with which the motion takes place. Although a second is only a very small
part of time, it nonetheless contains an infinite number of moments; in each moment there is now traversed
an infinitely small space, r, so that the series becomes $r + 2r + 3r + \ldots nr$. By integration one
obtains $r\int_0^n p\,dp = \frac{n^2 r}{2}$, $r\int_n^{2n} p\,dp = 3\frac{n^2 r}{2}$, $r\int_{2n}^{3n}
p\,dp = 5\frac{n^2 r}{2} \ldots$ or, when $\frac{n^2 r}{2}$ is set equal to g, the series: g, 3 g, 5 g, as
above indicated.

For the explanation of the phenomenon Hegel adduces a speculative ground instead of an efficient cause.
The fall-motion takes place with evenly increasing velocity, in order that natural reason may hereby get
occasion to square the time by means of the space. Such speculative motives can with right be reckoned to
``Bestimmungen, die durchaus ohne empirische Beglaubigung sind'' [determinations that are utterly without
empirical accreditation]. Just as little can the addition ``Nur das Suchen des Centrums ist im Fall die
absolute Seite'' [Only the seeking of the center is in the fall the absolute side] be said to hit the
real. When physics, on the contrary, seeks the efficient cause in the cooperation of matter's inertia and
the earth's attraction, the explanation has not only an empirical confirmation in the circumstance that
the calculation built thereon really holds, but has even brought out a character in matter which in all
experiments has shown itself constant, namely inertia. That inertia is not sluggishness, but on the
contrary that property of matter in virtue of which it is in equal
% ---- printed p.165 (PDF 173) ----
degree disposed both to motion and to rest, is seen precisely from the example before us. Set in motion by
an external cause, the earth's attraction, the falling body takes up the imparted motion into itself, and
keeps it; since now the continued attraction in each Now imparts to matter an equal motion, the motion must
grow evenly in each moment, and the velocity thus be evenly increasing. If there is now indicated the space
traversed in each finite definite part of time, in the 1st, 2nd, 3rd, etc. second, then it shows itself
rightly enough that the traversed spaces relate as the squares of the applied times; but this relation is
by no means restricted to the fall-motion, or ``ein Suchen des Centrums'' [a seeking of the center],
% [d/b glyph flag in the Danish: print ``forbi'' for ``fordi'' (= because)]
because the fall-motion is one thing falling thereunder.

If we now return to the formula $v = \frac{s}{t}$, it is clear that velocity is not held fast abstractly
as an empty concept, but that the real velocity presupposes a matter whose velocity can be determined. In
the unchangeable velocity, inertia then expresses the maintenance of the speed imparted by the force in the
moment of impact; whereas the changing velocity, on the contrary, presupposes a peculiar influence of a
moving force.

In that the analysis of reality in this way separates the factors, these factors cannot indeed, like the
logical reflexes of the analysis of essence, be transformed into immanent difference-determinations in an
absolute identity; but nonetheless these factors in their cooperation do furnish peculiar conceptual
determinations: reality is not conceptless, only that a separation of finitude belongs essentially in the
concept of reality.

Insofar as $s$ is a space already traversed in the time $t$, one cannot yet know whether this space is
traversed with changing or unchanging velocity. To determine the velocity, one must give the time an
increment, that
% ---- printed p.166 (PDF 174) ----
is to say: procure time for continued motion. If the increment to the time is now called h, one has the
series:
\[ f(t+h) = s + \frac{ds}{dt}\,\frac{h}{1} + \frac{d^2 s}{dt^2}\,\frac{h^2}{1.2}
+ \frac{d^3 s}{dt^3}\,\frac{h^3}{1.2.3} + \ldots. \]
If the increment h is a finite magnitude, and the velocity constant, the whole series will restrict itself
to $f(t+h) = s + \frac{ds}{dt}\frac{h}{1}$; if h is infinitely small, equal to dt, the same will hold, even
when the velocity is variable. On the contrary, when h is finite, and the velocity variable, the members of
the series can be increased to infinity.

Seen from the concept's point of view, the series indicates a totality of objective reflection-
determinations. The individual members can, in that s means the space, v the velocity, and $\varphi$ the
force, be rendered by different expressions corresponding to one another:
\[ \begin{aligned}
f(t+h) &= s + \frac{ds}{dt}\,\frac{h}{1} + \frac{d^2 s}{dt^2}\,\frac{h^2}{1.2} +
\frac{d^3 s}{dt^3}\,\frac{h^3}{1.2.3} + \\
&= .. + v\,\frac{h}{1} + \frac{dv}{dt}\,\frac{h^2}{1.2} + \frac{d^2 v}{dt^2}\,\frac{h^3}{1.2.3} + \\
&= .. + \quad .. \quad + \varphi\,\frac{h^2}{1.2} + \frac{d\varphi}{dt}\,\frac{h^3}{1.2.3} +
\end{aligned} \]
The categories \emph{Motion}, \emph{Velocity}, \emph{Force}, and \emph{Change in Force} are thus
determined by objective reflection. Immediately s designates a length of space; but in the equation
$s = f(t)$ this length of space is seen in reflection of a length of time; it owes indeed its existence to
time, insofar as it is namely traversed in a given time. Immediately v designates a velocity, but in the
expression $\frac{ds}{dt}$ the concept of velocity is determined by objective reflection of space and time
in the given relation. Immediately the force is designated by $\varphi$, next by reflection of velocity and
time in the definite relation $\frac{dv}{dt}$, finally, in that $v = \frac{ds}{dt}$ goes up in reflection of
space and time, by a doubled reflection in the relation
% ---- printed p.167 (PDF 175) ----
$\frac{d^2 s}{dt^2}$. In this way the changes of force can, with respect to the relation of space and time,
be determined as reflections of reflections to infinity.

The difference between analysis of essence and analysis of reality can be seen in ``the moment.'' In the
analysis of essence the moment is the pure Now, a dialectical playing between being and not-being; in the
analysis of reality the moment is, on the contrary, qualified: all the members following on s in the series
set up are sheer qualified ones, each for itself designating a peculiar disposition, a peculiar energy,
moments. In this qualification the separation is exact, and the concept of reality recognizable.

By excluding the speculative motives, asserting the incommensurable moment of finitude, and pursuing the
exact discretion to infinity, the analysis of reality has broken with the pure essence, and can on that
account not relate to the Idea as Idea.

\begin{center}{\bfseries \S\,32.}\end{center}

In his Physics Aristotle (cf. Brandis 2ter Th., 2 Abth. 2) has with great circumspection proved that,
since all motions are either circular or rectilinear or a mixture of both, the circular motion must be the
perfect. But, in that the Greek astronomers started from the metaphysical presupposition that the circular
motion, in which every point can be regarded as beginning, middle, and end, was the heavens' eternal
motion, and that the earth was this motion's center, it is remarkable to see how observations and
calculations gradually brought them into embarrassment with that speculative presupposition, and in
manifold particulars forced them to a break with the metaphysical.

Already the observation that the sun in winter moves away from a fixed star by a larger arc in equal time
than in summer was sufficient to prove that, if the motion were in itself even and uniform, the earth could
% ---- printed p.168 (PDF 176) ----
not be center for the circle in which the sun's yearly motion took place. Ptolemy therefore proposed to
explain the relation by letting the uniform motion take place, not about the earth as center, but about the
so-called \textit{punctum æquatorium} [equating point] or \textit{p. æquans} [equant]. It has in recent
times been proved that this \textit{punctum æquans} corresponds fairly well to the planetary orbit's second
focus, in that the motion, seen from there, approaches becoming uniform.

The decisive break between the astronomical and the speculative, which recent time has properly first
carried through, by no means proceeds from this, that the Ptolemaic system has had to give way to the
Copernican, or the Aristotelian circle to the Keplerian ellipse.

The Copernican world-structure is, if one will, an even more grandiose image of the Idea than the Ptolemaic;
wherefore Copernicus too appeals to the Idea. ``By no other ordering,'' he says enthusiastically (Humboldt:
Kosmos, 2nd volume), ``have I been able to find so admirable a symmetry in the universe, so harmonious a
connection of the orbits, as when I set \emph{the Sun}, this world-lantern (\textit{lucernam mundi} [the
lamp of the world]), steering the whole family of circling heavenly bodies (\textit{circumagentem gubernans
astrorum familiam} [governing and steering the circling family of the stars]) in the midst of the beautiful
temple of nature as upon a royal throne. \textit{Quis enim in hoc pulcherrimo templo lampadem hanc in alio
vel meliori loco poneret, quam unde totum simul possit illuminare?} [For who in this most beautiful temple
would place this lamp in another or better place than whence it can at once illuminate the whole?]''
(\textit{Nicol. Copern. de Revol. orbium coelestium lib. I cap. 10}).

That Speculation can find the ellipse just as rational as the circle, indeed even more rational, Hegel
proves.

``\textit{Was nun die Gestalt der Bahn betrifft, so ist der} \emph{Kreis} \textit{nur als die Bahn einer}
\emph{schlecht-gleichförmigen} \textit{Bewegung zu faßen.} \emph{Denkbar}\textit{, wie man es nennt, ist es
wohl, daß auch eine gleichförmig zu- und abnehmende Be-
% ---- printed p.169 (PDF 177) ----
wegung im Kreise geschehe. Aber diese Denkbarkeit oder Möglichkeit heißt nur eine abstracte
Vorstellbarkeit, welche das Bestimmte, worauf es ankommt, wegläßt, und daher nicht nur oberflächlich,
sondern falsch ist. Der Kreis ist die in sich zurückkehrende Linie, in der alle Radien} \emph{gleich}
\textit{sind: d. h. er ist durch den Radius vollkommen bestimmt; es ist dieß nur} \emph{Eine}\textit{, und
zwar die} \emph{ganze} \textit{Bestimmtheit. In der freien Bewegung aber, wo räumliche und zeitliche
Bestimmung in} \emph{Verschiedenheit}\textit{, in ein qvalitatives Verhältniß zu einander treten, tritt
nothwendig dieß Verhältniß} \emph{an dem Räumlichen} \textit{selbst als eine} \emph{Differenz} \textit{desselben
hervor, welche hiermit zwei Bestimmungen erfordert. Dadurch wird die Gestalt der in sich zurückgehenden
Bahn wesentlich eine} \emph{Ellipse}\textit{; --- das erste der Kepplerischen Gesetze} [As now for the shape
of the orbit, the \emph{circle} is to be grasped only as the orbit of a \emph{merely-uniform} motion.
\emph{Thinkable}, as one calls it, it may well be that a uniformly in- and decreasing motion too takes place
in a circle. But this thinkability or possibility means only an abstract representability, which omits the
determinate on which it depends, and is therefore not merely superficial, but false. The circle is the
line returning into itself, in which all radii are \emph{equal}: i.e. it is completely determined by the
radius; this is only \emph{one}, and indeed the \emph{whole} determinacy. But in the free motion, where
spatial and temporal determination step into \emph{difference}, into a qualitative relation to each other,
this relation necessarily steps forth \emph{on the spatial} itself as a \emph{difference} of it, which
hereby requires two determinations. Thereby the shape of the orbit returning into itself becomes
essentially an \emph{ellipse} --- the first of the Keplerian laws]''. (Naturphil. S. 104).

No, astronomy's and, on the whole, natural science's decisive break with Speculation is grounded in this,
that natural law everywhere necessarily must look at the efficient causes, whereas Speculation with its
metaphysical principles and essences will explain everything from purpose-causes. How little
purpose-causes help physics, F. Bacon has already expressed in the well-known proposition: \textit{Causarum
finalium inquisitio sterilis est, et tanquam virgo Deo consecrata nihil parit (De Augm. Scient. Lib. III)}
[The inquiry into final causes is sterile, and, like a virgin consecrated to God, bears nothing].

Leibniz, though a naturalist, held to the final causes (\textit{Bien loin d'exclure les causes finales ---
c'est de là qu'il faut tout déduire en Physique \ldots{} Extrait d'une lettre à Mr. Bayle sur un principe
général utile à l'explication des loix de la nature} [Far from excluding final causes --- it is from these
that one must deduce everything in physics \ldots{} Extract from a letter to Mr. Bayle on a general principle
useful for the explanation of the laws of nature]),
% [d/b glyph flag in the Danish: print ``forbi'' for ``fordi'' (= because)]
because he held to Speculation, and Hegel has still, as far as the astronomical phenomena are concerned,
such trust in the wisdom expressed in the concept that he even rebukes Newton and the astronomers,
% [d/b glyph flag in the Danish: print ``forbi'' for ``fordi'' (= because)]
because they, with respect to the motions of the heavenly bodies, adduce the same physically efficient
causes as those that hold in the finite on earth, whereas it must nonetheless in virtue
% ---- printed p.170 (PDF 178) ----
of the concept be settled that there is another causality in the whole than in the parts, other causes in
the heavenly than in the earthly.

``\textit{Da der Mond sechzig Durchmesser der Erde von der Erde entfernt ist, so wird also das Moment der
Attraction in seiner Bewegung danach bestimmt. Es wird dann gefunden, daß das, was die Attractivkraft der
Erde auf den Mond wirke (der} \textit{Sinus versus}\textit{, die} \textit{Sagitta}\textit{), zugleich den
ganzen Umlauf des Mondes bestimme: er falle eben so. Das mag richtig seyn.} [Since the moon is sixty
diameters of the earth distant from the earth, the moment of attraction in its motion is accordingly
determined. It is then found that that which the earth's attractive force acts upon the moon (the
\textit{sinus versus}, the \textit{sagitta}) determines at the same time the whole revolution of the moon:
it falls just so. That may be right.]'' (Naturphil. S. 111).

The moon thus stands, although it otherwise in its motion floats freely, like other heavenly bodies, yet
not high enough on the Idea's heaven to be emancipated from the finite, physically efficient causes.

``\textit{Das ist aber zunächst nur ein einzelner Fall, die Ausdehnung des empirischen Falls auf der Erde
auf den Mond. Von den Planeten ist dieß nicht gemeint, oder gälte nur von ihnen im Verhältniß zu ihren
Trabanten. Das ist also ein beschränkter Punkt.} [But that is at first only a single case, the extension of
the empirical fall on the earth to the moon. Of the planets this is not meant, or would hold only of them
in relation to their satellites. That is thus a restricted point.]'' (S. 111).

While, then, the bodies on earth are moved by physical causes, and the moon as satellite cannot either
raise itself to the speculative, the planetary system is, in its wholeness, on the contrary subjected to a
metaphysical causality in which the finite causes have vanished.

``\textit{Man sagt, den himmlischen Körpern kommt das Fallen zu. Sie fallen aber doch nicht in die Sonne;
so giebt man ihnen noch eine andere Bewegung, welche den Fall aufhält. Das ist sehr einfach verendlicht. So
schlagen Knaben mit dem Prügel einen Ball, der fallen will, auf die Seite. Es ist uns nicht geheuer, solche
Knabenverhältnisse auf diese freie Bewegung angewendet zu sehen.} [One says, falling belongs to the heavenly
bodies. But they do not fall into the sun; so one gives them yet another motion, which halts the fall. That
is very simply finitized. So boys strike with a club a ball that wants to fall, off to the side. It seems to
us uncanny to see such boys' relations applied to this free motion.]'' (Naturphil. S. 111).

It is then not the relation between the planet's own motion and the sun's attraction, not a cooperation of
centrifugal and centripetal force, that determines the orbit;
% ---- printed p.171 (PDF 179) ----
no, it is reason itself, the concept in its system of judgments and inferences, in one word: the Idea, that
moves the whole.

``\textit{Die Bewegung der Himmelskörper ist nicht ein solches Hin- und Hergezogenseyn, sondern die freie
Bewegung; sie gehen, wie die Alten sagten, als selige Götter einher. Die himmlische Körperlichkeit ist
nicht eine solche, welche das Princip der Ruhe oder Bewegung außer ihr hätte. Weil der Stein träge ist, die
ganze Erde aber aus Steinen besteht, und die anderen himmlischen Körper eben dergleichen sind, --- ist ein
Schluß, der die Eigenschaften des Ganzen denen des Theils gleichsetzt. Stoß, Druck, Widerstand, Reibung,
Ziehen und dergleichen gelten nur von einer andern Existens der Materie, als die himmlische Körperlichkeit.
Das Gemeinschaftliche Beider ist freilich die Materie, so wie ein guter Gedanke und ein schlechter beide
Gedanken sind: aber der schlechte nicht darum gut, weil der gute ein Gedanke ist.} [The motion of the
heavenly bodies is not such a being-drawn-hither-and-thither, but the free motion; they go, as the ancients
said, along like blessed gods. The heavenly corporeality is not such as would have the principle of rest or
motion outside it. Because the stone is inert, but the whole earth consists of stones, and the other
heavenly bodies are just the same --- is an inference that equates the properties of the whole with those of
the part. Impact, pressure, resistance, friction, pulling, and the like hold only of another existence of
matter than the heavenly corporeality. What is common to both is indeed matter, just as a good thought and a
bad are both thoughts: but the bad is not good on that account, because the good is a thought.]'' (Naturphil.
S. 97).

So long as one held fast as a decisive dilemma: either nature's order is an effect of blind chance, or else
a work of God's wisdom, it is comprehensible that men like Kepler, Newton, and Leibniz could not merely
unite a penetrating study of nature's laws with faith in a creating God, but even in decisive cases feel an
urge to lay this their religious presupposition at the ground of a deeper interpretation of the phenomena.
\textit{Finalem causam non tantum prodesse ad virtutem et pietatem, in Ethicà et Theologia naturali, sed
etiam in ipsa Physica ad inveniendum et detegendum abditas veritates. De ipsa natura}; Leibniz. [That the
final cause is of use not only for virtue and piety, in ethics and natural theology, but also in physics
itself for finding and detecting hidden truths. On Nature Itself; Leibniz.]

The more, on the contrary, the intermediate causes were traced out, the more the hypothetically necessary
(τὸ ἐξ ὑποθέσεως ἀναγκαῖον [that which is necessary on a hypothesis]) took contingency in seizure, the more
the thought pressed itself forward, that nature's order itself must be a work of nature.
% ---- printed p.172 (PDF 180) ----
It was then not a willing God, acting according to purposes, that gave the planet that first impulse, as
little as it was a willing God who arbitrarily commanded matter to act attractively. No, that matter
attracts follows from its nature. The original impulse, motion's natural cause, is itself a natural effect
of a preceding natural cause (the ring's rotation and bursting, according to the hypothesis), this cause
again an effect of a preceding one, and so on in infinity's cycle.

``I have searched the whole heaven,'' says the astronomer Lalande, ``and nowhere found a trace of God.'' ---
When the Emperor Napoleon asked Laplace why in his \textit{Traité de mécanique céleste} there was nowhere
talk of God, the famous astronomer is said to have given the answer: »\textit{Sire, je n'avais pas besoin de
cette hypothèse!}« [Sire, I had no need of this hypothesis!] From this answer one can conclude nothing
concerning Laplace's religiosity, but on the contrary much concerning his scientificity; he has not dwelt on
edifying general considerations, but interpreted facts; he has not immersed himself in dialectical
speculations, but in exact analyses, not intuited the pure essence, but tested the strict reality, not had
need of \textit{causæ finales} [final causes], but helped himself with \textit{causæ efficientes} [efficient
causes].

It is not only the application of the supernatural causality, or of the finite purposiveness (die äußere
Zweckmässigkeit [the outer purposiveness]), that natural science must reject; even the immanent, speculative
teleology has in the analysis of reality nothing to signify.

That the doctrine of ``die besondern Formationen der Erde'' [the particular formations of the earth] belongs
under physical geography, and that in physical geography there can be no talk of any higher teleology, Hegel
concedes. But a description of the physically mechanical is then also only of subordinate significance.
``\textit{Dieß erscheint zunächst als zufällig; aber die Thätigkeit des Begriffs ist, das als nothwendig
bestimmt}
% [trailing German quote continues into next batch (p.173); the opening `` stays open here —
%  it closes on p.173 — reproducing the source's page-spanning quote.]
% ---- printed p.173 (PDF 181) ----
\textit{zu faßen, was dem sinnlichen Bewußtseyn als zufällig erscheint.}'' (Naturphil. S. 441). [This
appears at first as contingent; but the activity of the concept is to grasp as necessary that which
appears to the sensible consciousness as contingent.]

The necessity of which there is here talk is not the blind fatal one that is dialectically one with the
contingent; no, the necessity in which the concept is to sublate the contingent has its deeper root in
the teleological. In the separation of land and sea, in the division of the new and the old world, the
Idea has revealed its wisdom:

``\textit{Ueberhaupt (Naturphil. S. 442) stellt die neue Welt die unausgebildete Entzweiung dar, ---
einen nördlichen und einen südlichen Theil in der Weise des Magneten: die alte aber die vollkommene
Entzweiung in drei Theile, deren Einer, Afrika, das gediegene Metall, das Lunarische, starr vor Hitze
ist, wo der Mensch in sich selbst verdumpft, --- der nicht ins Bewußtseyn tretende stumme Geist; der
andere, Asien, ist die bacchantisch kometarische Ausschweifung, die wild nur aus sich ausgebährende
Mitte, die formlose Erzeugung, ohne daß er über seine Mitte Meister werden könnte; der dritte aber,
Europa, bildet das Bewußtseyn, den vernünftigen Theil der Erde, das Gleichgewicht von Strömen und
Thälern und Gebirgen, --- dessen Mitte Deutschland ist.} [In general, the new world presents the
undeveloped division --- a northern and a southern part in the manner of the magnet: the old, however,
the complete division into three parts, of which one, Africa, the massive metal, the lunar, is rigid
from heat, where man becomes dull within himself --- the mute spirit not stepping into consciousness;
the second, Asia, is the bacchantic, cometary excess, the wildly self-generating middle, the formless
production, without his being able to become master over his middle; the third, however, Europe, forms
the consciousness, the rational part of the earth, the equilibrium of streams and valleys and
mountains --- whose middle is Germany.]''

\begin{center}{\bfseries \S\,33.}\end{center}

In the doctrine of Ideas matter is only a fleeting moment, a remainder of finitude vanishing in the
negation of the negation. Although Hegel and Aristotle ascribe to matter's being another significance
than Plato does, they are nonetheless all agreed in making matter into so formal a material that its
being according to possibility only constitutes the determining form's passive side.

That Plato ascribes to the fundamental parts in the four elements definite mathematical figures --- to
the fire-parts the pyramid, to the earth-parts the cube, to the air-parts the octahedron, to the
water-parts the icosahedron --- cannot possibly be misunderstood as though he hereby
% ---- printed p.174 (PDF 182) ----
thought of eternally existing, absolutely independent atoms; since he on the contrary grasps matter as
the not-being. Aristotle does indeed let matter step forth with a certain independence against form, in
that it makes form resistance, and brings it about that the one form falls out from itself into a
plurality of forms; but the role that matter hereby plays is nonetheless, as in the Hegelian
``Außersichseyn'' [being-outside-itself], ``Andersseyn'' [being-other], only that of passivity and
impotence: matter is in its being according to possibility a pure continuum, and its parts only
points of transition for the form's, i.e. the Idea's, infinite self-mediation.

It is otherwise in the exact analyses of reality. Without deciding what matter is in metaphysical
essentiality, Empiricism holds fast to the view that the real mass, though it may according to
possibility be divisible to infinity, nonetheless in reality always consists and must consist of finite
fundamental parts, since the presupposition of matter's finite discretion is an inviolable condition for
every natural-scientific investigation grounded on calculations and on weight and measure (cf.
\S\,10 S. 55 with \S\,11 S. 64).

To the difference that thus is between the idealistic and the empirical in the doctrine of matter
corresponds a still more thoroughgoing difference in the doctrine of form. However much the Aristotelian
form in its unity with matter otherwise diverges from Plato's for-itself-being Idea, it is nonetheless
this, to furnish type, παράδειγμα [paradigm], common to both. The type, the definite form, whose
realization is the last, is in concept and essence the proper first. In the formed matter the type is
teleologically realized by the typically acting principle of unity: the material constituents with
their finite reciprocity are here the subordinate. From the empirical standpoint the matter is seen
reversed; here it is precisely the material fundamental parts and their reciprocity on which the
emphasis falls. Metaphysically seen, the type is principle, physically seen, it is result;
metaphysically
% ---- printed p.175 (PDF 183) ----
seen, the aggregation is subordinate to the paradigm, physically seen, the paradigm is subordinate to
the aggregation.

``\textit{Die Thätigkeit in ihr Product übergegangen ist die Gestalt, und bestimmt als Krystall. In
dieser Totalität sind die differenten magnetischen Pole zur Neutralität reducirt, die abstracte
Linearität der ortbestimmenden Thätigkeit zur Fläche und Oberfläche des ganzen Körpers realisirt
\ldots{} Es wirkt die Eine Form, indem sie} α) \textit{die Kugel begrenzend, den Körper nach Außen
krystallisirt, und} β) \textit{die Punktualität gestaltend, seine innere Continuität durch und durch im
Durchgang der Blätter, d. h. in der Kerngestalt, krystallisirt} [The activity, having passed over into
its product, is the shape, and determined as crystal. In this totality the different magnetic poles are
reduced to neutrality, the abstract linearity of the place-determining activity realized into the
surface and outer surface of the whole body \ldots{} The one form acts, in that it α) bounding the
sphere, crystallizes the body outward, and β) shaping the punctuality, crystallizes its inner
continuity through and through in the passage of the leaves, i.e. in the core-shape.]'' (Hegel
Naturphil. S. 265).

From this speculative interpretation of the crystal it is seen that the typical form is in Hegel just as
independently self-active as in Aristotle. That this self-activity of the form must expressly be grasped
teleologically goes without saying.

``\textit{Der Magnet ist noch nicht zweckmäßig; denn seine Entzweiten sind noch nicht gleichgültig,
sondern nur rein Nothwendige für einander. Hier aber ist eine Einheit Gleichgültiger, oder Solcher,
deren Daseyn in seiner Beziehung frei von einander ist. Die Linien des Krystalls sind diese
Gleichgültigkeit: es kann eine von der andern getrennt werden, und sie bleiben; aber sie haben
schlechthin Bedeutung nur in Beziehung auf einander, --- der Zweck ist diese ihre Einheit und
Bedeutung.}

\textit{Indem der Krystall aber dieser ruhige Zweck ist, so ist die Bewegung ein Anderes, als sein
Zweck; der Zweck ist noch nicht als Zeit} [The magnet is not yet purposive; for its sundered ones are not yet
indifferent, but only purely necessary for one another. But here is a unity of indifferent ones, or of
such whose existence is in its relation free from each other. The lines of the crystal are this
indifference: one can be separated from the other, and they remain; but they have significance simply
only in relation to one another --- the purpose is this their unity and significance. But in that the
crystal is this resting purpose, the motion is an other than its purpose; the purpose is not yet as
time]''. (Naturphil. S. 267).

The real doctrine of crystals now knows indeed nothing at all of this ``ruhige Zweck'' [resting
purpose], and attains nothing either by appealing to the typically active form. On the contrary,
Empiricism, by seeing the crystallogonic in its intimate connection with the chemical, is everywhere ---
isomorphism, polymorphism --- attentive to the atoms' size, shape,
% ---- printed p.176 (PDF 184) ----
number, and arrangement; Empiricism understands that it can make progress in the interpretation of the
real crystal-formations only insofar as it gradually discovers the laws for the definite atoms' definite
aggregations.

Since the theory of aggregation plays so comprehensive a role in the science of experience, it goes
without saying that the category \emph{aggregation} must be taken now in a stricter, now in a wider, now
in a mechanical, now in a chemical, now in an organic sense.

``All plants and animals begin \ldots{} in the form of a simple cell; the cell itself is formed out of a
mother-substance (Cytoblastema), in such a way that several microscopic small grains gather in a group
(the nucleus), around which lay themselves others, less firm parts of the mother-substance and an outer
thin membrane. In crystallization there are separated out of a solution or, if one will, a
mother-liquid, homogeneous firmer substances and around these others homogeneous, the smallest
precipitated parts constantly observing a certain position to one another mutually, whereby there arises
a regular figure, different according to the substance's chemical nature. It is true that a certain
difference between that cell-formation and this crystallization cannot be mistaken. The cells are round,
soft, hollow, and filled with a liquid; the crystals angular, firm, only in certain cases containing a
crystallization-liquid and in comparatively small quantity''. (D. F. Eschricht: \emph{Læren om Livet}
[The Doctrine of Life], S. 35).

The difference between the organic points and the crystal-points may otherwise be ever so thoroughgoing;
so is nonetheless, according to what is hereby made clear, punctuality and aggregation a common
condition for crystal- and plant-formation.

In that natural science looks at the types, it can in the different types very well see an expression of
the teleological; but as natural science it must nonetheless in the very strictest sense assert the
principle that the realization of any type whatever is explained by the mutually attracting
% ---- printed p.177 (PDF 185) ----
and repelling fundamental parts' separation, union, and arrangement proceeding in such a way from the
parts' original disposition that a type is realized only in that the strict law is by force of the
hypothetically necessary fulfilled point for point.

If we now consider the difference, grounded in the nature of the matter, between an a priori doctrine of
Ideas and an empirical doctrine of reality, we can understand that each of these directions must be
related to a peculiar corresponding view of the world. That a one-sidedly continued Speculation must
dull the eye for the real is just as psychologically explicable as that a persevering analysis of
reality must make consciousness foreign to the Idea as Idea.

To this come still such idealistic blunders and fantastic encroachments as could not do otherwise than
bring the again-refreshed Materialism into a rage.

Thus supposedly scientific efforts to hold fast to God's supernatural essence in the form of speculative
Idea have, among Idealism's opponents, served only to nourish the delusion that science can in earnest
disprove God's being. ``\textit{Der außer- und übermenschliche Gott ist nichts Anderes, als das außer-
und übernatürliche Selbst, das seinen Schranken entrückte, über sein objectives Wesen gestellte
subjective Wesen des Menschen.} [The extra- and superhuman God is nothing other than the extra- and
supernatural self, the subjective essence of man withdrawn from its barriers, set above its objective
essence.]'' (Feuerbach).

Thus supposedly scientific attempts to prove the soul's immortality by force of the Idea's demand have
served only to call forth proofs of the soul's mortality by force of reality's conditions. ``\textit{Es
ist so unmöglich, daß ein unversehrtes Gehirn nicht denkt, wie es unmöglich ist, daß der Gedanke einen
andern Stoff, als dem Gehirn als seinem Träger angehöre.} [It is as impossible that an undamaged brain
does not think, as it is impossible that the thought belongs to another material than the brain as its
bearer.]'' (Moleschott).

Thus an unscientific application of the teleological has occasioned an equally unscientific denial of all
teleology in nature. ``\textit{Wenn der Hirsch lange Beine zum Laufen hat, so hat er dieselben nicht
deßwegen
% ---- printed p.178 (PDF 186) ----
erhalten, um schnell laufen zu können, sondern er läuft schnell, weil er lange Beine hat.} [When the
stag has long legs for running, it has not received them for that reason, in order to be able to run
fast, but it runs fast because it has long legs.]'' (Büchner).

Finally, in the strife for and against the reality of the ideal, the unconditional Materialism has been
occasioned to apply a means that will not thus rightly suit Idealism, namely a peculiar kind of natural
pithiness. ``\textit{Die Gedanken stehen in demselben Verhältniß zu dem Gehirn, wie die Galle zur Leber
oder der Urin zu den Nieren, um mich einigermaßen grob hier
% [d/b glyph flag in the German: print ``auszubrücken'' for ``auszudrücken'' (= to express)]
auszudrücken.} [The thoughts stand in the same relation to the brain as the bile to the liver or the
urine to the kidneys, to express myself somewhat coarsely here.]'' (C. Vogt).

\begin{center}\rule{0.35\linewidth}{0.4pt}\end{center}

\begin{center}
  {\large\bfseries C.}\par\vspace{2pt}
  {\bfseries The Validity of the Idea.}\par\vspace{4pt}
  \rule{2cm}{0.4pt}
\end{center}
\phantomsection
\addcontentsline{toc}{subsection}{\texorpdfstring{C.\quad The Validity of the Idea}{C. The Validity of the Idea}}

\begin{center}{\bfseries \S\,34.}\end{center}

\emph{To deny the Idea is to deny cognition's a priori presuppositions; to acknowledge these
presuppositions is to acknowledge the Idea.}

That Materialism in the 18th century must be regarded as a peculiar consequence of Locke's Sensualism
has already been dealt with in the theory of knowledge (S. 14); in what relation the more recent
Materialism, on the contrary, will set itself to the question of the a priori is still, as it seems,
rather unclear.

``\textit{Die Welt}'' [The world], it is said in a work: ``\textit{Kraft und Stoff}'' [Force and Matter]
(von Dr. Büchner, Frankfurt a. M. 1855, S. 6) ``\textit{oder der Stoff mit seinen Eigenschaften, die wir
Kräfte nennen, mußten von Ewigkeit sein, und werden in Ewigkeit sein müssen.} [or matter with its
properties, which we call forces, must have been from eternity, and will have to be in eternity.]''
Here the category ``\textit{Ewigkeit}'' [eternity] is now used, with an emphasis so apodictically
decisive as though this concept held a priori, and the a priori held unconditionally. Meanwhile the
author does not seem to be quite sure in his matter; that is to
% ---- printed p.179 (PDF 187) ----
say: he is, as a genuine Materialist, naturally eternally convinced of matter's eternity; but about the
concept of the eternal he is nonetheless not quite clear with himself. ``\textit{Freilich}'' [To be
sure], he adds, ``\textit{ist der Begriff `Ewig' ein solcher, der sich schwer mit unsern endlichen
Verstandeskräften zu vertragen scheint} [is the concept `eternal' one that seems to comport with
difficulty with our finite powers of understanding]''; \ldots{} here the materialist dogmatist seems to
have a touch of the Humean scepsis.

If Materialism is really to assert the scientific superiority it ascribes to itself, then it must
necessarily express itself clearly and definitely concerning the a priori. It must then be scientifically
demonstrable how it properly comes about that one can a second time approach Locke without being troubled
by Hume, a second time do homage to Sensualism without being assailed by Scepticism.

Materialism is in its immediate self-sufficiency dogmatic. So long as this Dogmatism straightforwardly
builds on its own assertions, it naturally need not trouble itself the least about the results won in the
theory of knowledge; but when it once again condescends to tread the scientific path, it will scarcely be
able to avoid serious negotiations with Hume and Kant. That Materialism with its crudities must be to a
thinker like Kant insufferable goes without saying. He himself says of his critique of reason in the
preface to the second edition: ``\textit{Durch diese kann allein dem} \emph{Materialismus, Fatalismus,
Atheismus}\textit{, dem freigeisterischen} \emph{Unglauben}\textit{, der} \emph{Schwärmerei}
\textit{und} \emph{Aberglauben}\textit{, die allgemein schädlich werden können \ldots{} selbst die Wurzel
abgeschnitten werden.} [Through this alone can the very root be cut off from \emph{materialism, fatalism,
atheism}, from freethinking \emph{unbelief}, from \emph{enthusiasm} and \emph{superstition}, which can
become generally harmful \ldots{}]''

But to wish to exclude Materialism by excluding Speculation and holding to Kant alone is nonetheless also
unscientific. If one first grants the a priori validity, one is also compelled to grant the dialectical
validity; but, when the Dialectic is carried through, the critique of reason is transformed.
% ---- printed p.180 (PDF 188) ----
The rubrics set apart in the Kantian schema (cf. S. 24---26) contain, when the matter is seen
essentially, by no means separate for-themselves-existing pieces of a finitely composite whole of
consciousness, but dialectical moments of a rational totality-unity. The pure intuitions and the pure
concepts of the understanding do not fall outside or beside the ideas of reason; for sense and
understanding belong just as surely under reason as intuitions and concepts under the Idea.

To see how the critical bound at all points in the analysis of essence gives way to the dialectical, we
can here, e.g., make a trial with the separation of analytic and synthetic judgments.

``Matter is an extended thing,'' this judgment is analytic; for to deny a predicate that straightforwardly
follows from the subject's concept would indeed be a contradiction. But that in this analytic judgment
there is at the same time a synthetic moment is just as incontrovertible; it follows, namely, from this,
that two different intuitions --- extension and matter --- are in the judgment united into unity. By
wholly removing the synthetic we got only the tautological: ``Matter is matter, the extended is an
extended thing,'' but in such tautologies the analytic too is then lacking.

``The straight line is the shortest path between two points,'' and ``$5 + 7 = 12$,'' these propositions
Kant grasps as synthetic judgments a priori. ``\textit{Denn ich nehme zuerst die Zahl 7, und indem ich
für den Begriff der 5 die Finger meiner Hand als Anschauung zu Hilfe nehme, so thue ich die Einheiten,
die ich vorher zusammennahm, um die Zahl 5 auszumachen, nun an jenem meinem Bilde nach und nach zur Zahl
7, und sehe so die Zahl 12 entspringen \ldots{} Der arithmetische Satz ist also jederzeit synthetisch ..
Eben so wenig sind irgend ein Grundsatz der reinen Geometrie analytisch.} [For I first take the number 7,
and in that, for the concept of the 5, I take the fingers of my hand as intuition to help, I add the
units that I previously gathered together to make out the number 5, now to that image of mine little by
little up to the number 7, and thus see the number 12 arise \ldots{} The arithmetical proposition is thus
at all times synthetic .. Just as little is any principle of pure geometry analytic.]'' (Krit. d. r.
Vernft. Einleitung).

That in all mathematical judgments there is a synthetic element can so much the less be denied, as it has
precisely
% ---- printed p.181 (PDF 189) ----
shown itself that an analytic judgment is impossible when one excludes all co-determination by the
synthetic; but that the mathematical judgments should therefore, as Kant supposes, straightforwardly
exclude the analytic, is so striking a misunderstanding that Hegel in his displeasure hereat has even
explained it by ``a caprice.''

What has moved Kant to so finite a separation of the analytic and the synthetic is now nonetheless no
caprice; an unconditional holding-fast to the critical bound is, what Hegel could easily have seen, here
as everywhere an inviolable consequence of Criticism itself.

As it goes with the reciprocal determination of the analytic and the synthetic, so it goes with all the
other reciprocal determinations: from the side of reality the critical holds, but from the side of
essence the dialectical.

The critical separation of the pure intuition without concept and the pure concept without intuition is
idealess and in its essence untrue: ``\textit{Raum und Zeit sind besonders dann keine
Gedankenbestimmungen, wenn man keine Gedanken dabei hat: aber ein Begriff, sobald man einen Begriff davon
hat.} [Space and time are especially then no thought-determinations, when one has no thoughts in the
matter: but a concept, as soon as one has a concept of it.]'' (Hegel). Concept and intuition are moments
for one another mutually, their dialectical unity is their truth: they are in the Idea one. The critical
separation of the categories of quantity, quality, relation, and modality has in the finite its
significance, but is, when it is a matter of the essence itself, idealess and untrue. To criticism's
finite separation corresponds an equally finite connection; but the Idea is just as little a sum of
finitely connected categories as a for-itself, i.e. over against the others existing, category of reason.
The Idea is, as Idea, the categories' inner, principled totality, the truth in the categorical
determinations. The critical separation of a thing in itself without properties and a phenomenon of
properties without a thing in itself is idealess and untrue; the truth is that the thing is determined by
its properties and the properties by the thing: the infinite reciprocal determination
% ---- printed p.182 (PDF 190) ----
grounds itself on the Idea. The critical separation of a subject for itself without object and an object
for itself without subject is from the side of essence the highest idealessness. To deny the identity of
the subjective and the objective is to deny the Idea and make all cognition of truth impossible.

As for the assessment and the use of the categories' schema, an essential difference shows itself between
one-sidedness and thoughtlessness. That Hume denies the Ideas' a priori essentiality is one-sided, but
consistent, since he on the whole denies the a priori; that Kant himself ascribes to the Ideas only
subjective validity is likewise one-sided, but consistent, since he does not either ascribe to his a
priori intuitions and categories objective significance. But to wish to use some of the a priori
categories in a merely subjective, others on the contrary at the same time in an objective sense, is
thoughtlessness.

If one considers the a priori categories of the understanding for themselves, scarcely anyone is so
inconsistent as to ascribe to the ground objective validity when he grants the consequence only
subjective validity, or to grant both ground and consequence objective significance when the concept of
causality is to have only subjective significance; on the contrary, it is a not so seldom, though in its
essence equally thoughtless, inconsistency that one without more applies the a priori intuitions and
concepts of the understanding objectively, while one nonetheless regards the Ideas as something so
subjective that they exist only in our heads, without in the least concerning the essence of things.

Outside the analysis of essence it does indeed at first glance look as though the Idea were either merely
a subjective point, or, as far as the objective is concerned, nonetheless so infinitely far away that the
human understanding could scarcely descry it. But the matter stands reversed; the objective Idea is
infinitely near. In every intellectual glance, in every a priori category, in every point where knowledge
and being are united in truth, the Idea acts.
% ---- printed p.183 (PDF 191) ----
He who in mathematics will deny the validity of the higher analyses must consistently end by denying the
validity of the mathematical axioms; and yet there is such a distance between these axioms and those
analyses that ``sound understanding,'' which involuntarily acknowledges the axioms, must work long before
it raises itself to insight into the analyses' inner infinity. In corresponding manner it stands with the
analysis of the Idea and the a priori categories.

That it is one and the same fundamental thought that runs through the Platonic, the Aristotelian, and the
Hegelian doctrine of Ideas, so that one cannot --- what nonetheless sometimes happens in thoughtlessness
--- praise the Platonic Dialectic and cast contempt on the Hegelian, or, as far as the Idea is concerned,
support oneself on Aristotle in order to combat Hegel, is, when we look at the principle, incontrovertible.

The Hegelian Logic is in its grandiose execution a scientific proof of the Idea's validity. One can grant
that there are in this work of genius, as far as individual points and portions are concerned, obscurities
and blunders that need corrections, and that corrections will especially be conditioned by sharp preceding
analyses of reality; but one cannot give up the fundamental thought without at the same time either wholly
giving up philosophy, and thereby at bottom all the sciences' ideal foundation, or making oneself guilty
of a striking thoughtlessness.

Even if it is an error that the philosophizer should in Speculation be able to make himself one with the
Idea, it is nonetheless no error that the Idea as Idea lets itself be presented only in dialectical
scientific Speculation; even if it is a misunderstanding that philosophy should be able to give a knowledge
absolute, corresponding to reality, it is nonetheless no misunderstanding that the absolute identity of
essence and reality should be in the Idea itself; even if it is untrue that the universal concepts should
not be able to hold except by force
% [translation continues: §34 (printed p.184) onward — C. The Validity of the Idea]

\end{document}
